% Lesson 23 — Grammar: Partitive expressions, e.g. a bunch of …, a cup of … — Anglais 1ère A — Cours signé OnBuch+
% Programme officiel MINESEC. Compiler avec : tectonic ENA23-grammar-partitive-expressions-e-g-a-bunch-of-a-cup-of.tex
\newcommand{\DOCMATIERE}{Anglais}
\newcommand{\DOCNIVEAU}{1ère A}
\newcommand{\DOCMODULE}{Chapter 2 : Economic Life and Occupations}
\newcommand{\DOCLECON}{ENA23}
\newcommand{\DOCTITRE}{Lesson 23 — Grammar: Partitive expressions, e.g. a bunch of …, a cup of …}
% OnBuch+ — préambule « cours signés » ENRICHI (compilé avec tectonic / XeLaTeX)
% Système visuel « Pop » d'OnBuch+ : fond crème (jamais blanc), bordures encre
% épaisses, ombres « dures » décalées, titres Archivo Black en capitales.
% Version MATHÉMATIQUES : graphes (pgfplots/tikz), boîtes pédagogiques
% (définition, propriété, méthode, attention, le savais-tu, exercice résolu,
% exercice, TP, bilan), page de garde et signature OnBuch+ ancrée en bas.
%
% Macros attendues dans main.tex AVANT \input{preamble} :
%   \DOCMATIERE  \DOCNIVEAU  \DOCMODULE  \DOCLECON (numéro)  \DOCTITRE
\documentclass[11pt]{article}

\usepackage{fontspec}
\usepackage[english,french]{babel} % français = langue principale ; l'anglais via \eng{..} / textebox
\usepackage[a4paper,margin=2cm,top=3.4cm,bottom=2.8cm,headheight=1.4cm,footskip=1.3cm]{geometry}
\usepackage{amsmath,amssymb,amsfonts}
\usepackage{mathtools} % \xmapsto, \coloneqq... souvent utilisés par les modèles
\usepackage{cancel} % simplifications barrées (bilans de forces)
% \abs{z} (module, valeur absolue) et \norm{u} (norme), utilisés par les modèles
\providecommand{\abs}{}\let\abs\relax\DeclarePairedDelimiter{\abs}{\lvert}{\rvert}
\providecommand{\norm}{}\let\norm\relax\DeclarePairedDelimiter{\norm}{\lVert}{\rVert}
\usepackage[table]{xcolor} % [table] : fournit aussi \rowcolors (alternance de lignes)
\usepackage{pagecolor}
\usepackage{tikz}
\usetikzlibrary{arrows.meta,positioning,calc,shapes.geometric,shapes.misc,shapes.arrows,decorations.pathmorphing,decorations.pathreplacing,decorations.markings,patterns,shadows,mindmap,trees,fit,backgrounds,matrix,intersections,angles,quotes}
\usepackage{pgfplots}
\usepackage[version=4]{mhchem} % équations nucléaires \ce{^{235}_{92}U}
\usepackage[european,siunitx]{circuitikz} % schémas électriques
\pgfplotsset{compat=1.18}
\usepgfplotslibrary{fillbetween,groupplots} % aires entre courbes (calcul intégral)
\usepackage{enumitem}
\usepackage{fancyhdr}
% [most] charge toutes les bibliothèques tcolorbox (listings, minted, raster,
% documentation...) et gonfle inutilement la mémoire TeX sur un document long
% (~300 boîtes) : on ne charge que ce qui est réellement utilisé.
\usepackage[skins,breakable]{tcolorbox}
\usepackage{graphicx}
\usepackage{colortbl}
\usepackage{booktabs}
\usepackage{tabularx}
\usepackage{diagbox} % case d'en-tête diagonale des tableaux à double entrée
% Colonnes extensibles centrée / gauche / droite (C, L, R), souvent utilisées par les modèles
\newcolumntype{C}{>{\centering\arraybackslash}X}
\newcolumntype{L}{>{\raggedright\arraybackslash}X}
\newcolumntype{R}{>{\raggedleft\arraybackslash}X}
\usepackage{array}
\usepackage{multicol}
\usepackage{multirow}
\usepackage{siunitx}
% parse-numbers=false : accepte aussi \SI{2,5 \times 10^{-3}}{...}
\sisetup{locale=FR,output-decimal-marker={,},group-minimum-digits=4,parse-numbers=false}
\DeclareSIUnit\ohm{\ensuremath{\symup{\Omega}}} % U+2126 absent de la police
\DeclareSIUnit\division{div} % réglages d'oscilloscope (ms/div, V/div)
\DeclareSIUnit\tr{tr} % tours (vitesse de rotation, ex. \SI{1500}{\tr\per\minute})
\DeclareSIUnit\molar{M} % concentration molaire (ex. \SI{3}{\molar}, \SI{1}{\micro\molar})
\DeclareSIUnit\an{an} % année (le modèle confond parfois avec \year, un compteur TeX) — ex. \SI{5}{\kilo\meter\per\an}
\DeclareSIUnit\FCFA{FCFA} % franc CFA (ex. \SI{500}{\FCFA\per\kilo\gram})
\DeclareSIUnit\degreeBrix{\textdegree Brix} % teneur en sucre d'un sirop/jus (réfractométrie)
\DeclareSIUnit\month{mois} % le modèle utilise parfois \month, un compteur TeX, comme unité de durée
\usepackage{qrcode}
% \degree : commande fréquemment utilisée par les modèles pour un angle ou une
% température en dehors de \SI{}{} (non fournie par les paquets chargés).
\providecommand{\degree}{^{\circ}}
% \qed (fin de démonstration), utilisé par les modèles sans amsthm
\providecommand{\qed}{\hfill$\square$}
% \barre : double barre des valeurs interdites dans un tableau de variations (tkz-tab)
\providecommand{\barre}{\ensuremath{\|}}
% environnement proof (amsthm non chargé) utilisé par les modèles
\ifdefined\proof\else\newenvironment{proof}[1][Démonstration]{\par\noindent\textit{#1.}\ }{\qed\par}\fi
\usepackage{lastpage}
\usepackage{hyperref}
\hypersetup{hidelinks}

% ============================================================
% Palette « Pop » OnBuch+ — mêmes hex que la classe P (plus_ui.dart)
% ============================================================
\definecolor{poporange}{HTML}{F59321}
\definecolor{popdark}{HTML}{E07A0C}
\definecolor{popcream}{HTML}{FFF6E8}
\definecolor{popink}{HTML}{1C1714}
\definecolor{poppurple}{HTML}{7A5AE0}
\definecolor{popgreen}{HTML}{2E9E6A}
\definecolor{poppink}{HTML}{E0457B}
\definecolor{popblue}{HTML}{2F7DE1}
\definecolor{popgold}{HTML}{C98A12}
\definecolor{popmuted}{HTML}{8A7F76}
\definecolor{popline}{HTML}{EADFD2}
\definecolor{popgray}{HTML}{8A7F76}
\colorlet{popgrey}{popgray}
% Alias tolérants (le modèle nomme parfois les couleurs différemment)
\colorlet{popwhite}{white}
\colorlet{popblack}{popink}
\colorlet{popred}{poppink}
\colorlet{popyellow}{popgold}
% Teintes claires (fonds de tableaux/encadrés) — le modèle nomme parfois
% les couleurs avec un suffixe L (ex. popblueL) sans les avoir définies.
\colorlet{poporangeL}{poporange!20}
\colorlet{popdarkL}{popdark!20}
\colorlet{poppurpleL}{poppurple!20}
\colorlet{popgreenL}{popgreen!20}
\colorlet{poppinkL}{poppink!20}
\colorlet{popblueL}{popblue!20}
\colorlet{popgoldL}{popgold!20}
\colorlet{popredL}{popred!20}
\colorlet{popgrayL}{popgray!20}
% Teintes claires pour fonds de tableaux / zones
\colorlet{popblueL}{popblue!10!white}
\colorlet{poporangeL}{poporange!14!white}
\colorlet{popgreenL}{popgreen!12!white}
\colorlet{poppurpleL}{poppurple!12!white}
\colorlet{poppinkL}{poppink!12!white}
\colorlet{popgoldL}{popgold!14!white}

\pagecolor{popcream}

% ============================================================
% Typographie
% ============================================================
\defaultfontfeatures{Ligatures=TeX}
\setmainfont{PlusJakartaSans-Medium.ttf}[Path=../../../fonts/,BoldFont=PlusJakartaSans-Bold.ttf]
% Symboles, API et emojis absents de la police principale : repli sur DejaVu Sans (caractères actifs)
\newfontface\symfont{DejaVuSans.ttf}[Path=../../../fonts/]
\makeatletter
\newcommand\symactive[1]{\catcode#1=13 \begingroup\lccode`\~=#1 \lowercase{\endgroup\def~}{{\symfont\char#1}}}
\newcommand\symblank[1]{\catcode#1=13 \begingroup\lccode`\~=#1 \lowercase{\endgroup\def~}{}}
\newcommand\symhyph[1]{\catcode#1=13 \begingroup\lccode`\~=#1 \lowercase{\endgroup\def~}{\mbox{-}}}
\newcount\symc
\newcommand\symrange[2]{\symc=#1 \@whilenum\symc<#2 \do{\expandafter\symactive\expandafter{\the\symc}\advance\symc by 1 }}
\newcommand\symblankrange[2]{\symc=#1 \@whilenum\symc<#2 \do{\expandafter\symblank\expandafter{\the\symc}\advance\symc by 1 }}
\makeatother
\symrange{"0250}{"02FF}\symactive{"2713}\symactive{"2714}\symactive{"2717}\symactive{"2718}\symactive{"2610}\symactive{"2611}\symactive{"2612}
\symhyph{"2011}\symhyph{"2E1A}\symblankrange{"1F300}{"1FAFF}\symblank{"FE0F}
% Maths en sans-serif (Fira Math) avec chiffres et lettres droites de Plus
% Jakarta, pour que les formules se fondent dans le texte.
\usepackage[math-style=ISO,bold-style=ISO]{unicode-math}
\setmathfont{FiraMath-Regular.otf}[Path=../../../fonts/]
\setmathfont{PlusJakartaSans-Medium.ttf}[Path=../../../fonts/,range={up/num,up/{Latin,latin},\mathup}]
\usepackage{icomma} % virgule décimale sans espace en mode math
% Caractères mathématiques Unicode absents des polices de texte (Plus Jakarta,
% Archivo Black) : rendus par leur équivalent mathématique, en texte comme en
% formule — sinon ils s'affichent en carré vide (ex. « Arithmétique dans ℤ »).
\usepackage{newunicodechar}
\newunicodechar{ℝ}{\ensuremath{\mathbb{R}}}
\newunicodechar{ℤ}{\ensuremath{\mathbb{Z}}}
\newunicodechar{ℂ}{\ensuremath{\mathbb{C}}}
\newunicodechar{ℕ}{\ensuremath{\mathbb{N}}}
\newunicodechar{ℚ}{\ensuremath{\mathbb{Q}}}
\newunicodechar{𝔻}{\ensuremath{\mathbb{D}}}
\newunicodechar{α}{\ensuremath{\alpha}}
\newunicodechar{β}{\ensuremath{\beta}}
\newunicodechar{γ}{\ensuremath{\gamma}}
\newunicodechar{δ}{\ensuremath{\delta}}
\newunicodechar{ε}{\ensuremath{\varepsilon}}
\newunicodechar{θ}{\ensuremath{\theta}}
\newunicodechar{λ}{\ensuremath{\lambda}}
\newunicodechar{μ}{\ensuremath{\mu}}
\newunicodechar{σ}{\ensuremath{\sigma}}
\newunicodechar{τ}{\ensuremath{\tau}}
\newunicodechar{φ}{\ensuremath{\varphi}}
\newunicodechar{ψ}{\ensuremath{\psi}}
\newunicodechar{ω}{\ensuremath{\omega}}
\newunicodechar{Σ}{\ensuremath{\Sigma}}
% majuscules produites par \MakeUppercase dans les titres (α-aminés -> Α)
\newunicodechar{Α}{\ensuremath{\alpha}}
\newunicodechar{Β}{\ensuremath{\beta}}
\newunicodechar{Γ}{\ensuremath{\Gamma}}
\newunicodechar{Δ}{\ensuremath{\Delta}}
\newunicodechar{Ε}{\ensuremath{\varepsilon}}
\newunicodechar{Θ}{\ensuremath{\Theta}}
\newunicodechar{Λ}{\ensuremath{\Lambda}}
\newunicodechar{Μ}{\ensuremath{\mu}}
\newunicodechar{Τ}{\ensuremath{\tau}}
\newunicodechar{Φ}{\ensuremath{\Phi}}
\newunicodechar{Ψ}{\ensuremath{\Psi}}
\newunicodechar{Ω}{\ensuremath{\Omega}}
\newunicodechar{↦}{\ensuremath{\mapsto}}
\newunicodechar{∈}{\ensuremath{\in}}
\newunicodechar{∉}{\ensuremath{\notin}}
\newunicodechar{∧}{\ensuremath{\wedge}}
\newunicodechar{⇔}{\ensuremath{\Leftrightarrow}}
\newunicodechar{⟺}{\ensuremath{\Longleftrightarrow}}
\newunicodechar{⇒}{\ensuremath{\Rightarrow}}
\newunicodechar{⟹}{\ensuremath{\Longrightarrow}}
\newunicodechar{∘}{\ensuremath{\circ}}
\newunicodechar{∪}{\ensuremath{\cup}}
\newunicodechar{∩}{\ensuremath{\cap}}
\newunicodechar{≡}{\ensuremath{\equiv}}
\newunicodechar{⊂}{\ensuremath{\subset}}
\newunicodechar{⊥}{\ensuremath{\perp}}
\newunicodechar{∃}{\ensuremath{\exists}}
\newunicodechar{∀}{\ensuremath{\forall}}
\newunicodechar{∛}{\ensuremath{\sqrt[3]{\,}}}
\newunicodechar{∜}{\ensuremath{\sqrt[4]{\,}}}
\newunicodechar{⃗}{\ensuremath{{}^{\to}}}
\newunicodechar{ⁿ}{\ifmmode^{n}\else\textsuperscript{\ensuremath{n}}\fi}
\newunicodechar{ˣ}{\ifmmode^{x}\else\textsuperscript{\ensuremath{x}}\fi}
\newunicodechar{ᵇ}{\ifmmode^{b}\else\textsuperscript{\ensuremath{b}}\fi}
\newunicodechar{ʳ}{\ifmmode^{r}\else\textsuperscript{\ensuremath{r}}\fi}
\newunicodechar{ᵘ}{\ifmmode^{u}\else\textsuperscript{\ensuremath{u}}\fi}
\newunicodechar{ʸ}{\ifmmode^{y}\else\textsuperscript{\ensuremath{y}}\fi}
\newunicodechar{ᵃ}{\ifmmode^{a}\else\textsuperscript{\ensuremath{a}}\fi}
\newunicodechar{ᵉ}{\ifmmode^{e}\else\textsuperscript{\ensuremath{e}}\fi}
\newunicodechar{ᵏ}{\ifmmode^{k}\else\textsuperscript{\ensuremath{k}}\fi}
\newunicodechar{ᵖ}{\ifmmode^{p}\else\textsuperscript{\ensuremath{p}}\fi}
\newunicodechar{ᴺ}{\ifmmode^{N}\else\textsuperscript{\ensuremath{N}}\fi}
\newunicodechar{ᶜ}{\ifmmode^{c}\else\textsuperscript{\ensuremath{c}}\fi}
\newunicodechar{ʰ}{\ifmmode^{h}\else\textsuperscript{\ensuremath{h}}\fi}
\newunicodechar{ᵠ}{\ifmmode^{q}\else\textsuperscript{\ensuremath{q}}\fi}
\newunicodechar{ₙ}{\ifmmode_{n}\else\textsubscript{\ensuremath{n}}\fi}
\newunicodechar{ₐ}{\ifmmode_{a}\else\textsubscript{\ensuremath{a}}\fi}
\newunicodechar{ᵢ}{\ifmmode_{i}\else\textsubscript{\ensuremath{i}}\fi}
\newunicodechar{ᵣ}{\ifmmode_{r}\else\textsubscript{\ensuremath{r}}\fi}
\newunicodechar{ₖ}{\ifmmode_{k}\else\textsubscript{\ensuremath{k}}\fi}
\newunicodechar{ᵦ}{\ifmmode_{\beta}\else\textsubscript{\ensuremath{\beta}}\fi}
\newunicodechar{ₓ}{\ifmmode_{x}\else\textsubscript{\ensuremath{x}}\fi}
\newfontfamily\popdisplay{ArchivoBlack-Regular.ttf}[Path=../../../fonts/]
\newfontfamily\poplogo{SpaceGrotesk-Bold.ttf}[Path=../../../fonts/]
\newfontfamily\popsemi{PlusJakartaSans-SemiBold.ttf}[Path=../../../fonts/]
\newfontfamily\popxbold{PlusJakartaSans-ExtraBold.ttf}[Path=../../../fonts/]
\newfontfamily\popxboldls{PlusJakartaSans-ExtraBold.ttf}[Path=../../../fonts/,LetterSpace=3.0]

\setlist[enumerate]{leftmargin=1.7em,itemsep=5pt,topsep=4pt}
\setlist[itemize]{leftmargin=1.7em,itemsep=4pt,topsep=4pt,label={\color{poporange}\textbullet}}
\setlength{\parindent}{0pt}
\setlength{\parskip}{5pt}
\renewcommand{\arraystretch}{1.25}

% pgfplots : style Pop par défaut
\pgfplotsset{
  every axis/.append style={axis line style={popink,line width=1pt},tick style={popink},
    grid style={popline,line width=0.5pt},label style={font=\small},tick label style={font=\footnotesize}},
  popcurve/.style={poporange,line width=1.8pt,no markers,smooth},
}
% Même style disponible dans un \draw TikZ simple (hors pgfplots)
\tikzset{popcurve/.style={poporange,line width=1.8pt}}
\tikzset{popgrid/.style={popline,very thin}}
\colorlet{popcurve}{poporange} % le nom du style est parfois utilisé comme couleur

% ============================================================
% Pill / squiggle
% ============================================================
\newcommand{\poppill}[3]{%
  \tikz[baseline=-0.55ex]\node[draw=popink,line width=1.2pt,rounded corners=2.6pt,%
    fill=#2,inner xsep=6.5pt,inner ysep=3pt]{%
    \popxboldls\fontsize{7}{7}\selectfont\color{#3}\MakeUppercase{#1}};%
}
\newcommand{\popsquiggle}[1]{%
  \tikz[baseline=-0.4ex]{
    \draw[#1,line width=2.6pt,line cap=round]
      plot [smooth] coordinates {(0,0.09) (0.34,0.01) (0.68,0.21) (1.02,0.01) (1.36,0.10)};
  }%
}
% Mot-clé surligné (marqueur orange sous le texte)
\newcommand{\cle}[1]{{\popxbold\color{popdark}#1}}

% ============================================================
% En-tête / pied
% ============================================================
\pagestyle{fancy}
\fancyhf{}
\renewcommand{\headrulewidth}{0pt}
\renewcommand{\footrulewidth}{0pt}
\lhead{\raisebox{-0.15ex}{\poplogo\fontsize{13}{13}\selectfont\color{popink}OnBuch\color{poporange}+}}
\rhead{\poppill{\DOCMATIERE\ \textbullet\ \DOCNIVEAU\ \textbullet\ Leçon \DOCLECON}{white}{popink}}
\lfoot{\popxbold\fontsize{7.6}{7.6}\selectfont\color{popmuted}\MakeUppercase{Cours signé OnBuch+}\ \textbullet\ {\popsemi\DOCTITRE}}
\rfoot{\tikz[baseline=-0.6ex]\node[rounded corners=8.5pt,fill=popink,minimum height=17pt,minimum width=17pt,inner xsep=5pt,inner sep=0]{\popxbold\fontsize{8}{8}\selectfont\color{popcream}\thepage\,/\,\pageref*{LastPage}};}

% ============================================================
% Titres
% ============================================================
\newcommand{\chaptitre}[2]{%
  \begin{tcolorbox}[enhanced,colback=poporange,colframe=popink,boxrule=2.6pt,arc=11pt,
      drop shadow=popink,left=15pt,right=15pt,top=12pt,bottom=11pt,before skip=3pt,after skip=18pt]
    \poppill{#2}{popink}{popcream}\par\vspace{9pt}
    {\raggedright\hyphenpenalty=10000\exhyphenpenalty=10000\popdisplay\fontsize{22}{24}\selectfont\color{popcream}\MakeUppercase{#1}\par}
  \end{tcolorbox}
}
% Section numérotée : pastille orange avec le numéro + titre Archivo
\newcounter{popsec}
\newcommand{\coursec}[1]{%
  \stepcounter{popsec}\setcounter{popsub}{0}%
  \par\vspace{16pt}\noindent
  \begin{minipage}{\linewidth}
  \tikz[baseline=-0.7ex]\node[draw=popink,line width=1.6pt,fill=poporange,rounded corners=4pt,minimum size=22pt,inner sep=0]{\popdisplay\fontsize{12}{12}\selectfont\color{popcream}\thepopsec};%
  \hspace{8pt}{\popdisplay\fontsize{15}{17}\selectfont\color{popink}\MakeUppercase{#1}}\par
  \vspace{4pt}\noindent{\color{poporange}\rule{36pt}{3.6pt}}
  \end{minipage}\par\nopagebreak\vspace{8pt}
}
\newcounter{popsub}
\newcommand{\courssub}[1]{%
  \stepcounter{popsub}%
  \par\vspace{9pt}\noindent{\popxbold\fontsize{11.5}{13}\selectfont\color{popink}{\color{poporange}\thepopsec.\thepopsub}\hspace{6pt}#1}\par\nopagebreak\vspace{4pt}
}

% ============================================================
% Boîtes pédagogiques — même recette : fond blanc, bordure épaisse colorée,
% ombre dure encre, badge plein. Titre optionnel [#1].
% ============================================================
\tcbset{popbase/.style={breakable,enhanced,colback=white,boxrule=2.3pt,arc=9pt,
  shadow={3pt}{-3pt}{0pt}{popink},left=14pt,right=14pt,top=11pt,bottom=11pt,before skip=11pt,after skip=15pt,
  fontupper=\popsemi\fontsize{10.6}{14.5}\selectfont\color{popink},
  fontlower=\popsemi\fontsize{10.6}{14.5}\selectfont\color{popink}}}
% Coche dessinée (le glyphe ✓ n'existe pas dans les polices du document)
\newcommand{\popcheck}[1]{\tikz[baseline=-0.5ex]\draw[#1,line width=1.6pt,line cap=round,line join=round](-0.13,0.0)--(-0.04,-0.09)--(0.13,0.1);}
\providecommand{\coche}{\popcheck{popgreen}}
\providecommand{\resultat}[1]{\par\smallskip\noindent\fcolorbox{popgreen}{popgreenL}{\parbox{\dimexpr\linewidth-2\fboxsep-2\fboxrule\relax}{\textbf{Résultat.} #1}}\par\smallskip}
\newcommand{\popboxhead}[3]{\poppill{#1}{#2}{white}\ifx\relax#3\relax\else\hspace{6pt}{\popxbold\fontsize{10.6}{12}\selectfont\color{popink}#3}\fi\par\vspace{7pt}}

\newtcolorbox{aretenir}[1][]{popbase,colframe=popgreen,before upper={\popboxhead{À retenir}{popgreen}{#1}}}
% Texte en anglais (document de lecture, dialogue, extrait) : cadre
% d'étude. Un titre optionnel (source, auteur) s'affiche dans l'en-tête.
\newtcolorbox{textebox}[1][]{popbase,colframe=popblue,before upper={\popboxhead{Text}{popblue}{#1}\begin{otherlanguage}{english}},after upper={\end{otherlanguage}}}
% Marques ✓ / ✗ (forme correcte / fautive) souvent utilisées par les modèles
\providecommand{\cross}{\textcolor{poppink}{\ensuremath{\times}}}
\providecommand{\xmark}{\cross}\providecommand{\crossmark}{\cross}\providecommand{\cmark}{\ensuremath{\checkmark}}
% Ligne de réponse / justification à compléter (inventée par les modèles)
\providecommand{\justification}[1]{\par\noindent\textit{Justification~:} \dotfill\par}
% \textipa (package tipa non chargé) : la police ne couvre pas l'API -> simple passage
\providecommand{\textipa}[1]{#1}
% Soulignement (ulem non chargé) : \uline, \ul
\providecommand{\uline}[1]{\underline{#1}}\providecommand{\ul}[1]{\underline{#1}}
% \glossaire{\item ...} inventée par les modèles : liste de glossaire
\providecommand{\glossaire}[1]{\begin{itemize}#1\end{itemize}}
% \alert (beamer) inventée par les modèles : texte mis en évidence
\providecommand{\alert}[1]{\textbf{\textcolor{poppink}{#1}}}
% variantes ASCII d'ordinaux écrites en macro
\providecommand{\iere}{\textsuperscript{re}}\providecommand{\ieme}{\textsuperscript{e}}\providecommand{\ere}{\textsuperscript{re}}\providecommand{\eme}{\textsuperscript{e}}\providecommand{\er}{\textsuperscript{er}}\providecommand{\ie}{\textsuperscript{e}}
% Ordinaux français écrits en macro par les modèles : 1\ière, 3\ième
\newcommand{\ière}{\textsuperscript{re}}\newcommand{\ième}{\textsuperscript{e}}
% \exemple (inventée par les modèles) : étiquette « Exemple : »
\providecommand{\exemple}{\textbf{Exemple~:}\ }
% \square utilisable aussi hors mode math (cases à cocher)
\AtBeginDocument{\let\squareMath\square\renewcommand{\square}{\ensuremath{\squareMath}}}
% Mot ou phrase en anglais à l'intérieur d'un paragraphe français
\newcommand{\eng}[1]{\ifmmode\text{\textit{#1}}\else\begingroup\selectlanguage{english}\itshape #1\endgroup\fi}
\let\esp\eng % alias toléré
% flèches d'intonation inventées par les modèles
\providecommand{\textsearrow}{}\renewcommand{\textsearrow}{\ensuremath{\searrow}}\providecommand{\textnearrow}{}\renewcommand{\textnearrow}{\ensuremath{\nearrow}}
% \fr{..} : mot français (inventé par les modèles)
\providecommand{\fr}[1]{\textit{#1}}
\newtcolorbox{exemplebox}[1][]{popbase,colframe=poporange,before upper={\popboxhead{Exemple}{poporange}{#1}}}
\newtcolorbox{definition}[1][]{popbase,colframe=popblue,before upper={\popboxhead{Définition}{popblue}{#1}}}
\newtcolorbox{propriete}[1][]{popbase,colframe=poppurple,before upper={\popboxhead{Propriété}{poppurple}{#1}}}
\newtcolorbox{methode}[1][]{popbase,colframe=popgold,before upper={\popboxhead{Méthode}{popgold}{#1}}}
\newtcolorbox{attention}[1][]{popbase,colframe=poppink,before upper={\popboxhead{Attention}{poppink}{#1}}}
\newtcolorbox{experience}[1][]{popbase,colframe=popblue,colback=popblueL,before upper={\popboxhead{Expérience}{popblue}{#1}}}
% « Le savais-tu ? » : carte violette pleine (contexte, culture, Cameroun)
\newtcolorbox{savaistu}[1][]{popbase,colback=poppurple,colframe=popink,
  fontupper=\popsemi\fontsize{10.4}{14.2}\selectfont\color{white},
  before upper={\poppill{Le savais-tu\,?}{white}{poppurple}\ifx\relax#1\relax\else\hspace{6pt}{\popxbold\color{white}#1}\fi\par\vspace{7pt}}}
% Exercice résolu : énoncé en haut, \tcblower puis la solution sur fond crème
\newcounter{popexo}\newcounter{popexr}
\newtcolorbox{exoresolu}[1][]{popbase,colframe=poporange,colbacklower=popcream,
  segmentation style={popink,line width=1.2pt,dashed},
  before upper={\stepcounter{popexr}\popboxhead{Exercice résolu \thepopexr}{poporange}{#1}},
  before lower={\poppill{Solution}{popink}{popcream}\par\vspace{6pt}}}
% Exercice (non résolu en place) — la correction est dans la partie Corrigés
\newtcolorbox{exercice}[1][]{popbase,colframe=popink,
  before upper={\stepcounter{popexo}\popboxhead{Exercice \thepopexo}{popink}{#1}}}
% Corrigé
\newtcolorbox{corrige}[1][]{popbase,colframe=popgreen,colback=popgreenL,
  before upper={\popboxhead{Corrigé}{popgreen}{#1}}}
% Objectifs / prérequis / bilan
\newtcolorbox{objectifs}{popbase,colframe=popink,colback=poporangeL,
  before upper={\popboxhead{Objectifs de la leçon}{popink}{}}}
\newtcolorbox{prerequis}{popbase,colframe=popink,colback=popblueL,
  before upper={\popboxhead{Prérequis}{popblue}{}}}
\newtcolorbox{bilan}{popbase,colframe=popink,colback=white,boxrule=2.8pt,
  before upper={\popboxhead{Fiche bilan}{poporange}{L'essentiel à connaître pour l'examen}}}
% Figure encadrée avec légende
\newtcolorbox{popfigure}[1][]{enhanced,breakable=false,colback=white,colframe=popline,boxrule=1.4pt,arc=8pt,
  left=10pt,right=10pt,top=10pt,bottom=8pt,before skip=10pt,after skip=12pt,halign=center,
  lower separated=false,halign lower=center,
  fontlower=\popsemi\fontsize{9}{11.5}\selectfont\color{popmuted},
  #1}
\newcommand{\legende}[1]{\par\vspace{6pt}{\popsemi\fontsize{9}{11.5}\selectfont\color{popmuted}#1}}

% ============================================================
% Signature OnBuch+
% ============================================================
\newcommand{\poplogoinline}[1]{{\poplogo\fontsize{#1}{#1}\selectfont\color{popink}OnBuch\color{poporange}+}}
\newcommand{\signatureonbuch}{%
  \begin{tcolorbox}[enhanced,colback=popink,colframe=popink,boxrule=2.2pt,arc=10pt,
      drop shadow=poporange,left=14pt,right=14pt,top=10pt,bottom=10pt,before skip=0pt,after skip=0pt]
    \tikz[baseline=-0.7ex]\node[circle,fill=popgreen,draw=popcream,line width=1.3pt,minimum size=22pt,inner sep=0]{\popcheck{white}};%
    \hspace{9pt}\begin{minipage}[c]{0.62\linewidth}
      {\popxbold\fontsize{12}{13}\selectfont\color{popcream}Signé\ \textbullet\ {\poplogo OnBuch\color{poporange}+}}\par\vspace{2pt}
      {\raggedright\popsemi\fontsize{8}{10.5}\selectfont\color{popline}\MakeUppercase{Équipe pédagogique OnBuch+}\\\MakeUppercase{Conforme au programme officiel MINESEC}\par}
    \end{minipage}\hfill
    {\poplogo\fontsize{22}{22}\selectfont\color{popcream}OnBuch\color{poporange}+}
  \end{tcolorbox}%
}

% ============================================================
% Page de garde
% #1 titre · #2 matière · #3 niveau · #4 module · #5 n° leçon · #6 durée ·
% #7 description / objectifs (texte ou itemize)
% ============================================================
\newcommand{\pagegarde}[7]{%
  \thispagestyle{empty}
  \newgeometry{a4paper,margin=1.7cm,top=1.6cm,bottom=1.4cm}%
  \begin{tikzpicture}[remember picture,overlay]
    \fill[poporange!18] ([xshift=-3.2cm,yshift=-3.6cm]current page.north east) circle (4.6cm);
    \fill[poppurple!14] ([xshift=2.4cm,yshift=7.6cm]current page.south west) circle (3.8cm);
  \end{tikzpicture}%
  {\poplogo\fontsize{30}{30}\selectfont\color{popink}OnBuch\color{poporange}+}\hfill
  \raisebox{0.6ex}{\poppill{Cours signé \textbullet\ Premium}{popink}{popcream}}\par
  \vspace{1pt}\popsquiggle{poppurple}\par
  \vspace{8pt}
  \begin{tcolorbox}[enhanced,colback=poporange,colframe=popink,boxrule=2.8pt,arc=13pt,
      drop shadow=popink,left=18pt,right=18pt,top=12pt,bottom=13pt,after skip=10pt]
    \poppill{#2 \textbullet\ #3}{popink}{popcream}\hspace{5pt}\poppill{#4}{white}{popink}\par\vspace{8pt}
    {\popdisplay\fontsize{14}{14}\selectfont\color{popink}LEÇON #5}\par\vspace{4pt}
    {\raggedright\hyphenpenalty=10000\exhyphenpenalty=10000\popdisplay\fontsize{25}{27}\selectfont\color{popcream}\MakeUppercase{#1}\par}
    \vspace{8pt}
    {\popxbold\fontsize{9.5}{9.5}\selectfont\color{popink}\MakeUppercase{Durée conseillée : #6}}
  \end{tcolorbox}
  \begin{tcolorbox}[enhanced,colback=white,colframe=popink,boxrule=2.2pt,arc=10pt,
      drop shadow=popink,left=16pt,right=16pt,top=10pt,bottom=10pt,after skip=8pt]
    \poppill{Dans cette leçon}{poppurple}{white}\par\vspace{5pt}
    \popsemi\fontsize{9.3}{12.6}\selectfont\color{popink}#7
  \end{tcolorbox}
  \noindent\begin{minipage}[t]{0.487\linewidth}
    \begin{tcolorbox}[enhanced,colback=popgreen,colframe=popink,boxrule=2.2pt,arc=10pt,
        drop shadow=popink,left=11pt,right=11pt,top=10pt,bottom=10pt,height=3.1cm,valign=center]
      \tcbox[colback=white,colframe=popink,boxrule=1.3pt,arc=5pt,boxsep=0pt,nobeforeafter,
        left=3pt,right=3pt,top=3pt,bottom=3pt,valign=center]{\qrcode[height=1.9cm]{https://play.google.com/store/apps/details?id=com.luvvix.onbuch&hl=fr}}%
      \hspace{8pt}\begin{minipage}[c]{0.5\linewidth}
        {\popdisplay\fontsize{11}{12.5}\selectfont\color{white}\MakeUppercase{Télécharge OnBuch}}\par\vspace{4pt}
        {\raggedright\popsemi\fontsize{8.4}{11}\selectfont\color{white}Gratuit \textbullet\ Google Play\\Tuteur IA, annales, concours, résultats\par}
      \end{minipage}
    \end{tcolorbox}
  \end{minipage}\hfill
  \begin{minipage}[t]{0.487\linewidth}
    \begin{tcolorbox}[enhanced,colback=poppurple,colframe=popink,boxrule=2.2pt,arc=10pt,
        drop shadow=popink,left=11pt,right=11pt,top=10pt,bottom=10pt,height=3.1cm,valign=center]
      \tikz[baseline=-0.7ex]\node[circle,fill=white,draw=popink,line width=1.3pt,minimum size=1.6cm,inner sep=0]{\popdisplay\fontsize{24}{24}\selectfont\color{poppurple}+};%
      \hspace{8pt}\begin{minipage}[c]{0.6\linewidth}
        {\popdisplay\fontsize{11}{12.5}\selectfont\color{white}\MakeUppercase{Passe à OnBuch+}}\par\vspace{4pt}
        {\raggedright\popsemi\fontsize{8.4}{11}\selectfont\color{white}Tous les cours signés, Tuteur IA illimité, fiches PDF — onglet OnBuch+ de l'app\par}
      \end{minipage}
    \end{tcolorbox}
  \end{minipage}
  \vspace{8pt}
  \popcontenu
  \vfill
  \signatureonbuch
  \restoregeometry
  \newpage
}

% Rangée « Ce cours contient » (page de garde)
\newcommand{\poptile}[3]{% #1 couleur #2 gros texte #3 légende
  \begin{tcolorbox}[enhanced,colback=#1,colframe=popink,boxrule=2pt,arc=9pt,shadow={2.5pt}{-2.5pt}{0pt}{popink},
      left=6pt,right=6pt,top=9pt,bottom=9pt,halign=center,valign=center,height=1.75cm,width=\linewidth,nobeforeafter]
    {\popdisplay\fontsize{12.5}{13}\selectfont\color{white}\MakeUppercase{#2}}\par\vspace{3pt}
    {\centering\popsemi\fontsize{7.8}{9.5}\selectfont\color{white}#3\par}
  \end{tcolorbox}}
\newcommand{\popcontenu}{%
  \par\noindent{\popxboldls\fontsize{7.5}{7.5}\selectfont\color{popmuted}CE COURS SIGNÉ CONTIENT}\par\vspace{7pt}
  \noindent\begin{minipage}[t]{0.235\linewidth}\poptile{popblue}{Cours}{complet, illustré, pas à pas}\end{minipage}\hfill
  \begin{minipage}[t]{0.235\linewidth}\poptile{poporange}{Méthodes}{et exercices résolus}\end{minipage}\hfill
  \begin{minipage}[t]{0.235\linewidth}\poptile{poppink}{Exercices}{type examen + corrigés}\end{minipage}\hfill
  \begin{minipage}[t]{0.235\linewidth}\poptile{popgreen}{Fiche bilan}{l'essentiel à retenir}\end{minipage}\par}

% Sommaire
\newtcolorbox{sommaire}{enhanced,colback=white,colframe=popink,boxrule=2.2pt,arc=10pt,
  drop shadow=popink,left=18pt,right=18pt,top=13pt,bottom=13pt,before skip=6pt,after skip=20pt,
  before upper={\poppill{Sommaire}{poppurple}{white}\par\vspace{9pt}\popsemi\fontsize{10.6}{14.5}\selectfont\color{popink}}}

% Signature de fin de cours — poussée tout en bas de la dernière page
\newcommand{\findecours}{%
  \par\vfill\nopagebreak
  \begin{center}\popsquiggle{poporange}\end{center}\vspace{4pt}
  \signatureonbuch
}
\AtBeginDocument{\def\llbracket{\mathopen{[\mkern-3.5mu[}}\def\rrbracket{\mathclose{]\mkern-3.5mu]}}} % intervalles d'entiers (glyphe ⟦ absent de la police)
% Glyphes absents de Fira Math : redéfinis à partir de glyphes disponibles
\AtBeginDocument{%
  \def\setminus{\mathbin{\backslash}}\let\smallsetminus\setminus
  \def\vdots{\mathord{\vbox{\baselineskip3.2pt\lineskiplimit0pt\kern4pt\hbox{.}\hbox{.}\hbox{.}}}}%
  \def\ddots{\mathinner{\mkern1mu\raise6.5pt\hbox{.}\mkern2mu\raise3.6pt\hbox{.}\mkern2mu\raise0.7pt\hbox{.}\mkern1mu}}%
  \def\triangle{\mathord{\scalebox{0.9}{$\symup{\Delta}$}}}%
  \def\nabla{\mathord{\raisebox{\depth}{\scalebox{1}[-1]{$\symup{\Delta}$}}}}%
  \def\checkmark{\popcheck{popdark}}%
  \def\star{\mathbin{\ast}}\def\bigstar{\mathord{\ast}}%
  \def\diamond{\mathbin{\scalebox{0.6}{$\Diamond$}}}%
  \def\bigcup{\mathop{\mathchoice{\raisebox{-0.3em}{\scalebox{1.7}{$\cup$}}}{\scalebox{1.25}{$\cup$}}{\cup}{\cup}}\displaylimits}%
  \def\bigcap{\mathop{\mathchoice{\raisebox{-0.3em}{\scalebox{1.7}{$\cap$}}}{\scalebox{1.25}{$\cap$}}{\cap}{\cap}}\displaylimits}%
  \def\lesseqqgtr{\mathrel{\lessgtr}}\def\bigoplus{\mathop{\oplus}\displaylimits}%
}
\providecommand{\ssi}{si et seulement si} % abréviation fréquente
\AtBeginDocument{\providecommand{\wideparen}{\overparen}} % arc de cercle

\begin{document}

\pagegarde{Lesson 23 — Grammar: Partitive expressions, e.g. a bunch of …, a cup of …}{Anglais}{1ère A}{Chapter 2 : Economic Life and Occupations}{ENA23}{2 h}{Cette leçon de grammaire présente les expressions partitives anglaises (a cup of, a piece of, a bottle of, a bunch of...) qui permettent de quantifier les noms indénombrables. Les élèves apprennent à choisir la bonne expression selon le nom, à éviter les erreurs typiques des francophones et à utiliser ces structures dans des situations concrètes de la vie économique (marché, boutique, recettes).
\vspace{4pt}
\begin{multicols}{2}\small
\begin{itemize}
\item À la fin de cette leçon, je sais distinguer les noms dénombrables des noms indénombrables en anglais.
\item À la fin de cette leçon, je sais identifier une expression partitive dans une phrase.
\item À la fin de cette leçon, je sais construire correctement la structure : quantité + partitif + of + nom indénombrable.
\item À la fin de cette leçon, je sais associer le bon partitif au bon nom (a loaf of bread, a bunch of bananas, a piece of advice...).
\item À la fin de cette leçon, je sais mettre les expressions partitives au pluriel (two cups of tea, three bottles of oil).
\item À la fin de cette leçon, je sais éviter les erreurs fréquentes des francophones (*an information, *a bread, *advices).
\end{itemize}\end{multicols}}

\begin{sommaire}
\begin{multicols}{2}
\begin{enumerate}
\item Observation et découverte
\item La règle : structure et formation
\item Les principaux partitifs par catégories
\item Comparaison français / anglais
\item Emplois en contexte et cas particuliers
\item Exercices d'application et de transformation
\item Activité d'intégration
\item Méthodes et exercices résolus
\item Exercices
\item Corrigés des exercices
\item Fiche bilan
\end{enumerate}
\end{multicols}
\end{sommaire}

\begin{objectifs}
\begin{itemize}
\item À la fin de cette leçon, je sais distinguer les noms dénombrables des noms indénombrables en anglais.
\item À la fin de cette leçon, je sais identifier une expression partitive dans une phrase.
\item À la fin de cette leçon, je sais construire correctement la structure : quantité + partitif + of + nom indénombrable.
\item À la fin de cette leçon, je sais associer le bon partitif au bon nom (a loaf of bread, a bunch of bananas, a piece of advice...).
\item À la fin de cette leçon, je sais mettre les expressions partitives au pluriel (two cups of tea, three bottles of oil).
\item À la fin de cette leçon, je sais éviter les erreurs fréquentes des francophones (*an information, *a bread, *advices).
\item À la fin de cette leçon, je sais utiliser les partitifs dans des phrases et un court texte (liste de courses, dialogue au marché).
\item À la fin de cette leçon, je sais comparer le partitif anglais avec les articles partitifs français (du, de la, des).
\end{itemize}
\end{objectifs}

\begin{prerequis}
\begin{itemize}
\item Les articles a/an, the et l'article zéro (Seconde)
\item Les noms dénombrables et indénombrables : notion de base (Seconde)
\item Some, any, much, many, a lot of (Seconde)
\item Le pluriel des noms réguliers et irréguliers
\item Le vocabulaire de base de la nourriture et des courses (Chapter 2, lessons précédentes)
\end{itemize}
\end{prerequis}

\newpage

\coursec{Observation et découverte}

\courssub{Une situation de communication : la liste de courses de Sarah}

It is Saturday morning at Mokolo Market in Yaoundé. Sarah, a Lower Sixth student, is helping her mother write a shopping list in English for her uncle who lives in Bamenda. She writes: \eng{I need two rices, three breads and one water.} Her uncle laughs on the phone: \eng{In English, you can't count rice, bread or water directly — you need a special word before them!} Sarah is confused: in French, she says \emph{du riz}, \emph{de l'eau} without any problem. How can she count things that cannot be counted? Today, we discover the secret: \cle{partitive expressions} like \eng{a bag of rice}, \eng{a loaf of bread} and \eng{a bottle of water}.

\textbf{Problématique.} En français, nous disons naturellement \emph{du pain}, \emph{de l'huile}, \emph{un riz} n'existe pas non plus chez nous… et pourtant, lorsque nous parlons anglais, nous avons tendance à « compter » ces mots : \eng{two breads}, \eng{three oils}, \eng{an information}. Pourquoi l'anglais refuse-t-il ces formes ? Quel mot faut-il placer devant \eng{bread}, \eng{oil}, \eng{rice} pour exprimer une quantité précise ? C'est toute la question des \cle{expressions partitives} que nous allons observer aujourd'hui.

\courssub{Écoutons et lisons : au marché}

\begin{textebox}[At the market — a dialogue]
\textbf{Sarah:} Good morning, Madam! I need a bag of rice, a bunch of bananas, a tin of tomatoes and a bottle of palm oil.

\textbf{Trader:} Here you are. Anything else?

\textbf{Sarah:} Yes, a piece of soap and two loaves of bread, please. Oh, and can I also have a cup of groundnuts and a bar of chocolate for my little brother?

\textbf{Trader:} Of course! Is that all?

\textbf{Sarah:} Let me check my list… a bag of rice, a bunch of bananas, a tin of tomatoes, a bottle of palm oil, a piece of soap, two loaves of bread, a cup of groundnuts, a bar of chocolate. Yes, that's everything. How much is it?

\textbf{Trader:} Five thousand francs, my dear. You speak very good English!

\textbf{Sarah:} Thank you, Madam. My uncle in Bamenda taught me that in English, you never say “two breads” — you say “two loaves of bread”!

\textbf{Trader:} Your uncle is right. Have a nice day!
\end{textebox}

\textbf{Glossaire.}

\begin{center}
\begin{tabularx}{\linewidth}{|X|X|X|}
\hline
\rowcolor{popblueL} Mot anglais & Nature & Traduction française \\
\hline
trader & nom & commerçante, vendeuse \\
\hline
a bunch of bananas & expression & un régime de bananes \\
\hline
a tin of tomatoes & expression & une boîte (conserve) de tomates \\
\hline
palm oil & nom & huile de palme \\
\hline
a piece of soap & expression & un morceau de savon (une savonnette) \\
\hline
a loaf of bread & expression & une miche de pain, un pain \\
\hline
groundnuts & nom pluriel & arachides \\
\hline
a bar of chocolate & expression & une tablette de chocolat \\
\hline
\end{tabularx}
\end{center}

\textbf{Questions de compréhension (reading).}
\begin{enumerate}
\item Where does the dialogue take place?
\item Name three things Sarah buys and the quantity of each.
\item What did Sarah's uncle teach her?
\item In your own words: why does Sarah say \eng{two loaves of bread} and not \eng{two breads}?
\end{enumerate}

\courssub{Repérage : les groupes « quantité + of + nom »}

Reprenons le dialogue et soulignons les groupes qui expriment une \emph{quantité} :

\begin{exemplebox}[Ce que Sarah achète]
\begin{itemize}
\item \eng{a \textbf{bag} of rice} — un sac de riz
\item \eng{a \textbf{bunch} of bananas} — un régime de bananes
\item \eng{a \textbf{tin} of tomatoes} — une boîte de tomates
\item \eng{a \textbf{bottle} of palm oil} — une bouteille d'huile de palme
\item \eng{a \textbf{piece} of soap} — un morceau de savon
\item \eng{two \textbf{loaves} of bread} — deux pains
\item \eng{a \textbf{cup} of groundnuts} — une tasse d'arachides
\item \eng{a \textbf{bar} of chocolate} — une tablette de chocolat
\end{itemize}
\end{exemplebox}

\textbf{Observation guidée.} Réponds aux questions suivantes :
\begin{enumerate}
\item Quel petit mot revient dans \emph{tous} ces groupes ? (\eng{of})
\item Le mot avant \eng{of} désigne-t-il la marchandise elle-même, ou son contenant, sa portion, sa forme ?
\item Comment met-on ces groupes au pluriel ? Observe : \eng{a loaf of bread} $\rightarrow$ \eng{two loaves of bread}. Quel mot prend la marque du pluriel : \eng{loaf} ou \eng{bread} ?
\end{enumerate}

\begin{aretenir}[Réponses à l'observation guidée]
\begin{enumerate}
\item Le mot invariant \eng{of} (de) relie le mot de quantité au nom.
\item Le mot avant \eng{of} (\eng{bag, bunch, tin, bottle, piece, loaf, cup, bar}) désigne le \textbf{contenant}, la \textbf{portion}, la \textbf{forme} ou le \textbf{groupe} — jamais la matière elle-même.
\item C'est le \textbf{mot de quantité} (\eng{loaf} $\rightarrow$ \eng{loaves}) qui prend le pluriel. Le nom indénombrable (\eng{bread}) reste \textbf{invariable} au singulier.
\end{enumerate}
\end{aretenir}

\begin{attention}[La question piège : peut-on dire “two breads” ?]
Non ! En anglais, on ne dit jamais \eng{two breads}, \eng{three oils}, \eng{a rice}, \eng{an information}, \eng{a money}. Les noms \eng{bread}, \eng{oil}, \eng{rice}, \eng{water}, \eng{money}, \eng{information}, \eng{advice}, \eng{news}, \eng{furniture}, \eng{homework} sont \cle{indénombrables} (\emph{uncountable nouns}) : ils n'ont pas de pluriel et ne peuvent pas être précédés directement de \eng{a/an} ni d'un nombre. Pour les « compter », l'anglais compte leur \emph{contenant}, leur \emph{portion} ou leur \emph{forme} : \eng{two loaves of bread}, \eng{three bottles of oil}, \eng{a bag of rice}, \eng{a piece of information}.
\end{attention}

\courssub{Le constat : noms dénombrables et noms indénombrables}

\begin{definition}[Countable nouns et uncountable nouns]
Un \cle{nom dénombrable} (\emph{countable noun}) désigne quelque chose qu'on peut compter une par une : il a un singulier et un pluriel, et on peut le précéder de \eng{a/an} ou d'un nombre : \eng{a mango, two mangoes} ; \eng{a banana, three bananas}.

Un \cle{nom indénombrable} (\emph{uncountable noun} ou \emph{mass noun}) désigne une matière, un liquide, une notion abstraite qu'on ne peut pas compter directement : \eng{rice, water, bread, oil, money, information, advice, news, furniture, homework}. Il n'a \textbf{pas de pluriel} et ne prend \textbf{jamais} \eng{a/an} ni un nombre devant lui.
\end{definition}

\begin{center}
\begin{tabularx}{\linewidth}{|X|X|X|}
\hline
\rowcolor{popblueL} Français & English & Remarque \\
\hline
une mangue / deux mangues & a mango / two mangoes & dénombrable : pluriel normal \\
\hline
du riz / un sac de riz & rice / a bag of rice & indénombrable : jamais \eng{a rice} \\
\hline
de l'eau / trois bouteilles d'eau & water / three bottles of water & indénombrable : jamais \eng{three waters} \\
\hline
un pain / deux pains & bread / two loaves of bread & indénombrable : jamais \eng{two breads} \\
\hline
un conseil & a piece of advice & \eng{advice} est indénombrable ! \\
\hline
une information & a piece of information & \eng{information} est indénombrable ! \\
\hline
de l'argent & money & jamais \eng{a money} ni \eng{monies} (courant) \\
\hline
un meuble / des meubles & furniture & indénombrable : \eng{a piece of furniture} \\
\hline
un devoir / des devoirs & homework & indénombrable : jamais \eng{a homework} \\
\hline
\end{tabularx}
\end{center}

\begin{definition}[Partitive expression]
A \cle{partitive expression} is a word or phrase (\eng{a cup of}, \eng{a piece of}, \eng{a bottle of}) placed \emph{before} an uncountable noun to express a quantity or a portion of it. Structure : \textbf{(nombre +) mot de quantité + of + nom indénombrable}.
\end{definition}

\begin{exemplebox}[Exemples traduits]
\eng{a bottle of water} = une bouteille d'eau

\eng{a piece of advice} = un conseil

\eng{a bunch of keys} = un trousseau de clés

\eng{a loaf of bread} = un pain (une miche de pain)

\eng{two cups of tea} = deux tasses de thé

\eng{a piece of news} = une nouvelle
\end{exemplebox}

\coursec{La règle : structure et formation}

\begin{propriete}[Structure de l'expression partitive]
Pour quantifier un nom indénombrable, on utilise la structure fixe :
\[
\textbf{Quantificateur (singulier ou pluriel) + of + Nom indénombrable (toujours au singulier)}
\]
\begin{itemize}
\item Le \textbf{quantificateur} (mot de quantité) s'accorde en nombre : \eng{a bottle} $\rightarrow$ \eng{two bottles}.
\item La préposition \textbf{\eng{of}} est \textbf{obligatoire} : jamais \eng{a cup tea}, toujours \eng{a cup \textbf{of} tea}.
\item Le \textbf{nom indénombrable} ne change \textbf{jamais de forme} (pas de \eng{-s}).
\end{itemize}
\end{propriete}

\begin{exemplebox}[Application de la règle]
\begin{tabular}{lll}
Singulier & Pluriel & Traduction \\
\hline
\eng{a loaf of bread} & \eng{two loaves of bread} & un pain / deux pains \\
\eng{a bottle of oil} & \eng{three bottles of oil} & une bouteille d'huile / trois bouteilles d'huile \\
\eng{a piece of advice} & \eng{several pieces of advice} & un conseil / plusieurs conseils \\
\eng{a grain of rice} & \eng{many grains of rice} & un grain de riz / beaucoup de grains de riz \\
\end{tabular}
\end{exemplebox}

\begin{attention}[Piège : l'accord du verbe]
Le verbe s'accorde avec le \textbf{quantificateur} (le sujet grammatical), pas avec le nom indénombrable.
\eng{A bottle of water \textbf{is} on the table.} (singulier)
\eng{Two bottles of water \textbf{are} on the table.} (pluriel)
\end{attention}

\coursec{Les principaux partitifs par catégories}

\courssub{Tableau de référence : quantificateurs courants classés par notion}

\begin{center}
\begin{tabularx}{\linewidth}{|X|X|X|}
\hline
\rowcolor{popblueL} Catégorie & Expressions partitives (Quantificateur + of) & Exemples typiques (Noms indénombrables) \\
\hline
\textbf{Contenants} & \eng{a bottle of}, \eng{a tin/can of}, \eng{a jar of}, \eng{a packet/pack of}, \eng{a box of}, \eng{a bag of}, \eng{a cup of}, \eng{a glass of}, \eng{a mug of}, \eng{a bowl of}, \eng{a pot of} & \eng{water, oil, milk, jam, rice, sugar, flour, tea, coffee, soup, yoghurt} \\
\hline
\textbf{Portions / Morceaux} & \eng{a piece of}, \eng{a slice of}, \eng{a chunk of}, \eng{a lump of}, \eng{a bit of}, \eng{a segment of} & \eng{bread, cake, meat, cheese, chocolate, advice, information, news, furniture, wood, coal, sugar} \\
\hline
\textbf{Formes spécifiques} & \eng{a loaf of}, \eng{a bar of}, \eng{a sheet of}, \eng{a roll of}, \eng{a tube of}, \eng{a stick of}, \eng{a block of} & \eng{bread, chocolate, soap, paper, tape, toothpaste, butter, ice, cement} \\
\hline
\textbf{Groupes / Collections} & \eng{a bunch of}, \eng{a cluster of}, \eng{a bunch of}, \eng{a crowd of}, \eng{a flock of}, \eng{a herd of}, \eng{a pile of}, \eng{a heap of}, \eng{a stack of}, \eng{a set of} & \eng{bananas, grapes, keys, flowers, people, birds, cattle, clothes, books, plates, tools} \\
\hline
\textbf{Mesures / Quantités vagues} & \eng{a drop of}, \eng{a pinch of}, \eng{a spoonful of}, \eng{a handful of}, \eng{a load of}, \eng{a lot of}, \eng{plenty of}, \eng{a great deal of} & \eng{water, oil, salt, sugar, flour, sand, money, time, work, fun} \\
\hline
\end{tabularx}
\end{center}

\begin{savaistu}[Partitifs et noms dénombrables au pluriel]
Certains partitifs (surtout \textbf{groupes} et \textbf{mesures}) s'utilisent aussi devant des noms dénombrables au pluriel : \eng{a bunch of bananas}, \eng{a crowd of people}, \eng{a handful of coins}, \eng{a lot of mangoes}. Le nom qui suit \eng{of} est alors au pluriel, mais la règle de l'accord du verbe sur le quantificateur reste valable : \eng{A bunch of bananas \textbf{is} on the table.}
\end{savaistu}

\coursec{Comparaison français / anglais et cas particuliers}

\courssub{Tableau comparatif : pièges à éviter pour les francophones}

\begin{center}
\begin{tabularx}{\linewidth}{|X|X|X|}
\hline
\rowcolor{popblueL} Français (structure) & Anglais correct & Erreur fréquente (interférence) \\
\hline
du pain / un pain & \eng{some bread} / \eng{a loaf of bread} & \eng{a bread}, \eng{two breads} \\
\hline
de l'eau / une eau & \eng{some water} / \eng{a bottle of water} & \eng{a water}, \eng{two waters} \\
\hline
du riz / un riz & \eng{some rice} / \eng{a bag of rice} & \eng{a rice}, \eng{three rices} \\
\hline
un conseil & \eng{a piece of advice} & \eng{an advice} \\
\hline
des conseils & \eng{some advice} / \eng{pieces of advice} & \eng{advices} \\
\hline
une information & \eng{a piece of information} & \eng{an information} \\
\hline
des informations & \eng{some information} / \eng{pieces of information} & \eng{informations} \\
\hline
un meuble & \eng{a piece of furniture} & \eng{a furniture} \\
\hline
un cheveu & \eng{a hair} (dénombrable !) & \eng{a piece of hair} (possible mais \eng{a hair} est standard) \\
\hline
de la chance & \eng{luck} (indénombrable) & \eng{a luck} \\
\hline
un travail / des travaux & \eng{work} (indénombrable) / \eng{a job} (dénombrable) & \eng{a work} (sens « œuvre d'art » seulement) \\
\hline
\end{tabularx}
\end{center}

\courssub{Cas particuliers et exceptions notables}

\begin{propriete}[Noms à double nature : sens dénombrable vs indénombrable]
Certains noms changent de sens selon qu'ils sont utilisés comme dénombrables ou indénombrables. L'expression partitive s'applique \textbf{uniquement au sens indénombrable}.
\end{propriete}

\begin{center}
\begin{tabularx}{\linewidth}{|X|X|X|}
\hline
\rowcolor{popblueL} Nom & Sens indénombrable (partitif requis) & Sens dénombrable (a/an + pluriel -s) \\
\hline
\eng{chocolate} & \eng{a bar of chocolate} (matière) & \eng{a chocolate} (bonbon individuel), \eng{two chocolates} \\
\hline
\eng{coffee / tea} & \eng{a cup of coffee} (liquide) & \eng{a coffee} (une tasse de café au resto), \eng{two coffees} \\
\hline
\eng{glass} & \eng{a sheet of glass} (matière) & \eng{a glass} (verre à boire), \eng{glasses} (lunettes) \\
\hline
\eng{paper} & \eng{a sheet of paper} (matière) & \eng{a paper} (journal / document / épreuve), \eng{papers} \\
\hline
\eng{iron} & \eng{a bar of iron} (métal) & \eng{an iron} (fer à repasser) \\
\hline
\eng{hair} & \eng{hair} (chevelure, masse) & \eng{a hair} (un cheveu unique) \\
\hline
\eng{time} & \eng{time} (durée, temps) & \eng{a time} (une fois / occasion) \\
\hline
\end{tabularx}
\end{center}

\begin{attention}[Faux amis « partitifs » : \eng{a piece of} $\neq$ \emph{un morceau de} toujours]
\begin{itemize}
\item \eng{A piece of advice / information / news / furniture / luggage / equipment} = un conseil / une information / une nouvelle / un meuble / un bagage / un équipement. \textbf{Pas de traduction « morceau »}.
\item \eng{A piece of cake / meat / bread / chocolate} = un morceau de gâteau / viande / pain / chocolat. Ici « morceau » est correct.
\end{itemize}
\end{attention}

\begin{savaistu}[Le cas de \eng{news}, \eng{mathematics}, \eng{physics}, \eng{politics}]
En anglais, \eng{news} est \textbf{indénombrable} malgré son « s » final ! On dit \eng{a piece of news} (une nouvelle) et \eng{the news is good} (singulier), jamais \eng{a new} ni \eng{the news are good}. De même, \eng{mathematics}, \eng{physics}, \eng{economics}, \eng{politics}, \eng{gymnastics} se construisent au singulier : \eng{Mathematics is difficult.} À l'inverse, \eng{trousers}, \eng{scissors}, \eng{glasses} (lunettes), \eng{jeans}, \eng{shorts} sont toujours au pluriel : on dit \eng{a pair of trousers} (un pantalon), \eng{a pair of scissors} (une paire de ciseaux) — encore un partitif (\eng{pair of}) !
\end{savaistu}

\coursec{Emplois en contexte et cas particuliers}

\courssub{Choisir le bon quantificateur : méthode}

\begin{methode}[Comment choisir l'expression partitive appropriée ?]
\begin{enumerate}
\item \textbf{Identifie le nom} : est-il indénombrable (matière, liquide, abstraction) ou dénombrable ?
\item \textbf{Si indénombrable} : détermine la \textbf{nature de la quantité} voulue :
\begin{itemize}
\item Contenant standard $\rightarrow$ \eng{bottle, tin, jar, packet, bag, box, cup, glass, bowl}.
\item Portion découpée $\rightarrow$ \eng{slice} (tranche), \eng{piece} (morceau), \eng{chunk} (gros morceau).
\item Forme spécifique $\rightarrow$ \eng{loaf} (pain), \eng{bar} (tablette, barre), \eng{sheet} (feuille), \eng{roll} (rouleau), \eng{tube}.
\item Groupe naturel $\rightarrow$ \eng{bunch} (régime, bouquet), \eng{cluster} (grappe).
\item Mesure de cuisine $\rightarrow$ \eng{spoonful, pinch, drop, handful}.
\end{itemize}
\item \textbf{Si dénombrable pluriel} (groupe) : \eng{bunch of bananas}, \eng{crowd of people}, \eng{pile of books}.
\item \textbf{Vérifie l'accord} : pluriel sur le quantificateur (\eng{two slices}), \eng{of} inchangé, nom indénombrable au singulier.
\end{enumerate}
\end{methode}

\courssub{Exemples en contexte camerounais et international}

\begin{exemplebox}[Au restaurant / snack à Douala]
\textbf{Client:} I'd like \eng{a bottle of Mutzig}, \eng{a plate of ndolé} and \eng{a piece of grilled fish}, please. \\
\textbf{Waiter:} Would you like \eng{a glass of water} or \eng{a cup of tea} with that? \\
\textbf{Client:} Just \eng{a slice of pineapple} for dessert.
\end{exemplebox}

\begin{exemplebox}[À l'école / administration]
\begin{itemize}
\item The teacher gave us \eng{a piece of advice} about the Baccalaureate.
\item We need \eng{a sheet of paper} for the exam.
\item There is \eng{a lot of homework} this weekend.
\item The principal made \eng{an important announcement} (dénombrable) / gave \eng{us some important information}.
\end{itemize}
\end{exemplebox}

\begin{exemplebox}[Courses au supermarché (Buea / Yaoundé)]
\eng{Shopping list:} \eng{two bags of rice, a tin of sardines, a jar of mayonnaise, a packet of spaghetti, a loaf of bread, a bar of soap, a tube of toothpaste, a bunch of plantains, a bottle of palm oil, a kilo of tomatoes.}
\emph{Note :} \eng{a kilo of} fonctionne comme un partitif de mesure (quantificateur = \eng{kilo}).
\end{exemplebox}

\coursec{Exercices d'application et de transformation}

\begin{exoresolu}[Corrige la liste de courses de Sarah]
Énoncé. Sarah a écrit : \eng{I need two rices, three breads, one water and four soaps.} Corrige sa liste à l'aide d'expressions partitives adaptées.
\tcblower
Solution détaillée. Les quatre noms \eng{rice}, \eng{bread}, \eng{water} et \eng{soap} sont indénombrables : on ne peut pas les précéder directement d'un nombre. On intercale un mot de quantité, qui lui porte le pluriel :

\eng{I need two \textbf{bags} of rice, three \textbf{loaves} of bread, one \textbf{bottle} of water and four \textbf{bars} of soap.}

(J'ai besoin de deux sacs de riz, de trois pains, d'une bouteille d'eau et de quatre savonnettes.)

Remarque : \eng{loaf} a un pluriel irrégulier : \eng{loaves} (comme \eng{leaf/leaves}, \eng{knife/knives}).
\end{exoresolu}

\begin{exoresolu}[Transformation : du dénombrable au partitif]
Énoncé. Réécris les phrases en utilisant une expression partitive pour le mot souligné.
\begin{enumerate}
\item She bought \underline{two chocolates}. (une tablette)
\item He gave me \underline{an advice}.
\item We need \underline{three breads} for breakfast.
\item There are \underline{many furnitures} in the room.
\item She drank \underline{a water}.
\end{enumerate}
\tcblower
Solution détaillée.
\begin{enumerate}
\item She bought \eng{a bar of chocolate}. (\eng{chocolate} indénombrable = matière)
\item He gave me \eng{a piece of advice}. (\eng{advice} indénombrable)
\item We need \eng{three loaves of bread} for breakfast. (\eng{bread} indénombrable)
\item There is \eng{a lot of furniture} (ou \eng{many pieces of furniture}) in the room. (\eng{furniture} indénombrable, verbe au singulier !)
\item She drank \eng{a glass of water} (ou \eng{a bottle of water}). (\eng{water} indénombrable)
\end{enumerate}
\end{exoresolu}

\begin{attention}[Erreurs fréquentes des francophones — Récapitulatif]
\begin{itemize}
\item Ne traduis jamais \emph{un pain} par \eng{a bread} : dis \eng{a loaf of bread} ou \eng{some bread}.
\item \emph{Des informations} ne se dit pas \eng{informations} : \eng{information} est toujours singulier ; dis \eng{some information} ou \eng{pieces of information}.
\item \emph{Un conseil} = \eng{a piece of advice}, jamais \eng{an advice}.
\item N'oublie pas le \eng{of} : on dit \eng{a cup of tea}, jamais \eng{a cup tea}.
\item Le pluriel se met sur le mot de quantité, pas sur le nom indénombrable : \eng{two bottles of oil}, jamais \eng{two bottle of oils}.
\item Verbe au singulier avec quantificateur singulier : \eng{A lot of work \textbf{is} waiting.} / \eng{Lots of people \textbf{are} waiting.}
\end{itemize}
\end{attention}

\begin{experience}[Mini-activité orale : At the market — Role-play]
Par groupes de deux : l'élève A est le client, l'élève B est le vendeur au marché (Mokolo, Marché Central, ou Marché de Buea). Le client demande cinq produits en utilisant obligatoirement une expression partitive \textbf{différente} pour chacun (ex. \eng{a bag of…}, \eng{a tin of…}, \eng{a bunch of…}, \eng{a bottle of…}, \eng{a piece of…}, \eng{a loaf of…}, \eng{a bar of…}). Le vendeur répond : \eng{Here you are. Anything else?} Puis échangez les rôles. Le reste de la classe note les expressions entendues et vérifie qu'aucun nom indénombrable n'est « compté » directement (pas de \eng{two rices}, \eng{three oils}).
\end{experience}

\begin{experience}[Production écrite guidée : My shopping list]
Rédige une liste de courses en anglais pour préparer un repas typique camerounais (ex. \emph{ndolé}, \emph{eru}, \emph{riz sauce arachide}, \emph{poulet DG}). Utilise au moins \textbf{huit} expressions partitives différentes (contenants, portions, formes, groupes, mesures). Présente ta liste à l'oral en justifiant le choix de chaque quantificateur.
\end{experience}

\begin{exercice}[À toi de jouer — Entraînement complet]
\begin{enumerate}
\item \textbf{Classification.} Classe ces noms en deux colonnes : \emph{Countable} ou \emph{Uncountable} : \eng{mango, rice, money, banana, information, bread, key, water, advice, tomato, furniture, homework, coin, news, luggage}.
\item \textbf{Complétion.} Complète avec le quantificateur approprié (choisis dans la boîte) : \eng{loaf, bar, piece, bottle, bunch, cup, slice, sheet, tin, bag}.
\begin{itemize}
\item[a.] a \ldots of bread
\item[b.] a \ldots of keys
\item[c.] a \ldots of palm oil
\item[d.] a \ldots of news
\item[e.] two \ldots of tea
\item[f.] a \ldots of chocolate
\item[g.] a \ldots of paper
\item[h.] a \ldots of tomatoes (conserve)
\item[i.] a \ldots of bananas
\item[j.] a \ldots of rice
\end{itemize}
\item \textbf{Traduction.} Traduis en anglais :
\begin{enumerate}
\item Je voudrais un morceau de savon, deux bouteilles d'huile et un régime de bananes, s'il vous plaît.
\item Il n'y a plus de pain. Achète trois miches à la boulangerie.
\item Elle m'a donné un bon conseil et une information importante.
\item Nous avons mangé une tranche de gâteau et bu une tasse de café.
\item Il y a beaucoup de monde au marché aujourd'hui.
\end{enumerate}
\item \textbf{Correction d'erreurs.} Corrige les phrases suivantes (une erreur par phrase) :
\begin{enumerate}
\item I need two breads for the sandwich.
\item She gave me an advice.
\item There are many furnitures in the classroom.
\item Can I have a water, please?
\item He ate two chocolates (meaning: two bars).
\item A bunch of flowers are on the table.
\item I bought a tin of tomato.
\item She needs a piece of luggages.
\end{enumerate}
\end{enumerate}
\end{exercice}

\begin{corrige}[Exercice « À toi de jouer »]
\begin{enumerate}
\item \textbf{Countable :} \eng{mango, banana, key, tomato, coin}. \textbf{Uncountable :} \eng{rice, money, information, bread, water, advice, furniture, homework, news, luggage}.
\item \textbf{Complétion :}
a. \eng{loaf} \quad b. \eng{bunch} \quad c. \eng{bottle} \quad d. \eng{piece} \quad e. \eng{cups} \quad f. \eng{bar} \quad g. \eng{sheet} \quad h. \eng{tin} \quad i. \eng{bunch} \quad j. \eng{bag}
\item \textbf{Traduction :}
\begin{enumerate}
\item \eng{I would like a piece of soap, two bottles of oil and a bunch of bananas, please.}
\item \eng{There is no more bread. Buy three loaves at the bakery.}
\item \eng{She gave me a piece of advice and a piece of important information.}
\item \eng{We ate a slice of cake and drank a cup of coffee.}
\item \eng{There is a crowd of people at the market today.} (ou \eng{There are a lot of people...})
\end{enumerate}
\item \textbf{Correction d'erreurs :}
\begin{enumerate}
\item \eng{I need two \textbf{loaves of} bread for the sandwich.}
\item \eng{She gave me \textbf{a piece of} advice.}
\item \eng{There is \textbf{much furniture} (ou \textbf{many pieces of furniture}) in the classroom.}
\item \eng{Can I have \textbf{a glass of} water, please?}
\item \eng{He ate two \textbf{bars of} chocolate.}
\item \eng{A bunch of flowers \textbf{is} on the table.} (sujet = \eng{a bunch}, singulier)
\item \eng{I bought a tin of \textbf{tomatoes}.} (pluriel après \eng{of} pour noms dénombrables en conserve)
\item \eng{She needs a piece of \textbf{luggage}.} (\eng{luggage} indénombrable, pas de \eng{s})
\end{enumerate}
\end{enumerate}
\end{corrige}

\begin{savaistu}[Pour aller plus loin : \eng{some} / \eng{any} vs partitifs précis]
Quand la quantité est indéterminée ou non précisée, on utilise \eng{some} (affirmatif) / \eng{any} (négatif/interrogatif) + nom indénombrable \textbf{sans} \eng{of} : \eng{I need some rice.} \eng{Do you have any water?} Les expressions partitives (\eng{a bag of}, \eng{a bottle of}) servent à \textbf{préciser} la quantité (mesure, contenant, portion). C'est la différence entre « du riz » (quantité vague) et « un sac de riz » (quantité précise).
\end{savaistu}

\coursec{Observation et découverte}

\begin{textebox}[Dialogue au marché Mokolo (Yaoundé)]
\textbf{Context:} Mrs. Atangana is shopping at Mokolo market. She talks to a seller and a shop assistant.

\textbf{Seller:} Good morning, Madam. What can I get for you today?

\textbf{Mrs. Atangana:} Good morning. I need \eng{a bag of rice}, \eng{two bottles of palm oil} and \eng{a loaf of bread}.

\textbf{Seller:} Certainly. Here is the rice — fifty kilogrammes. And here are \eng{two bottles of red palm oil}. Do you need anything else?

\textbf{Mrs. Atangana:} Yes, my son needs \eng{a piece of chalk} for school and \eng{a box of chalk} for his teacher. Also, give me \eng{a bunch of plantains} and \eng{three bars of soap}.

\textbf{Shop assistant (at the buvette next door):} Welcome! What would you like?

\textbf{Mrs. Atangana:} \eng{Two cups of tea}, please. And \eng{a glass of water}.

\textbf{Shop assistant:} Here you are. \eng{Two teas} for you. Anything else?

\textbf{Mrs. Atangana:} No, thank you. How much is everything?
\end{textebox}

\begin{experience}[Observation guidée]
Travaillez par paires. Relisez le dialogue et répondez aux questions ci-dessous.
\begin{enumerate}
\item Relevez toutes les expressions qui indiquent une \textbf{quantité} ou une \textbf{portion} (ex. : \eng{a bag of}, \eng{two bottles of}). Combien en trouvez-vous ?
\item Dans chaque expression relevée, identifiez le \textbf{mot qui précède \eng{of}} (ex. : \eng{bag}, \eng{bottles}). Comment appelle-t-on ce mot ? (Indice : c'est un « partitif »).
\item Observez le nom qui suit \eng{of} (ex. : \eng{rice}, \eng{palm oil}, \eng{chalk}). Ces noms prennent-ils un \eng{-s} au pluriel ? Pourquoi ?
\item Comparez : \eng{two bottles of palm oil} (dialogue vendeur) et \eng{Two teas} (dialogue buvette). Quelle différence remarquez-vous sur la structure ? Laquelle est la forme grammaticale complète ?
\item Regardez l'accord du verbe dans : « \eng{Here \textbf{are} two bottles…} » et « \eng{Two teas \textbf{for} you} » (phrase elliptique). Avec quoi le verbe s'accorde-t-il ?
\end{enumerate}
\end{experience}

\coursec{La règle : structure et formation}

Dans la section précédente, nous avons observé des phrases comme \eng{a cup of tea}, \eng{a piece of chalk} ou \eng{two bottles of oil} dans le dialogue au marché. Nous avons remarqué que les noms \eng{tea}, \eng{chalk} et \eng{oil} ne prennent jamais de \eng{-s}, alors que \eng{cup}, \eng{piece} et \eng{bottle} peuvent en prendre un. Il est temps maintenant de dégager la règle générale qui gouverne toutes ces expressions : la structure partitive.

\courssub{La structure générale : quantity + partitive + of + uncountable noun}

\begin{definition}[Le nom indénombrable (uncountable noun)]
Un nom \cle{indénombrable} (ou \eng{uncountable noun} / \eng{mass noun}) désigne une matière, un liquide, une substance ou une notion abstraite qu'on ne peut pas compter unité par unité : \eng{water} (eau), \eng{rice} (riz), \eng{bread} (pain), \eng{money} (argent), \eng{information} (information), \eng{furniture} (meubles), \eng{advice} (conseil), \eng{chalk} (craie). En anglais, ces noms n'ont \textbf{pas de pluriel} (*\eng{rices}, *\eng{informations} sont faux) et ne s'emploient \textbf{jamais avec} \eng{a} ou \eng{an} directement (*\eng{a bread}, *\eng{an advice} sont faux).
\end{definition}

Pour compter ou quantifier un indénombrable, l'anglais passe obligatoirement par un \cle{partitif} (ou \eng{partitive noun} / \eng{measure word}), c'est-à-dire un mot qui indique un contenant, une portion, une unité de mesure ou un groupe : \eng{cup}, \eng{piece}, \eng{bottle}, \eng{bunch}, \eng{loaf}, \eng{sheet}, \eng{bar}, \eng{slice}, \eng{bag}, \eng{glass}, etc.

\begin{propriete}[Règle de formation du partitif]
La structure partitive se construit ainsi :

\textbf{QUANTITY + PARTITIVE + OF + UNCOUNTABLE NOUN}

\begin{itemize}
\item Au \textbf{singulier} : \eng{a} (ou \eng{an}, \eng{one}) + partitif (singulier) + \eng{of} + nom indénombrable (invariable).\\
\eng{a cup of tea} « une tasse de thé » ; \eng{a piece of chalk} « un morceau de craie » ; \eng{a bottle of oil} « une bouteille d'huile » ; \eng{a loaf of bread} « un pain / une miche de pain ».
\item Au \textbf{pluriel} : nombre (\eng{two}, \eng{three}, \eng{many}…) + partitif\textbf{-s} + \eng{of} + nom indénombrable (qui reste \textbf{invariable}).\\
\eng{two cups of tea} « deux tasses de thé » ; \eng{three pieces of chalk} « trois morceaux de craie » ; \eng{several bags of rice} « plusieurs sacs de riz ».
\end{itemize}

\textbf{Règle d'or :} Seul le partitif prend la marque du pluriel (\eng{-s} ou \eng{-es}), jamais le nom indénombrable. On dit \eng{three cups of tea}, jamais *\eng{three teas} (sauf cas particulier « commande au café », voir plus bas).
\end{propriete}

\begin{popfigure}
\begin{tikzpicture}[node distance=0.4cm, every node/.style={font=\small}]
\node[draw=popblue, fill=popblueL, rounded corners, minimum width=2.2cm, minimum height=1.4cm, align=center] (n1) {\textbf{Quantity}\\ \eng{two}};
\node[draw=poporange, fill=poporangeL, rounded corners, minimum width=2.4cm, minimum height=1.4cm, align=center, right=of n1] (n2) {\textbf{Partitive}\\ \eng{bottles}};
\node[draw=poppurple, fill=poppurpleL, rounded corners, minimum width=1.2cm, minimum height=1.4cm, align=center, right=of n2] (n3) {\textbf{Link}\\ \eng{of}};
\node[draw=popgreen, fill=popgreenL, rounded corners, minimum width=3.2cm, minimum height=1.4cm, align=center, right=of n3] (n4) {\textbf{Uncountable Noun}\\ \eng{palm oil}};
\draw[-{Stealth}, thick, popdark] (n1) -- (n2);
\draw[-{Stealth}, thick, popdark] (n2) -- (n3);
\draw[-{Stealth}, thick, popdark] (n3) -- (n4);
\node[below=0.35cm of n2, align=center, text=popmuted, font=\small] {\eng{two bottles of palm oil} = « deux bouteilles d'huile de palme »};
\node[below=0.8cm of n4, align=center, text=popgreen, font=\small] {\eng{palm oil} : invariable, pas de \eng{-s}};
\end{tikzpicture}
\legende{La structure partitive : la quantité porte sur le partitif, le nom indénombrable reste invariable.}
\end{popfigure}

\courssub{Du singulier au pluriel : le \eng{-s} va sur le partitif}

C'est le point capital de la leçon. Quand on passe au pluriel, la marque du pluriel se déplace sur le partitif, et le nom indénombrable ne change jamais de forme. Attention aux partitifs à pluriel irrégulier.

\begin{exemplebox}[Singulier et pluriel : exemples traduits]
\begin{itemize}
\item \eng{a glass of water} = « un verre d'eau » ; \eng{two glasses of water} = « deux verres d'eau » (\eng{glass} $\rightarrow$ \eng{glasses}).
\item \eng{a sheet of paper} = « une feuille de papier » ; \eng{five sheets of paper} = « cinq feuilles de papier ».
\item \eng{a loaf of bread} = « un pain / une miche de pain » ; \eng{two loaves of bread} = « deux pains » (\eng{loaf} $\rightarrow$ \eng{loaves}).
\item \eng{a bar of soap} = « une barre de savon » ; \eng{three bars of soap} = « trois barres de savon ».
\item \eng{a piece of advice} = « un conseil » ; \eng{several pieces of advice} = « plusieurs conseils ».
\item \eng{a bag of rice} = « un sac de riz » ; \eng{ten bags of rice} = « dix sacs de riz ».
\item \eng{a bottle of palm oil} = « une bouteille d'huile de palme » ; \eng{two bottles of palm oil} = « deux bouteilles d'huile de palme ».
\end{itemize}
\end{exemplebox}

\begin{attention}[Erreur fréquente des francophones : le pluriel interdit sur l'indénombrable]
En français, on peut parfois avoir l'impression de « compter » la matière : « deux riz », « trois meubles », « des informations ». En anglais, c'est \textbf{impossible} avec un indénombrable. Les formes suivantes sont toutes \textbf{fausses} :
\begin{itemize}
\item *\eng{two rices} $\rightarrow$ \eng{two bags of rice} / \eng{two kilos of rice} « deux sacs / kilos de riz »
\item *\eng{three furnitures} $\rightarrow$ \eng{three pieces of furniture} « trois meubles »
\item *\eng{many informations} $\rightarrow$ \eng{much information} ou \eng{many pieces of information} « beaucoup d'informations »
\item *\eng{an advice} $\rightarrow$ \eng{a piece of advice} « un conseil »
\item *\eng{a bread} $\rightarrow$ \eng{a loaf of bread} ou \eng{a slice of bread} « un pain / une tranche de pain »
\item *\eng{two chalks} $\rightarrow$ \eng{two pieces of chalk} / \eng{two boxes of chalk} « deux morceaux / boîtes de craie »
\end{itemize}
Retenez : un nom indénombrable anglais ne prend \textbf{jamais} \eng{a/an} ni \eng{-s}.
\end{attention}

\begin{popfigure}
\begin{tikzpicture}[node distance=1.8cm, every node/.style={font=\small}]
\node[draw=popblue, fill=popblueL, rounded corners, minimum width=4.8cm, minimum height=1.3cm, align=center] (sing) {\textbf{Singular}\\ \eng{a loaf of bread}};
\node[draw=popgreen, fill=popgreenL, rounded corners, minimum width=4.8cm, minimum height=1.3cm, align=center, right=of sing] (plur) {\textbf{Plural}\\ \eng{two \textbf{loaves} of bread}};
\draw[-{Stealth}, very thick, poporange] (sing) -- node[above, align=center, text=popdark, font=\small]{le \eng{-s} va sur\\ \eng{loaf} $\rightarrow$ \eng{loaves}} (plur);
\node[below=0.6cm of sing, align=center, text=poppink, font=\small] {\eng{bread} : jamais de \eng{-s}\\ *\eng{breads} est faux};
\node[below=0.6cm of plur, align=center, text=popgreen, font=\small] {\eng{bread} reste invariable\\ \eng{two loaves of bread}};
\end{tikzpicture}
\legende{Au pluriel, la marque du pluriel se fixe sur le partitif (\eng{loaf} $\rightarrow$ \eng{loaves}), jamais sur le nom indénombrable (\eng{bread}).}
\end{popfigure}

\courssub{L'accord du verbe : matière ou contenants ?}

Avec une structure partitive, le verbe s'accorde selon le noyau du sujet : le partitif (le contenant/la portion) ou la matière (l'indénombrable) si le partitif est absent.

\begin{propriete}[Accord du verbe avec les partitifs]
\begin{itemize}
\item Si le sujet est la \textbf{matière seule} (nom indénombrable sans partitif), le verbe est au \textbf{singulier} :\\
\eng{The rice is good.} « Le riz est bon. »\\
\eng{This information is useful.} « Cette information est utile. »\\
\eng{Palm oil costs a lot.} « L'huile de palme coûte cher. »
\item Si le sujet est une \textbf{structure partitive} (quantité + partitif + \eng{of} + indénombrable), le verbe s'accorde avec le \textbf{partitif} (le noyau du syntagme nominal) :\\
\eng{A bottle of oil \textbf{is} on the table.} (singulier : \eng{a bottle})\\
\eng{Two bottles of oil \textbf{are} on the table.} (pluriel : \eng{two bottles})\\
\eng{A lot of sugar \textbf{is} needed.} (singulier : \eng{a lot} traité comme quantité unique) — mais \eng{A lot of students \textbf{are} here.} (pluriel : \eng{students} dénombrable).
\end{itemize}
Autrement dit : dans \eng{[Quantity + Partitive] of [Uncountable]}, le verbe s'accorde avec \textbf{Partitive}.
\end{propriete}

\begin{center}
\begin{tabularx}{\linewidth}{|X|X|X|}
\hline
\rowcolor{popblueL} \textbf{Phrase anglaise} & \textbf{Traduction} & \textbf{Justification de l'accord} \\
\hline
\eng{The water is cold.} & L'eau est froide. & Sujet = \eng{water} (indénombrable) $\rightarrow$ singulier \\
\hline
\eng{A glass of water is enough.} & Un verre d'eau suffit. & Sujet = \eng{a glass} (singulier) $\rightarrow$ \eng{is} \\
\hline
\eng{Two glasses of water are on the tray.} & Deux verres d'eau sont sur le plateau. & Sujet = \eng{two glasses} (pluriel) $\rightarrow$ \eng{are} \\
\hline
\eng{Three pieces of chalk were in the box.} & Trois morceaux de craie étaient dans la boîte. & Sujet = \eng{three pieces} (pluriel) $\rightarrow$ \eng{were} \\
\hline
\eng{A lot of money was spent.} & Beaucoup d'argent a été dépensé. & \eng{A lot} (quantité) $\rightarrow$ singulier \\
\hline
\end{tabularx}
\end{center}

\courssub{Cas particulier : \eng{two teas}, \eng{two coffees} au restaurant}

Il existe une exception d'usage, bien connue dans les pays anglophones : dans les cafés, restaurants et \eng{buvettes}, on entend souvent \eng{Two teas, please!} ou \eng{Three coffees and one hot chocolate.} Ici, \eng{tea} et \eng{coffee} sont employés comme des raccourcis commerciaux (ellipse) pour \eng{two cups of tea} et \eng{three cups of coffee}.

\begin{attention}[Usage commercial à signaler, pas à imiter à l'écrit scolaire]
\eng{Two teas, please.} = « Deux thés, s'il vous plaît. » (= \eng{two cups of tea}). Cet usage est \textbf{accepté à l'oral}, dans un contexte de commande (\eng{at the café}, \eng{at the buvette}). Mais dans un devoir écrit, une rédaction, un examen ou une situation formelle, écrivez \textbf{toujours} la forme complète : \eng{two cups of tea}. Et cet usage ne s'étend \textbf{jamais} à d'autres indénombrables : *\eng{two rices}, *\eng{three furnitures}, *\eng{four advices} restent \textbf{faux} partout.
\end{attention}

\courssub{Les partitifs avec des noms dénombrables au pluriel (groupement)}

La structure partitive ne sert pas qu'avec les indénombrables. Certains partitifs de \textbf{groupement} (\eng{collective partitives}) s'emploient avec des noms dénombrables au pluriel : le partitif est au singulier, mais le nom qui suit \eng{of} est au pluriel (car il s'agit d'un ensemble d'unités).

\begin{exemplebox}[Partitifs de groupement + nom dénombrable au pluriel]
\begin{itemize}
\item \eng{a bunch of flowers} = « un bouquet de fleurs » (\eng{flowers} pluriel)
\item \eng{a bunch of keys} = « un trousseau de clés » (\eng{keys} pluriel)
\item \eng{a bunch of bananas / plantains} = « une régime de bananes / plantains »
\item \eng{a group of students} = « un groupe d'élèves » (\eng{students} pluriel)
\item \eng{a pair of shoes} = « une paire de chaussures » (\eng{shoes} pluriel, \eng{pair} singulier)
\item \eng{a pair of trousers / jeans / scissors} = « un pantalon / un jean / des ciseaux » (noms \eng{pluralia tantum})
\item \eng{a crowd of people} = « une foule de gens » (\eng{people} pluriel)
\item \eng{a team of players} = « une équipe de joueurs »
\item \eng{a set of tools} = « un jeu d'outils »
\end{itemize}
Au pluriel, c'est encore le partitif qui change : \eng{two bunches of flowers} « deux bouquets de fleurs », \eng{three pairs of shoes} « trois paires de chaussures », \eng{several groups of students}.
\end{exemplebox}

\begin{center}
\begin{tabularx}{\linewidth}{|X|X|X|}
\hline
\rowcolor{popblueL} \textbf{Français} & \textbf{English} & \textbf{Analyse grammaticale} \\
\hline
une tasse de thé & \eng{a cup of tea} & \eng{tea} indénombrable, invariable \\
\hline
deux tasses de thé & \eng{two cups of tea} & \eng{-s} sur \eng{cup} seulement \\
\hline
un morceau de craie & \eng{a piece of chalk} & pas de *\eng{a chalk} \\
\hline
un bouquet de fleurs & \eng{a bunch of flowers} & \eng{flowers} dénombrable au pluriel \\
\hline
une paire de chaussures & \eng{a pair of shoes} & \eng{shoes} pluriel ; \eng{pair} singulier \\
\hline
un verre d'eau & \eng{a glass of water} & \eng{of} = « de / d' » \\
\hline
trois pains & \eng{three loaves of bread} & \eng{loaf} $\rightarrow$ \eng{loaves} (irrégulier) \\
\hline
\end{tabularx}
\end{center}

\courssub{Méthode : comment choisir la bonne structure}

\begin{methode}[Choisir entre \eng{a/an} + nom et partitif + \eng{of} + nom]
\begin{enumerate}
\item \textbf{Identifiez le nom :} Est-il \textbf{dénombrable} ? Peut-on dire \eng{one [nom], two [nom]s} ? (ex. \eng{book}, \eng{student}, \eng{flower}, \eng{banana}).
\item \textbf{Si OUI (dénombrable) :} Utilisez directement \eng{a/an} + nom (singulier) ou nombre + nom \textbf{-s} (pluriel). \eng{a flower}, \eng{three students}, \eng{a bunch of flowers} (si notion de groupe).
\item \textbf{Si NON (indénombrable : matière, liquide, abstraction) :} Vous \textbf{devez} utiliser un partitif : \textbf{contenant / portion / mesure} + \eng{of} + nom. Choisissez le partitif adapté au contexte : \eng{a cup of tea}, \eng{a slice of bread}, \eng{a piece of advice}, \eng{a bag of rice}, \eng{a bottle of oil}, \eng{a sheet of paper}.
\item \textbf{Mettez le pluriel au bon endroit :} Pour plusieurs unités, ajoutez le \eng{-s} (ou \eng{-es} / forme irrégulière) au \textbf{partitif uniquement} : \eng{two cups of tea}, \eng{three loaves of bread}. Jamais *\eng{two teas}.
\item \textbf{Accordez le verbe :} Le verbe s'accorde avec le \textbf{partitif} (sujet grammatical) : \eng{Two bottles of oil are…} ; \eng{The oil is…}.
\end{enumerate}
\end{methode}

\begin{exoresolu}[Application : construire le partitif correct]
Énoncé. Complétez chaque phrase avec le partitif correctement formé, à partir des éléments donnés entre parenthèses.
\begin{enumerate}
\item The teacher asked for \ldots\ (three / piece / chalk).
\item My mother bought \ldots\ (two / bag / rice) at the market in Mokolo.
\item \ldots\ (a / bunch / flower) is a nice gift for a visitor.
\item \ldots\ (two / bottle / palm oil) \ldots\ (be) on the kitchen table.
\item She gave me \ldots\ (a / piece / advice) that changed my life.
\item We need \ldots\ (five / sheet / paper) for the printer.
\end{enumerate}
\tcblower
Solution détaillée.
\begin{enumerate}
\item \eng{three pieces of chalk} — \eng{chalk} est indénombrable ; le pluriel porte sur \eng{piece} $\rightarrow$ \eng{pieces} ; \eng{chalk} reste invariable. « Le professeur a demandé trois morceaux de craie. »
\item \eng{two bags of rice} — *\eng{two rices} est impossible ; on compte les sacs : \eng{bags}. « Ma mère a acheté deux sacs de riz au marché de Mokolo. »
\item \eng{A bunch of flowers} — \eng{flowers} est dénombrable au pluriel après le partitif de groupement ; le verbe \eng{is} s'accorde avec \eng{bunch} (singulier). « Un bouquet de fleurs est un beau cadeau pour un visiteur. »
\item \eng{Two bottles of palm oil are} — le verbe s'accorde avec \eng{bottles} (pluriel) $\rightarrow$ \eng{are}, et non avec \eng{oil}. « Deux bouteilles d'huile de palme sont sur la table de la cuisine. »
\item \eng{a piece of advice} — \eng{advice} indénombrable ; pas de *\eng{an advice}. « Elle m'a donné un conseil qui a changé ma vie. »
\item \eng{five sheets of paper} — \eng{paper} indénombrable (matière) ; \eng{sheet} prend le \eng{-s}. « Nous avons besoin de cinq feuilles de papier pour l'imprimante. »
\end{enumerate}
\end{exoresolu}

\begin{exercice}[Entraînement : transformation et choix]
\begin{enumerate}
\item \textbf{Passage au pluriel.} Transformez les phrases suivantes au pluriel (quantité indiquée entre parenthèses). Faites les accords nécessaires.
\begin{itemize}
\item[a)] \eng{A cup of coffee is on the table.} (three)
\item[b)] \eng{There is a loaf of bread in the kitchen.} (two)
\item[c)] \eng{She needs a piece of furniture for her room.} (several)
\item[d)] \eng{A bottle of water costs 500 francs.} (four)
\item[e)] \eng{This is a useful piece of information.} (many)
\end{itemize}
\item \textbf{Choix du partitif.} Complétez avec le partitif approprié parmi la liste : \eng{bar / bottle / bunch / cup / loaf / piece / sheet / slice}.
\begin{itemize}
\item[a)] Can I have a \ldots\ of chocolate, please?
\item[b)] He drank two \ldots\ of tea at the buvette.
\item[c)] She cut a \ldots\ of bread for breakfast.
\item[d)] We bought a \ldots\ of palm oil at the supermarket.
\item[e)] The secretary typed three \ldots\ of paper.
\item[f)] A \ldots\ of plantains costs 1000 francs in Buea.
\item[g)] Please give me a \ldots\ of soap to wash my hands.
\item[h)] I need a \ldots\ of advice before I decide.
\end{itemize}
\item \textbf{Traduction (Français $\rightarrow$ Anglais).}
\begin{itemize}
\item[a)] Trois morceaux de craie.
\item[b)] Deux sacs de riz.
\item[c)] Une tranche de pain.
\item[d)] Plusieurs conseils.
\item[e)] Une bouteille d'huile de palme.
\item[f)] Cinq feuilles de papier.
\item[g)] Un bouquet de fleurs.
\item[h)] Deux verres d'eau.
\end{itemize}
\item \textbf{Correction d'erreurs.} Corrigez les phrases suivantes (toutes contiennent une erreur de structure partitive ou d'accord).
\begin{itemize}
\item[a)] *\eng{She bought two rices at the market.}
\item[b)] *\eng{Three furnitures are in the room.}
\item[c)] *\eng{He gave me an advice.}
\item[d)] *\eng{Two bottle of oil are on the shelf.}
\item[e)] *\eng{A bunch of flower is nice.}
\item[f)] *\eng{There are many informations in this book.}
\item[g)] *\eng{Two breads, please.} (contexte : boulangerie)
\item[h)] *\eng{Five piece of chalk are missing.}
\end{itemize}
\end{enumerate}
\end{exercice}

\begin{corrige}[Exercice n°1 -- Passage au pluriel]
\begin{itemize}
\item[a)] \eng{Three cups of coffee are on the table.}
\item[b)] \eng{There are two loaves of bread in the kitchen.}
\item[c)] \eng{She needs several pieces of furniture for her room.}
\item[d)] \eng{Four bottles of water cost 500 francs.} (ou \eng{cost 500 francs each} selon sens, mais \eng{four bottles} sujet pluriel $\rightarrow$ \eng{cost})
\item[e)] \eng{These are many useful pieces of information.} (ou \eng{much useful information} si quantité, mais consigne = pluriel partitif).
\end{itemize}
\end{corrige}

\begin{corrige}[Exercice n°2 -- Choix du partitif]
\begin{itemize}
\item[a)] \eng{bar} (\eng{a bar of chocolate})
\item[b)] \eng{cups} (\eng{two cups of tea})
\item[c)] \eng{slice} / \eng{loaf} (\eng{a slice of bread} = tranche ; \eng{a loaf of bread} = pain entier)
\item[d)] \eng{bottle} (\eng{a bottle of palm oil})
\item[e)] \eng{sheets} (\eng{three sheets of paper})
\item[f)] \eng{bunch} (\eng{A bunch of plantains})
\item[g)] \eng{bar} (\eng{a bar of soap})
\item[h)] \eng{piece} (\eng{a piece of advice})
\end{itemize}
\end{corrige}

\begin{corrige}[Exercice n°3 -- Traduction]
\begin{itemize}
\item[a)] \eng{Three pieces of chalk}
\item[b)] \eng{Two bags of rice}
\item[c)] \eng{A slice of bread} (ou \eng{a loaf of bread} selon contexte)
\item[d)] \eng{Several pieces of advice} (ou \eng{much advice})
\item[e)] \eng{A bottle of palm oil}
\item[f)] \eng{Five sheets of paper}
\item[g)] \eng{A bunch of flowers}
\item[h)] \eng{Two glasses of water}
\end{itemize}
\end{corrige}

\begin{corrige}[Exercice n°4 -- Correction d'erreurs]
\begin{itemize}
\item[a)] \eng{She bought two bags of rice at the market.} (\eng{rice} indénombrable)
\item[b)] \eng{Three pieces of furniture are in the room.} (\eng{furniture} indénombrable)
\item[c)] \eng{He gave me a piece of advice.} (\eng{advice} indénombrable)
\item[d)] \eng{Two bottles of oil are on the shelf.} (pluriel sur \eng{bottle} + accord verbe \eng{are})
\item[e)] \eng{A bunch of flowers is nice.} (\eng{flowers} pluriel après \eng{bunch of} ; verbe \eng{is} accordé à \eng{bunch})
\item[f)] \eng{There is much information in this book.} (ou \eng{many pieces of information})
\item[g)] \eng{Two loaves of bread, please.} (contexte boulangerie : forme complète ; *\eng{two breads} faux hors commande café pour \eng{tea/coffee})
\item[h)] \eng{Five pieces of chalk are missing.} (pluriel sur \eng{piece} + accord \eng{are})
\end{itemize}
\end{corrige}

\begin{aretenir}[L'essentiel de la règle : Expressions partitives]
\begin{itemize}
\item \textbf{Structure de base :} \textbf{Quantité + Partitif + \eng{of} + Nom indénombrable}.
\item \textbf{Pluriel :} Seul le \textbf{partitif} prend le \eng{-s} (\eng{-es} / irrégulier). Le nom indénombrable reste \textbf{invariable}.
\eng{two cups of tea} \checkmark \quad *\eng{two teas} \cross (sauf café/restaurant).
\item \textbf{Accord du verbe :} Le verbe s'accorde avec le \textbf{partitif} (sujet grammatical).
\eng{Two bottles of oil \textbf{are}…} \quad \eng{The oil \textbf{is}…}
\item \textbf{Partitifs de groupement (noms dénombrables) :} \eng{a bunch of flowers}, \eng{a pair of shoes}, \eng{a group of students}. Le nom après \eng{of} est au \textbf{pluriel}. Le verbe s'accorde avec le partitif (\eng{a bunch \textbf{is}}, \eng{two bunches \textbf{are}}).
\item \textbf{Principaux partitifs à connaître :} \eng{a cup of}, \eng{a glass of}, \eng{a bottle of}, \eng{a piece of}, \eng{a slice of}, \eng{a loaf of}, \eng{a bar of}, \eng{a sheet of}, \eng{a bag of}, \eng{a bunch of}, \eng{a pair of}, \eng{a group of}, \eng{a box of}, \eng{a jar of}, \eng{a tin/can of}, \eng{a kilo/pound of}.
\end{itemize}
\end{aretenir}

\begin{savaistu}[Le saviez-vous ? : Mesures et culture]
\begin{itemize}
\item Au Royaume-Uni et dans les pays du Commonwealth (dont le Cameroun), on utilise souvent le système métrique (\eng{kilogramme}, \eng{litre}, \eng{metre}) mais les vieilles unités persistent : \eng{a pint of milk} (environ 568 ml), \eng{a pound of sugar} (environ 454 g). En cuisine nigériane ou ghanéenne, on achète souvent \eng{a mudu of rice} (mesure locale) ou \eng{a paint bucket of garri}.
\item L'expression \eng{a piece of cake} signifie littéralement « un morceau de gâteau », mais idiomatiquement « c'est très facile » (\eng{It's a piece of cake!} = « C'est du gâteau ! » / « C'est facile ! »).
\item \eng{A bunch of} s'emploie aussi familièrement pour « un groupe de gens » : \eng{They are a nice bunch} = « Ce sont des gens sympas / un chouette groupe ».
\end{itemize}
\end{savaistu}

\coursec{Les principaux partitifs par catégories}

Dans la section précédente, nous avons établi la règle générale : pour exprimer une quantité, l'anglais utilise la structure \cle{a + nom de contenant/mesure + of + nom}. Mais quels contenants ? Quelles mesures ? C'est précisément l'objet de cette section : apprendre à choisir le \emph{bon} partitif selon la nature du nom qu'il introduit. Nous allons classer les partitifs en grandes familles : les contenants, les portions, les formes et mesures, les groupes, et enfin les partitifs abstraits.

\courssub{La règle : structure et formation}

\begin{propriete}[Structure de base des partitifs]
Pour quantifier un nom indénombrable (substance, matière, notion abstraite) ou un nom dénombrable vendu en vrac, l'anglais utilise obligatoirement un \cle{partitif} :\\
\eng{a/an + nom de quantité/contenant + of + nom}\\
Exemples : \eng{a cup of tea}, \eng{a piece of advice}, \eng{a bunch of plantains}.
\end{propriete}

\begin{attention}[Of ne se traduit pas toujours par « de »]
En français, on dit « une tasse \emph{de} thé », « un morceau \emph{de} pain ». En anglais, \eng{of} est \emph{obligatoire} après le partitif, même si le français omet la préposition (ex. \eng{a box matches} $\rightarrow$ incorrect). Ne dites jamais \eng{a cup tea} ni \eng{a piece advice}.
\end{attention}

\courssub{Observation : une liste de courses}

Commençons par observer un document authentique de la vie quotidienne : la liste de courses qu'une mère de famille de Buea a écrite avant d'aller au marché.

\begin{textebox}[Shopping list for the week]
Shopping list for the week — Mrs Njie, Buea

1 bag of rice (50 kg, for the whole family)

2 loaves of bread from the bakery on Molyko Street

1 bunch of plantains — not too ripe, please!

3 tins of sardines for the children's sandwiches

1 packet of sugar and 1 packet of tea

2 bars of soap for the laundry

1 bottle of palm oil from the market

1 box of matches and 1 tin of tomatoes

1 piece of cloth for my daughter's school uniform

A few slices of meat for Sunday's lunch

Don't forget: a pair of scissors for the tailor!
\end{textebox}

\textbf{Glossaire :} \eng{loaf} (n., miche de pain, pain entier) ; \eng{bunch} (n., régime, bouquet) ; \eng{ripe} (adj., mûr) ; \eng{tin} (n., boîte de conserve) ; \eng{bar} (n., savonnette, barre) ; \eng{laundry} (n., lessive) ; \eng{cloth} (n., tissu, pagne) ; \eng{tailor} (n., tailleur).

\textbf{Questions de compréhension :}
\begin{enumerate}
\item How many loaves of bread does Mrs Njie want to buy?
\item What does she need for the children's sandwiches?
\item Which partitive is used with “plantains”? And with “soap”?
\end{enumerate}

Observez : chaque nom est précédé d'un partitif différent. On ne dit pas \eng{a rice} ni \eng{a soap} : ces noms sont indénombrables ou vendus en vrac, et le partitif est \emph{obligatoire} pour les compter. Voyons maintenant ces familles une par une.

\courssub{Les contenants (containers)}

\begin{definition}[Contenants]
Un \cle{contenant} est un objet qui \emph{contient} une substance : bouteille, tasse, verre, boîte, sac, paquet. En anglais, on nomme d'abord le contenant, puis \eng{of}, puis le contenu : \eng{a bottle of water} = « une bouteille d'eau ».
\end{definition}

\begin{center}
\begin{tabularx}{\linewidth}{|X|X|X|}\hline
\rowcolor{popblueL} Partitif anglais & Français & Exemples typiques \\ \hline
a bottle of & une bouteille de & water, oil, wine, palm oil \\ \hline
a cup of & une tasse de & tea, coffee, garri \\ \hline
a glass of & un verre de & water, juice, beer \\ \hline
a tin / a can of & une boîte de & tomatoes, sardines, milk \\ \hline
a bag of & un sac de & rice, flour, cement \\ \hline
a packet of & un paquet de & sugar, biscuits, tea \\ \hline
a box of & une boîte de & matches, chalk, chocolates \\ \hline
\end{tabularx}
\end{center}

\begin{exemplebox}[Contenants en contexte]
\eng{She bought a bottle of palm oil at Mokolo market.} « Elle a acheté une bouteille d'huile de palme au marché Mokolo. »

\eng{Could I have a cup of tea, please?} « Puis-je avoir une tasse de thé, s'il vous plaît ? »

\eng{The workers carried two bags of cement to the lorry.} « Les ouvriers ont porté deux sacs de ciment jusqu'au camion. »

\eng{We opened a tin of sardines for lunch.} « Nous avons ouvert une boîte de sardines pour le déjeuner. »
\end{exemplebox}

\begin{attention}[Tin ou can ?]
\eng{Tin} est britannique, \eng{can} est américain : \eng{a tin of tomatoes} (GB) = \eng{a can of tomatoes} (US). Au Cameroun anglophone, on entend les deux. Attention aussi : \eng{a box of} se traduit par « une boîte de » (allumettes, craies), tandis que \eng{a tin of} désigne une boîte \emph{métallique} de conserve. Ne confondez pas avec le français « une boîte de sardines » qui correspond à \eng{a tin of sardines}.
\end{attention}

\courssub{Les portions et les morceaux}

C'est la famille la plus importante, car elle permet de \emph{compter les noms indénombrables}.

\begin{definition}[A piece of]
\cle{A piece of} (un morceau de, un bout de) est le partitif le plus polyvalent : il s'emploie avec les noms concrets (\eng{bread, cake, meat, chalk, paper}) mais aussi avec de nombreux noms \emph{abstraits indénombrables} : \eng{a piece of advice} (un conseil), \eng{a piece of information} (une information), \eng{a piece of news} (une nouvelle), \eng{a piece of furniture} (un meuble), \eng{a piece of luggage} (un bagage).
\end{definition}

\begin{propriete}[Les portions principales]
\begin{itemize}
\item \cle{a slice of} = une tranche de (pain, viande, gâteau) : \eng{a slice of bread}.
\item \cle{a loaf of bread} = un pain (entier) ; pluriel irrégulier : \eng{two loaves of bread}.
\item \cle{a lump of} = un morceau compact de (sucre, terre) : \eng{a lump of sugar} = « un morceau de sucre ».
\end{itemize}
\end{propriete}

\begin{exemplebox}[Exemples traduits]
\eng{a slice of bread} = « une tranche de pain »

\eng{Would you like a piece of cake?} « Voulez-vous un morceau de gâteau ? »

\eng{The teacher gave me a piece of chalk.} « Le professeur m'a donné un bout de craie. »

\eng{Let me give you a piece of advice.} « Laisse-moi te donner un conseil. »

\eng{She put two lumps of sugar in her tea.} « Elle a mis deux morceaux de sucre dans son thé. »
\end{exemplebox}

\begin{attention}[Le pluriel irrégulier de loaf]
\eng{A loaf of bread} a un pluriel irrégulier : \cle{two loaves of bread} (prononcé « louvze »). Comme \eng{wife} $\rightarrow$ \eng{wives} et \eng{leaf} $\rightarrow$ \eng{leaves}, le \emph{-f} final devient \emph{-ves}. Erreur fréquente : \eng{two loafs} — à bannir absolument !
\end{attention}

\courssub{Les formes et les mesures (shapes and measures)}

Ces partitifs décrivent la \emph{forme} ou la \emph{quantité minimale} d'une substance.

\begin{center}
\begin{tabularx}{\linewidth}{|X|X|X|}\hline
\rowcolor{popblueL} Partitif anglais & Français & Exemples typiques \\ \hline
a bar of & une savonnette / une barre de & soap, chocolate, gold \\ \hline
a sheet of & une feuille de & paper \\ \hline
a drop of & une goutte de & water, rain, oil \\ \hline
a grain of & un grain de & rice, sand \\ \hline
a spoonful of & une cuillerée de & sugar, salt \\ \hline
\end{tabularx}
\end{center}

\begin{exemplebox}[Exemples traduits]
\eng{a bar of soap} = « une savonnette » ; \eng{a bar of chocolate} = « une tablette de chocolat »

\eng{a drop of rain} = « une goutte de pluie » ; \eng{There isn't a drop of oil left.} « Il ne reste pas une goutte d'huile. »

\eng{a grain of rice} = « un grain de riz » ; \eng{Not a grain of rice was wasted.} « Pas un grain de riz n'a été gaspillé. »

\eng{Add a spoonful of salt to the sauce.} « Ajoute une cuillerée de sel à la sauce. »

\eng{Please give me a sheet of paper.} « Donne-moi une feuille de papier, s'il te plaît. »
\end{exemplebox}

\begin{propriete}[Formation des mesures en -ful]
Les mesures en \cle{-ful} (poignée, cuillerée, bouche pleine) forment leur pluriel sur le nom de base : \eng{a spoonful of sugar} $\rightarrow$ \eng{two spoonfuls of sugar} ; \eng{a handful of groundnuts} $\rightarrow$ \eng{two handfuls of groundnuts}. On écrit \eng{spoonfuls}, jamais \eng{spoonsful} à l'examen.
\end{propriete}

\courssub{Les groupes et les collections}

\begin{definition}[Partitifs de groupe]
Ces partitifs rassemblent des éléments \emph{semblables} en un ensemble : \cle{a bunch of} (régime, bouquet, trousseau), \cle{a pair of} (paire), \cle{a crowd of} (foule), \cle{a team of} (équipe), \cle{a flock of} (troupeau de moutons, volée d'oiseaux), \cle{a herd of} (troupeau de bovins).
\end{definition}

\begin{exemplebox}[Exemples traduits]
\eng{a bunch of plantains} = « un régime de plantains » ; \eng{a bunch of keys} = « un trousseau de clés » ; \eng{a bunch of flowers} = « un bouquet de fleurs » ; \eng{a bunch of grapes} = « une grappe de raisin »

\eng{a pair of scissors} = « une paire de ciseaux » ; \eng{a pair of shoes} = « une paire de chaussures » ; \eng{a pair of trousers} = « un pantalon » ; \eng{a pair of glasses} = « une paire de lunettes »

\eng{A crowd of people gathered at the stadium.} « Une foule de personnes s'est rassemblée au stade. »

\eng{A flock of sheep crossed the road near Bamenda.} « Un troupeau de moutons a traversé la route près de Bamenda. »

\eng{The Fulani shepherd led a herd of cattle to the river.} « Le berger peul a conduit un troupeau de bœufs à la rivière. »
\end{exemplebox}

\begin{attention}[A pair of : singulier ou pluriel ?]
\eng{Trousers, scissors, glasses} sont toujours au \emph{pluriel} en anglais, mais \eng{a pair of} est \emph{singulier} : \eng{This pair of scissors is sharp.} « Cette paire de ciseaux est bien aiguisée. » Au pluriel : \eng{two pairs of shoes}. Ne dites jamais \eng{a scissors} ni \eng{a trousers} !
\end{attention}

\courssub{Les partitifs abstraits}

Certains noms abstraits sont indénombrables en anglais alors que leur équivalent français est dénombrable. C'est l'un des pièges majeurs du Bac.

\begin{propriete}[Noms abstraits indénombrables + partitif]
\eng{advice, information, news, work, furniture, luggage} sont indénombrables. Pour les compter : \cle{a piece of advice} (un conseil), \cle{a piece of information} (une information), \cle{a piece of news} (une nouvelle), \cle{a piece of work} (un travail), \cle{an item of news/clothing/furniture} (une nouvelle, un vêtement, un meuble), \cle{a bit of luck/time} (un peu de chance, un peu de temps).
\end{propriete}

\begin{exemplebox}[Exemples traduits]
\eng{That's an interesting piece of news!} « C'est une nouvelle intéressante ! »

\eng{I need a bit of time to finish my homework.} « J'ai besoin d'un peu de temps pour finir mes devoirs. »

\eng{With a bit of luck, we shall pass the exam.} « Avec un peu de chance, nous réussirons l'examen. »

\eng{Each item of furniture was carried upstairs.} « Chaque meuble a été monté à l'étage. »
\end{exemplebox}

\begin{attention}[Les faux amis indénombrables]
Erreurs classiques des francophones : \eng{an advice}, \eng{an information}, \eng{a news}, \eng{furnitures}, \eng{luggages} sont \textbf{impossibles}. Dites : \eng{a piece of advice}, \eng{a piece of information}, \eng{a piece of news}. Rappelez-vous : en français « un conseil, une information » se comptent ; en anglais, jamais sans partitif.
\end{attention}

\courssub{Figure récapitulative : les quatre familles}

\begin{popfigure}
\begin{center}
\begin{tikzpicture}[
fam/.style={rectangle, rounded corners, draw=#1, fill=#1!15, text width=3.1cm, align=center, font=\small, inner sep=4pt},
titre/.style={rectangle, rounded corners, draw=#1, fill=#1, text=white, text width=3.1cm, align=center, font=\small\bfseries, inner sep=3pt}]
\node[titre=popblue] (t1) at (0,0) {Containers};
\node[fam=popblue, below=0.15cm of t1, align=center]{a bottle of oil\\ a cup of tea\\ a glass of juice\\ a tin of sardines\\ a bag of rice};
\node[titre=poporange] (t2) at (3.9,0) {Portions};
\node[fam=poporange, below=0.15cm of t2, align=center]{a piece of cake\\ a slice of bread\\ a loaf of bread\\ a lump of sugar\\ a piece of chalk};
\node[titre=popgreen] (t3) at (7.8,0) {Shapes \& measures};
\node[fam=popgreen, below=0.15cm of t3, align=center]{a bar of soap\\ a sheet of paper\\ a drop of rain\\ a grain of rice\\ a spoonful of salt};
\node[titre=poppurple] (t4) at (11.7,0) {Groups};
\node[fam=poppurple, below=0.15cm of t4, align=center]{a bunch of keys\\ a pair of shoes\\ a crowd of people\\ a flock of sheep\\ a herd of cattle};
\end{tikzpicture}
\end{center}
\legende{Les quatre grandes familles de partitifs, avec cinq exemples chacune.}
\end{popfigure}

\courssub{Tableau récapitulatif : les 25 partitifs les plus fréquents au Bac}

\begin{center}
\begin{tabularx}{\linewidth}{|X|X|}\hline
\rowcolor{popblueL} English & Français \\ \hline
a bottle of palm oil & une bouteille d'huile de palme \\ \hline
a cup of tea & une tasse de thé \\ \hline
a glass of water & un verre d'eau \\ \hline
a tin of sardines & une boîte de sardines \\ \hline
a bag of rice & un sac de riz \\ \hline
a packet of sugar & un paquet de sucre \\ \hline
a box of matches & une boîte d'allumettes \\ \hline
a piece of bread & un morceau de pain \\ \hline
a piece of advice & un conseil \\ \hline
a piece of information & une information \\ \hline
a piece of news & une nouvelle \\ \hline
a slice of meat & une tranche de viande \\ \hline
a loaf of bread (loaves) & un pain \\ \hline
a lump of sugar & un morceau de sucre \\ \hline
a bar of soap & une savonnette \\ \hline
a bar of chocolate & une tablette de chocolat \\ \hline
a sheet of paper & une feuille de papier \\ \hline
a drop of rain & une goutte de pluie \\ \hline
a grain of rice & un grain de riz \\ \hline
a spoonful of salt & une cuillerée de sel \\ \hline
a bunch of plantains & un régime de plantains \\ \hline
a pair of scissors & une paire de ciseaux \\ \hline
a crowd of people & une foule de personnes \\ \hline
a flock of sheep & un troupeau de moutons \\ \hline
a herd of cattle & un troupeau de bœufs \\ \hline
\end{tabularx}
\end{center}

\courssub{Exercice résolu et stratégie de mémorisation}

\begin{exoresolu}[Complétez avec le partitif correct]
Énoncé — Complete with the correct partitive: \emph{bar, bunch, drop, grain, loaf, pair, piece, sheet, spoonful, tin}.

1. She bought a \ldots\ of soap and a \ldots\ of sardines. 2. There was not a \ldots\ of water in the house. 3. He gave me a \ldots\ of paper and a \ldots\ of good advice. 4. The farmer sold a \ldots\ of plantains at the market. 5. Add a \ldots\ of salt, please. 6. I need a new \ldots\ of shoes. 7. Not a single \ldots\ of rice was left in the bag.
\tcblower
Solution détaillée :

1. \eng{a bar of soap} (savonnette) ; \eng{a tin of sardines} (conserve métallique). 2. \eng{a drop of water} (quantité minimale de liquide). 3. \eng{a sheet of paper} (feuille) ; \eng{a piece of advice} (\emph{advice} est indénombrable). 4. \eng{a bunch of plantains} (régime). 5. \eng{a spoonful of salt} (cuillerée). 6. \eng{a pair of shoes} (paire). 7. \eng{a grain of rice} (grain).
\end{exoresolu}

\begin{methode}[Comment mémoriser les partitifs]
\begin{enumerate}
\item N'apprenez jamais un partitif avec sa traduction mot à mot : associez-le à une \emph{image mentale} — une bouteille pour \eng{bottle}, une tranche pour \eng{slice}, une grappe ou un régime pour \eng{bunch}.
\item Classez les partitifs en quatre familles (contenants, portions, mesures, groupes) et apprenez cinq exemples par famille.
\item Mémorisez à part les cas pièges : \eng{a piece of advice/information/news} et le pluriel \eng{loaves}.
\item Réutilisez-les à l'oral : décrivez chaque jour votre liste de courses ou votre repas en anglais (\eng{a cup of coffee}, \eng{a slice of bread}, \eng{a piece of meat}).
\end{enumerate}
\end{methode}

\begin{savaistu}[Les partitifs s'adaptent aux réalités locales]
Au Nigeria et dans les régions anglophones du Cameroun, on entend souvent au marché \eng{a cup of garri} ou \eng{a basin of garri} : le gari se vend à la tasse ou à la cuvette ! De même, on dira \eng{a basin of tomatoes} ou \eng{a bucket of groundnuts}. Les partitifs anglais ne sont pas figés : ils épousent les contenants réels de la vie quotidienne. À vous de créer les vôtres : \eng{a plate of ndolé}, \eng{a pot of eru}...
\end{savaistu}

\begin{exercice}[À vous de jouer]
1. Traduisez : « Ma mère a acheté deux pains, un régime de bananes et une bouteille d'huile. » 2. Mettez au pluriel : \eng{a loaf of bread, a box of matches, a piece of news, a pair of trousers}. 3. Trouvez le partitif qui convient : a \ldots\ of chocolate, a \ldots\ of sand, a \ldots\ of cattle, a \ldots\ of chalk.
\end{exercice}

\begin{corrige}[Exercice « À vous de jouer »]
1. \eng{My mother bought two loaves of bread, a bunch of bananas and a bottle of oil.} 2. \eng{loaves of bread, boxes of matches, pieces of news, pairs of trousers.} 3. \eng{a bar of chocolate, a grain of sand, a herd of cattle, a piece of chalk.}
\end{corrige}

\coursec{Comparaison français / anglais}

\courssub{Transition depuis la classification des partitifs}

Dans la section précédente, nous avons répertorié les principaux partitifs anglais par catégories (contenants, mesures, portions, formes). Nous avons vu que l'anglais utilise invariablement la structure \cle{a/an + nom de quantité + of + nom de substance}. Nous allons maintenant confronter cette mécanique à la logique du français, car c'est précisément à cette interface que se nichent les erreurs les plus tenaces pour un apprenant francophone. Le français et l'anglais ne découpent pas la réalité matérielle de la même manière : ce qui est « comptable » dans une langue ne l'est pas nécessairement dans l'autre.

\courssub{Deux logiques de quantification : article partitif vs nom de portion}

\textbf{Observation.} Comparez ces deux échanges au marché de Buea :

\begin{textebox}[Dialogue au marché : quantification en anglais et en français]
-- \eng{How much bread would you like?}
-- \eng{I'll take a loaf, please. And a bunch of bananas.}
-- « Combien de pain voulez-vous ? »
-- « Je prendrai un pain, s'il vous plaît. Et une grappe de bananes. »
\end{textebox}

\textbf{Analyse.}
\begin{itemize}
    \item En \textbf{français}, l'article partitif (\cle{du, de la, de l', des}) ou l'article indéfini (\cle{un, une}) suffit souvent à exprimer une quantité indéterminée ou une unité : \eng{Je bois \textbf{de l'eau}} / \eng{Je mange \textbf{du pain}} / \eng{Je veux \textbf{un pain}}.
    \item En \textbf{anglais}, les noms de substances (\cle{water, bread, rice, meat}) sont \cle{indénombrables (uncountable)}. On ne peut pas dire *\eng{a water}, *\eng{a bread}, *\eng{an advice}. Il est \textbf{obligatoire} d'insérer un \textbf{nom de portion ou de contenant} (partitive noun) suivi de \cle{of} : \eng{a \textbf{glass} of water}, \eng{a \textbf{loaf} of bread}, \eng{a \textbf{piece} of advice}.
\end{itemize}

\begin{propriete}[Règle de comparaison : quantification française vs anglaise]
\textbf{Français :} quantification par \textbf{article} (partitif \textit{du/de la} ou indéfini \textit{un/une}) + \textbf{nom nu}.\par
\textbf{Anglais :} quantification par \textbf{nom concret de portion ou de contenant} (\textit{a glass, a slice, a piece, a loaf, a bunch}) + \textbf{of} + \textbf{nom de substance}.\par
\medskip
\textbf{Conséquence majeure :} le partitif anglais est \textbf{toujours un nom concret} (un objet, une forme, une mesure), jamais un simple article. On ne peut pas individualiser une unité avec \textit{some} seul : \textit{some} exprime une quantité floue (\eng{some water} = « de l'eau »), alors que \eng{a glass of water} individualise une portion précise (« un verre d'eau »).
\end{propriete}

\courssub{Les « noms-pièges » : indénombrables en anglais, dénombrables en français}

C'est le point noir numéro un pour les francophones. Une série de noms très fréquents sont \textbf{comptables en français} (on peut dire \textit{un conseil, deux informations, trois meubles}) mais \textbf{strictement indénombrables en anglais}. Ils exigent systématiquement \cle{a piece of} (ou \cle{an item of}, \cle{a bit of} à l'oral familier).

\begin{center}
\begin{tabularx}{\linewidth}{|X|X|X|}
\hline
\rowcolor{popblueL}
\textbf{Français (dénombrable)} & \textbf{Anglais (indénombrable)} & \textbf{Solution partitive correcte} \\
\hline
un conseil & advice & \eng{a piece of advice} \\
deux conseils & advice & \eng{two pieces of advice} \\
\hline
une information & information & \eng{a piece of information} \\
trois informations & information & \eng{three pieces of information} \\
\hline
une nouvelle & news & \eng{a piece of news} \\
\hline
un meuble & furniture & \eng{a piece of furniture} \\
deux meubles & furniture & \eng{two pieces of furniture} \\
\hline
un bagage & luggage / baggage & \eng{a piece of luggage} \\
\hline
un devoir (scolaire) & homework & \eng{a piece of homework} \\
\hline
une somme d'argent & money & \eng{a sum of money} / \eng{an amount of money} \\
\hline
un pain / une tranche de pain & bread & \eng{a loaf of bread} / \eng{a slice of bread} \\
\hline
un travail / un emploi & work & \eng{a piece of work} / \eng{a job} \\
\hline
un équipement & equipment & \eng{a piece of equipment} \\
\hline
\end{tabularx}
\end{center}

\begin{exemplebox}[Exemples traduits : noms-pièges au quotidien]
\eng{The teacher gave us \textbf{a useful piece of advice} about the exam.} \\
« Le professeur nous a donné un conseil utile pour l'examen. »

\eng{She bought \textbf{two pieces of furniture} for her new flat in Yaoundé.} \\
« Elle a acheté deux meubles pour son nouvel appartement à Yaoundé. »

\eng{He carried \textbf{a heavy piece of luggage} up the stairs.} \\
« Il a porté un lourd bagage en haut de l'escalier. »

\eng{I have \textbf{a lot of homework} to do tonight.} \\
« J'ai beaucoup de devoirs à faire ce soir. »

\eng{\textbf{A piece of information} is missing from your form.} \\
« Une information manque sur votre formulaire. »
\end{exemplebox}

\begin{attention}[Erreurs fatales : *\eng{an advice}, *\eng{an information}, *\eng{a furniture}, *\eng{a luggage}, *\eng{a homework}, *\eng{a bread}, *\eng{a money}]
\textbf{Jamais} d'article indéfini \cle{a/an} directement devant ces noms. \textbf{Jamais} de \cle{-s} au pluriel : *\eng{advices}, *\eng{informations}, *\eng{furnitures}, *\eng{luggages}, *\eng{homeworks}, *\eng{breads}, *\eng{moneys} sont des formes inexistantes.

\medskip
\textbf{Stratégies de contournement (à l'oral comme à l'écrit) :}
\begin{enumerate}
    \item \textbf{Le partitif :} \eng{a piece of advice / information / furniture / luggage / homework / equipment} ; \eng{an item of luggage / furniture}.
    \item \textbf{Le verbe équivalent :} \eng{He \textbf{advised} me to leave} (au lieu de *\eng{He gave me an advice}) ; \eng{Please \textbf{inform} us} (au lieu de *\eng{Give us an information}).
    \item \textbf{Un nom comptable proche :} \eng{a \textbf{job}} (pour « un emploi », car \eng{work} est indénombrable) ; \eng{a \textbf{loaf} of bread}, \eng{a \textbf{slice} of bread}, \eng{a \textbf{sum} of money}.
\end{enumerate}
\end{attention}

\courssub{Le cas inverse : dénombrables en anglais, partitif ou collectif en français}

Quelques noms courants fonctionnent à l'envers : ils sont comptables en anglais mais traités comme une masse (partitif) en français, ou inversement.

\begin{propriete}[Cas inverses notables]
\begin{itemize}
    \item \textbf{Grapes (raisin) :} en anglais, \cle{grape} est un nom comptable (\eng{a grape, two grapes}). Une grappe se dit \eng{a bunch of grapes}. En français, on achète \textit{du raisin} (partitif) ou \textit{une grappe de raisin}.
    \item \textbf{Spaghetti / pasta / macaroni :} en anglais, ce sont des noms \textbf{indénombrables} (\eng{spaghetti is...}). On dit \eng{a plate of spaghetti}, \eng{a bowl of pasta}. En français, on utilise souvent le pluriel : \textit{des spaghettis, des pâtes}.
    \item \textbf{Fruit :} souvent indénombrable en anglais (\eng{a piece of fruit}, \eng{eat more fruit}), alors que le français dit couramment \textit{des fruits} au pluriel.
    \item \textbf{Hair / cheveu :} \eng{hair} est indénombrable quand il désigne la chevelure : \eng{She has long hair}. Un cheveu isolé se dit \eng{a hair} : \eng{There is a hair in my soup}. Le français oppose \textit{des cheveux} (collectif) à \textit{un cheveu}.
\end{itemize}
\end{propriete}

\begin{exemplebox}[Grapes et spaghetti : attention au piège]
\eng{I bought \textbf{a bunch of grapes} at the market.} (Correct : \textit{grapes} au pluriel après \textit{a bunch of}.) \\
« J'ai acheté une grappe de raisin au marché. »

\eng{Would you like \textbf{a plate of spaghetti}?} (Correct : \textit{spaghetti} invariable.) \\
« Voulez-vous une assiette de spaghettis ? »

\eng{She found \textbf{a hair} in her salad.} (Un seul cheveu.) \\
« Elle a trouvé un cheveu dans sa salade. »

\eng{He has \textbf{grey hair}.} (La chevelure = masse indénombrable.) \\
« Il a les cheveux gris. »
\end{exemplebox}

\courssub{Tableau de synthèse contrastif}

Le schéma ci-dessous résume visuellement les quatre grands profils de correspondance entre la quantification française et anglaise.

\begin{popfigure}
\begin{tikzpicture}[
    box/.style={rectangle, draw=popblue, thick, rounded corners, fill=popblueL, minimum width=3.6cm, minimum height=1.2cm, align=center, font=\small},
    labelbox/.style={rectangle, draw=poporange, thick, rounded corners, fill=poporangeL, minimum width=3.6cm, minimum height=0.9cm, align=center, font=\small\bfseries},
    arrow/.style={-Stealth, thick, popdark},
]
% Titres de colonnes
\node[labelbox] (fr) at (0,4) {Français};
\node[labelbox] (en) at (5.5,4) {Anglais (nature du nom)};
\node[labelbox] (sol) at (11,4) {Forme correcte};

% Ligne 1 : masses standard
\node[box, align=center] (fr1) at (0,2.5){du / de la / de l' + nom de masse\\ \textit{(de l'eau, du riz)}};
\node[box, align=center] (en1) at (5.5,2.5){nom indénombrable\\ \textit{(water, rice)}};
\node[box, align=center] (sol1) at (11,2.5){\eng{some water / rice}\\ ou \eng{a glass / a bowl of water / rice}};

% Ligne 2 : noms-pièges
\node[box, align=center] (fr2) at (0,1){un / une + nom comptable\\ \textit{(un conseil, un meuble)}};
\node[box, align=center] (en2) at (5.5,1){nom \textbf{indénombrable}\\ \textit{(advice, furniture)}};
\node[box, align=center] (sol2) at (11,1){\eng{\textbf{a piece of} advice / furniture}\\ jamais *\eng{an advice}};

% Ligne 3 : objets individualisés
\node[box, align=center] (fr3) at (0,-0.5){un / une + nom comptable\\ \textit{(un pain, une grappe)}};
\node[box, align=center] (en3) at (5.5,-0.5){indénombrable (\textit{bread})\\ ou comptable (\textit{grape})};
\node[box, align=center] (sol3) at (11,-0.5){\eng{a \textbf{loaf / slice} of bread}\\ \eng{a \textbf{bunch of} grapes}};

% Ligne 4 : cas inverses
\node[box, align=center] (fr4) at (0,-2){des + pluriel\\ \textit{(des spaghettis, des cheveux)}};
\node[box, align=center] (en4) at (5.5,-2){nom indénombrable\\ \textit{(spaghetti, hair)}};
\node[box, align=center] (sol4) at (11,-2){\eng{a plate / a bowl of spaghetti}\\ \eng{hair} (chevelure) / \eng{a hair} (un seul)};

% Flèches
\draw[arrow] (fr1) -- (en1);
\draw[arrow] (en1) -- (sol1);
\draw[arrow] (fr2) -- (en2) node[midway, above, font=\tiny] {piège};
\draw[arrow] (en2) -- (sol2) node[midway, above, font=\tiny] {obligatoire};
\draw[arrow] (fr3) -- (en3);
\draw[arrow] (en3) -- (sol3);
\draw[arrow] (fr4) -- (en4) node[midway, above, font=\tiny] {inversion};
\draw[arrow] (en4) -- (sol4);
\end{tikzpicture}
\legende{Synthèse visuelle : quatre profils de correspondance français / anglais pour la quantification.}
\end{popfigure}

\courssub{Focus culturel et linguistique : le cas de \eng{news}}

\begin{savaistu}[Pourquoi \eng{news} prend-il le singulier ?]
Le mot \eng{news} ressemble à un pluriel (il finit par \textit{-s}), mais c'est un \textbf{nom indénombrable singulier}. Il vient de l'anglais moyen \textit{newes} (« nouvelles choses »), figé au fil du temps en collectif singulier.
\begin{itemize}
    \item \eng{The news \textbf{is} good today.} — « Les nouvelles sont bonnes aujourd'hui. »
    \item \eng{This \textbf{is} a surprising piece of news.} — « C'est une nouvelle surprenante. »
\end{itemize}
\textbf{Autres mots en \textit{-s} qui prennent le singulier :} \eng{mathematics, physics, economics, politics, measles, billiards, darts}. Ils se construisent tous avec un verbe au singulier : \eng{Mathematics \textbf{is} difficult for many students.} (« Les mathématiques sont difficiles pour beaucoup d'élèves. »)
\end{savaistu}

\courssub{Exercice résolu : transformation et correction}

\begin{exoresolu}[Correction d'erreurs de quantification (francophonismes)]
\textbf{Énoncé :} corrigez les phrases suivantes. Elles contiennent toutes des erreurs typiques de transfert du français vers l'anglais (article \textit{a/an} devant un indénombrable, pluriel en \textit{-s} interdit, mauvais partitif).

\begin{enumerate}
    \item \eng{She gave me \textbf{an advice} for my interview.}
    \item \eng{We bought \textbf{two furnitures} for the living room.}
    \item \eng{He has \textbf{long hairs} and a beard.}
    \item \eng{I need \textbf{an information} about the train schedule.}
    \item \eng{She cooked \textbf{a spaghetti} for dinner.}
    \item \eng{They have \textbf{many luggages}.}
    \item \eng{Can I have \textbf{a bread} for my sandwich?}
    \item \eng{He earned \textbf{a lot of moneys} last year.}
\end{enumerate}

\tcblower
\textbf{Solution détaillée :}

\begin{enumerate}
    \item \eng{She gave me \textbf{a piece of advice} for my interview.} \\
    (\textit{Advice} est indénombrable $\rightarrow$ \textbf{a piece of}. Autre solution avec le verbe : \eng{She \textbf{advised} me for my interview.})
    \item \eng{We bought \textbf{two pieces of furniture} for the living room.} \\
    (\textit{Furniture} est indénombrable $\rightarrow$ \textbf{pieces of}. Jamais *\eng{furnitures}.)
    \item \eng{He has \textbf{long hair} and a beard.} \\
    (\textit{Hair} désignant la chevelure est indénombrable et singulier. \eng{A hair} = un seul cheveu.)
    \item \eng{I need \textbf{a piece of information} about the train schedule.} \\
    (\textit{Information} est indénombrable $\rightarrow$ \textbf{a piece of}. Autre solution : \eng{Please \textbf{inform} me about the train schedule.})
    \item \eng{She cooked \textbf{spaghetti} (ou \textbf{a plate of spaghetti}) for dinner.} \\
    (\textit{Spaghetti} est indénombrable en anglais. Jamais *\eng{a spaghetti}.)
    \item \eng{They have \textbf{a lot of luggage} (ou \textbf{much luggage}).} \\
    (\textit{Luggage} est indénombrable. Jamais *\eng{luggages}. \textit{Many} est réservé aux comptables ; avec un indénombrable, on utilise \textit{much} ou \textit{a lot of}.)
    \item \eng{Can I have \textbf{a slice of bread} (ou \textbf{a loaf of bread}) for my sandwich?} \\
    (\textit{Bread} est indénombrable. *\eng{A bread} n'existe pas. \eng{A slice} = une tranche ; \eng{a loaf} = un pain entier.)
    \item \eng{He earned \textbf{a lot of money} last year.} \\
    (\textit{Money} est indénombrable. Jamais *\eng{moneys}. On utilise \textit{much} ou \textit{a lot of}.)
\end{enumerate}
\end{exoresolu}

\courssub{Application immédiate : mini-production écrite}

\begin{exercice}[Mise en situation : au supermarché et à la maison]
Complétez les phrases avec le partitif anglais correct, choisi dans la liste suivante : \eng{a piece of, a slice of, a loaf of, a bunch of, a cup of, a bar of, a tube of, an item of, a large sum of, a hair}. Traduisez ensuite chaque phrase en français.

\begin{enumerate}
    \item For breakfast, I usually have \_\_\_\_\_\_\_\_\_\_ toast and \_\_\_\_\_\_\_\_\_\_ tea.
    \item Could you buy \_\_\_\_\_\_\_\_\_\_ bananas and \_\_\_\_\_\_\_\_\_\_ bread at the market?
    \item The mechanic gave me \_\_\_\_\_\_\_\_\_\_ advice about the engine.
    \item She spent \_\_\_\_\_\_\_\_\_\_ money on her new phone.
    \item There is \_\_\_\_\_\_\_\_\_\_ in my soup!
    \item We need \_\_\_\_\_\_\_\_\_\_ equipment for the camping trip.
    \item He ate \_\_\_\_\_\_\_\_\_\_ chocolate after lunch.
    \item Please squeeze \_\_\_\_\_\_\_\_\_\_ toothpaste onto the brush.
\end{enumerate}
\end{exercice}

\begin{corrige}[Exercice : au supermarché et à la maison]
\begin{enumerate}
    \item \eng{a slice of} toast and \eng{a cup of} tea. — « Une tranche de pain grillé et une tasse de thé. »
    \item \eng{a bunch of} bananas and \eng{a loaf of} bread. — « Une grappe de bananes et un pain. »
    \item \eng{a piece of} advice. — « Le mécanicien m'a donné un conseil au sujet du moteur. »
    \item \eng{a large sum of} money. — « Elle a dépensé une grosse somme d'argent pour son nouveau téléphone. »
    \item \eng{a hair}. — « Il y a un cheveu dans ma soupe ! » (singulier, car un seul cheveu).
    \item \eng{an item of} equipment (ou \eng{a piece of} equipment). — « Nous avons besoin d'un équipement pour le camping. »
    \item \eng{a bar of} chocolate. — « Il a mangé une tablette de chocolat après le déjeuner. »
    \item \eng{a tube of} toothpaste. — « Mets, s'il te plaît, un tube de dentifrice sur la brosse. » (au sens figuré : presse le tube ; littéralement « presse un tube de dentifrice sur la brosse »).
\end{enumerate}
\end{corrige}

\coursec{Emplois en contexte et cas particuliers}

Nous avons vu dans la section précédente que le français et l'anglais partagent la même logique partitive (\eng{a bottle of water} = « une bouteille d'eau »), mais que l'anglais impose cette structure là où le français peut s'en passer. Il reste maintenant à répondre à la question essentielle : \textbf{quand et comment utilise-t-on ces expressions dans la vraie vie ?} C'est l'objet de cette section : les situations concrètes — le marché, la cuisine, la description — puis les cas particuliers qui font la différence à l'examen.

\courssub{Au marché et à la boutique : commander et demander les prix}

\begin{textebox}[At the market — a short dialogue]
\textbf{Customer:} Good morning, madam. How much is a bag of rice?\\
\textbf{Trader:} A small bag is 5,000 francs, and a big bag is 12,000 francs.\\
\textbf{Customer:} I see. And how much is a bottle of groundnut oil?\\
\textbf{Trader:} One litre costs 1,500 francs.\\
\textbf{Customer:} All right. I'd like a tin of tomatoes, a packet of salt and two bars of soap, please.\\
\textbf{Trader:} Here you are. Anything else? A bunch of plantains, perhaps? They are very fresh.\\
\textbf{Customer:} Yes, give me a bunch of plantains and a piece of that smoked fish. How much is everything?\\
\textbf{Trader:} That will be 9,200 francs.\\
\textbf{Customer:} Here is the money. Thank you very much.\\
\textbf{Trader:} You're welcome. Come again!
\end{textebox}

\textbf{Glossaire :} \emph{trader} (n.) : commerçant(e) ; \emph{bar of soap} : savonnette ; \emph{bunch} (n.) : régime, bouquet ; \emph{smoked fish} : poisson fumé ; \emph{Here you are} : voici (en tendant quelque chose) ; \emph{Anything else?} : autre chose ?

\begin{propriete}[Demander le prix et commander : les structures clés]
\begin{itemize}
\item \textbf{Demander le prix :} \eng{How much is a bag of rice?} / \eng{How much are two bottles of oil?} — le verbe s'accorde avec le \cle{partitif} (singulier ou pluriel), pas avec le nom indénombrable.
\item \textbf{Commander poliment :} \eng{I'd like a tin of tomatoes, please.} (= I would like…). On peut aussi dire \eng{Can I have a packet of sugar?} ou \eng{Give me a piece of fish, please.}
\item \textbf{Annoncer le total :} \eng{That will be 9,200 francs.}
\end{itemize}
\end{propriete}

\begin{exemplebox}[Exemples traduits]
\eng{How much is a bag of cement?} « Combien coûte un sac de ciment ? »\\
\eng{I'd like a tin of milk and a loaf of bread.} « Je voudrais une boîte de lait et une miche de pain. »\\
\eng{Can I have a glass of water, please?} « Puis-je avoir un verre d'eau, s'il vous plaît ? »\\
\eng{Two bottles of palm wine cost 1,000 francs at this shop.} « Deux bouteilles de vin de palme coûtent 1 000 francs dans cette boutique. »
\end{exemplebox}

\begin{attention}[Accord du verbe : le piège classique]
Dans \eng{How much \textbf{is} a bag of rice?}, le verbe est au singulier parce que \eng{bag} est singulier, même si \eng{rice} est indénombrable. Mais attention : \eng{How much \textbf{are} two bags of rice?} — le verbe devient pluriel avec le partitif pluriel. Ne dites jamais \eng{How much are a bag of rice?}.\\
Autre piège : \eng{how much} s'emploie avec les noms indénombrables et pour les prix ; \eng{how many} avec les noms dénombrables pluriels : \eng{How many bottles of oil did you buy?} « Combien de bouteilles d'huile as-tu achetées ? »
\end{attention}

\courssub{Dans les recettes de cuisine}

\begin{textebox}[Recipe — Groundnut soup]
Pour a bottle of water into the pot. Add a spoonful of salt, two cups of groundnut paste and a tin of tomatoes. Stir well and serve with a bunch of boiled plantains.
\end{textebox}

\textbf{Glossaire :} \emph{to pour} (v.) : verser ; \emph{pot} (n.) : marmite ; \emph{spoonful} (n.) : cuillerée ; \emph{to stir} (v.) : remuer ; \emph{to serve} (v.) : servir.

\begin{propriete}[Les partitifs de la cuisine]
Les recettes anglaises utilisent des \cle{mesures} précises devant les noms indénombrables :
\begin{itemize}
\item \eng{a spoonful of salt} = une cuillerée de sel (notez le suffixe \cle{-ful} : \eng{cupful}, \eng{mouthful}, \eng{handful} ; pluriel : \eng{spoonfuls})
\item \eng{a cup of flour} = une tasse de farine
\item \eng{a tin of tomatoes} = une boîte de tomates
\item \eng{a slice of bread} = une tranche de pain
\item \eng{a pinch of pepper} = une pincée de poivre
\end{itemize}
L'impératif (\eng{Pour… Add… Stir…}) est le mode habituel des recettes, comme en français (« Versez… Ajoutez… »).
\end{propriete}

\begin{exemplebox}[Exemples traduits]
\eng{Add a spoonful of salt and a cup of flour.} « Ajoutez une cuillerée de sel et une tasse de farine. »\\
\eng{Cut the meat into small pieces.} « Coupez la viande en petits morceaux. »\\
\eng{She put a lump of sugar in my tea.} « Elle a mis un morceau de sucre dans mon thé. »
\end{exemplebox}

\courssub{Dans les descriptions}

Les partitifs servent aussi à décrire des scènes, notamment avec \eng{there is / there are} :

\begin{exemplebox}[Exemples traduits]
\eng{There is a crowd of people in front of the bank.} « Il y a une foule de gens devant la banque. »\\
\eng{There is a group of students near the school gate.} « Il y a un groupe d'élèves près du portail de l'école. »\\
\eng{A team of doctors arrived in Bamenda yesterday.} « Une équipe de médecins est arrivée à Bamenda hier. »\\
\eng{There is a pile of books on the teacher's desk.} « Il y a une pile de livres sur le bureau du professeur. »
\end{exemplebox}

\begin{attention}[There is ou there are ?]
Avec un partitif singulier, on dit toujours \eng{there \textbf{is}}, même si le nom qui suit est pluriel : \eng{There \textbf{is} a bunch of flowers on the table} (et non \eng{there are}). Le verbe s'accorde avec \eng{bunch}, le nom de tête. En revanche : \eng{There \textbf{are} two bunches of flowers on the table.}
\end{attention}

\courssub{Partitifs avec les noms dénombrables pluriels}

Jusqu'ici, les partitifs introduisaient surtout des noms \textbf{indénombrables} (\eng{water, rice, advice}). Mais ils s'emploient aussi devant des noms \textbf{dénombrables au pluriel} pour former des collections :

\begin{center}
\begin{tabularx}{\linewidth}{|X|X|X|}
\hline
\rowcolor{popblueL} Partitif & English & Français \\
\hline
\eng{a bunch of} & a bunch of flowers / keys & un bouquet de fleurs / un trousseau de clés \\
\hline
\eng{a pair of} & a pair of shoes / scissors / trousers & une paire de chaussures / de ciseaux / un pantalon \\
\hline
\eng{a set of} & a set of tools / rules & un jeu d'outils / un ensemble de règles \\
\hline
\eng{a packet of} & a packet of cigarettes / biscuits & un paquet de cigarettes / de biscuits \\
\hline
\eng{a group of} & a group of students & un groupe d'élèves \\
\hline
\eng{a crowd of} & a crowd of people & une foule de gens \\
\hline
\eng{a team of} & a team of players & une équipe de joueurs \\
\hline
\end{tabularx}
\end{center}

\begin{attention}[\eng{A pair of} ou \eng{a couple of} ?]
Ne confondez pas ces deux expressions :
\begin{itemize}
\item \eng{\textbf{a pair of}} = deux éléments \textbf{liés} qui vont ensemble ou forment un seul objet : \eng{a pair of shoes} (une paire de chaussures), \eng{a pair of scissors} (des ciseaux), \eng{a pair of trousers} (un pantalon), \eng{a pair of glasses} (des lunettes).
\item \eng{\textbf{a couple of}} = deux éléments \textbf{quelconques}, environ deux : \eng{a couple of friends} (deux amis), \eng{a couple of days} (deux ou trois jours).
\end{itemize}
On dit \eng{a couple of friends}, jamais \eng{a pair of friends} !
\end{attention}

\courssub{Registre familier et registre soutenu}

\begin{propriete}[Choisir le bon registre]
\begin{itemize}
\item \textbf{Registre familier} (conversation quotidienne) : \eng{a bit of} + indénombrable (\eng{a bit of time} = un peu de temps, \eng{a bit of luck} = un peu de chance, \eng{a bit of money}) ; \eng{a lot of} et \eng{lots of} (beaucoup de) pour les quantités.
\item \textbf{Registre soutenu} (presse, dissertation, écrit formel) : \eng{an item of news} (une nouvelle), \eng{a piece of research} (une recherche), \eng{a piece of information} (une information), \eng{a piece of advice} (un conseil), \eng{a work of art} (une œuvre d'art).
\end{itemize}
\end{propriete}

\begin{exemplebox}[Exemples traduits]
\eng{Wait a bit — I need a bit of time.} « Attends un peu — j'ai besoin d'un peu de temps. »\\
\eng{She gave me a piece of advice.} « Elle m'a donné un conseil. »\\
\eng{That is an interesting item of news.} « C'est une nouvelle intéressante. »\\
\eng{There were lots of people at the football match.} « Il y avait beaucoup de monde au match de football. »
\end{exemplebox}

\courssub{Questions et négations}

Les partitifs fonctionnent normalement dans les questions et les phrases négatives, avec \eng{there is / there isn't} :

\begin{exemplebox}[Exemples traduits]
\eng{Is there a bottle of oil left?} « Reste-t-il une bouteille d'huile ? »\\
\eng{There isn't a drop of water in the house.} « Il n'y a pas une goutte d'eau dans la maison. »\\
\eng{There isn't a drop of water left.} « Il ne reste pas une goutte d'eau. »\\
\eng{Have you got a piece of paper?} « As-tu une feuille de papier ? »\\
\eng{There isn't a grain of truth in his story.} « Il n'y a pas un grain de vérité dans son histoire. »
\end{exemplebox}

Notez l'expression négative forte \eng{not a drop of} / \eng{not a grain of} / \eng{not a word of} : \eng{She didn't say a word.} « Elle n'a pas dit un mot. »

\courssub{Frise des contextes d'emploi}

\begin{popfigure}
\begin{center}
\begin{tikzpicture}[font=\small]
\draw[-{Stealth}, thick, popline] (0,0) -- (12.5,0);
\foreach \x/\titre/\ex in {1.5/{Market}/a bag of rice, 4.8/{Recipe}/a cup of flour, 8.1/{Description}/a crowd of people, 11.2/{Abstract}/a piece of advice}{
  \fill[popblue] (\x,0) circle (0.09);
  \node[above, align=center, text=popdark] at (\x,0.15) {\textbf{\titre}};
  \node[below, align=center, text=popmuted] at (\x,-0.15) {\emph{\ex}};
}
\end{tikzpicture}
\end{center}
\legende{Les quatre grands contextes d'emploi des partitifs, du plus concret au plus abstrait.}
\end{popfigure}

\courssub{Pour aller plus loin : les partitifs collectifs d'animaux (hors programme strict)}

\begin{savaistu}[Un bonus utile pour le Bac]
Les collectifs d'animaux ne figurent pas explicitement dans le programme de Première, mais ils apparaissent \textbf{souvent dans les textes de compréhension écrite} du baccalauréat. Les connaître, c'est gagner des points faciles :
\begin{itemize}
\item \eng{a flock of sheep} = un troupeau de moutons (aussi : \eng{a flock of birds} = une volée d'oiseaux)
\item \eng{a herd of cattle} = un troupeau de bovins (aussi : \eng{a herd of elephants})
\item \eng{a swarm of bees} = un essaim d'abeilles
\item \eng{a pack of wolves} = une meute de loups (aussi : \eng{a pack of dogs})
\item \eng{a school of fish} = un banc de poissons
\end{itemize}
Exemple : \eng{A swarm of bees attacked the farmers near Buea.} « Un essaim d'abeilles a attaqué les agriculteurs près de Buéa. »
\end{savaistu}

\begin{methode}[Réussir les « fill in the blanks » au Bac]
Dans un exercice de type \emph{fill in the blanks} (texte à trous), procédez en trois étapes :
\begin{enumerate}
\item \textbf{Repérez le nom} qui suit le trou : est-il indénombrable (\eng{water, news, bread, money}) ou dénombrable pluriel (\eng{flowers, sheep, tools}) ?
\item \textbf{Déterminez la logique} : s'agit-il d'un \emph{contenant} (\eng{bottle, cup, tin, bag, glass}), d'une \emph{portion} (\eng{piece, slice, drop, lump, spoonful}) ou d'une \emph{collection} (\eng{bunch, pair, set, flock, herd}) ?
\item \textbf{Vérifiez l'accord} : \eng{a} + partitif singulier ; si le sujet est au pluriel, mettez le partitif au pluriel (\eng{two bottles of oil}).
\end{enumerate}
Exemple : \eng{She bought a \ldots{} of bread.} → \eng{bread} est indénombrable, on veut une portion → \eng{a \textbf{loaf} of bread} ou \eng{a \textbf{slice} of bread}.
\end{methode}

\begin{exoresolu}[Exercice résolu — fill in the blanks]
Énoncé : complétez avec le partitif qui convient : \emph{bunch, drop, piece, pair, flock, tin}.\\
1. There isn't a \ldots{} of water left in the tank.\\
2. My father bought a \ldots{} of shoes for my brother.\\
3. The teacher gave us a useful \ldots{} of advice before the exam.\\
4. A \ldots{} of sheep was crossing the road near Bamenda.\\
5. She offered me a beautiful \ldots{} of flowers.\\
6. Open a \ldots{} of tomatoes for the sauce, please.
\tcblower
Solution détaillée :\\
1. \textbf{drop} — \eng{water} est indénombrable ; la négation \eng{there isn't} appelle la petite portion : \eng{There isn't a drop of water left.} « Il ne reste pas une goutte d'eau. »\\
2. \textbf{pair} — deux éléments liés : \eng{a pair of shoes} « une paire de chaussures » (pas \eng{a couple of shoes} !).\\
3. \textbf{piece} — \eng{advice} est indénombrable en anglais : \eng{a piece of advice} « un conseil ».\\
4. \textbf{flock} — collectif des moutons : \eng{a flock of sheep} « un troupeau de moutons ».\\
5. \textbf{bunch} — collection de fleurs : \eng{a bunch of flowers} « un bouquet de fleurs ».\\
6. \textbf{tin} — contenant métallique : \eng{a tin of tomatoes} « une boîte de tomates ».
\end{exoresolu}

\begin{exercice}[À vous de jouer]
Traduisez en anglais : 1. Combien coûte un sac de riz ? 2. Je voudrais une bouteille d'huile et un paquet de sel. 3. Il y a une foule de gens devant la banque. 4. Ajoutez une cuillerée de sel. 5. Elle m'a donné un conseil. 6. Il ne reste pas une goutte d'eau dans la maison.
\end{exercice}

\begin{corrige}[Exercice « À vous de jouer »]
1. \eng{How much is a bag of rice?} 2. \eng{I'd like a bottle of oil and a packet of salt.} 3. \eng{There is a crowd of people in front of the bank.} 4. \eng{Add a spoonful of salt.} 5. \eng{She gave me a piece of advice.} 6. \eng{There isn't a drop of water left in the house.}
\end{corrige}

\coursec{Exercices d'application et de transformation}

Nous venons d'étudier les emplois en contexte et les cas particuliers des expressions partitives : vous savez désormais qu'on ne compte pas directement \eng{bread}, \eng{information} ou \eng{money}, mais qu'on compte \eng{slices of bread}, \eng{pieces of information} et \eng{sums of money}. Il est temps de vérifier et de consolider ces acquis par un entraînement complet, du plus simple (repérer) au plus exigeant (produire un texte type Bac). Avant de commencer, voici la démarche générale à adopter face à n'importe quelle phrase contenant un nom.

\begin{methode}[Comment choisir et corriger : la démarche en 4 étapes]
Face à une phrase à compléter ou à corriger, suivez toujours le même raisonnement :
\begin{enumerate}
\item \cle{Repérer} le nom : est-il dénombrable (\eng{a banana, two bananas}) ou indénombrable (\eng{bread, water, information}) ?
\item Si le nom est indénombrable, \cle{supprimer} tout article indéfini (\eng{a/an}) ou tout \eng{-s} fautif : on ne dit jamais \eng{an information} ni \eng{informations}.
\item \cle{Insérer} le partitif adapté à la nature du nom : \eng{a piece of information}, \eng{a loaf of bread}, \eng{a glass of water}.
\item \cle{Vérifier} l'accord du verbe : il s'accorde avec le partitif, pas avec le nom indénombrable.
\end{enumerate}
\end{methode}

\begin{popfigure}
\begin{tikzpicture}[
node distance=1.1cm,
startstop/.style={rectangle, rounded corners, draw=popblue, fill=popblueL, text width=4.2cm, align=center, minimum height=1cm, font=\small},
decision/.style={diamond, draw=poporange, fill=poporangeL, text width=3.2cm, align=center, inner sep=1pt, font=\small},
process/.style={rectangle, draw=poppurple, fill=poppurpleL, text width=4.6cm, align=center, minimum height=1cm, font=\small},
arrow/.style={-{Stealth}, thick, popdark}]
\node[startstop, align=center] (start){Nom à exprimer\\ \eng{bread ? information ? banana ?}};
\node[decision, below=of start] (count) {Dénombrable ?};
\node[process, below left=1.1cm and 1.2cm of count, align=center] (yes){OUI : \eng{a/an} + nom\\ \eng{a banana, two bananas}};
\node[process, below right=1.1cm and 1.2cm of count, align=center] (no){NON : partitif + \eng{of} + nom\\ \eng{a piece of bread}};
\node[decision, below=1.1cm of no] (plural) {Pluriel ?};
\node[process, below=of plural, align=center] (s){\eng{-s} sur le partitif uniquement\\ \eng{two pieces of bread}};
\draw[arrow] (start) -- (count);
\draw[arrow] (count) -| node[above left, font=\small\itshape]{oui} (yes);
\draw[arrow] (count) -| node[above right, font=\small\itshape]{non} (no);
\draw[arrow] (no) -- (plural);
\draw[arrow] (plural) -- node[right, font=\small\itshape]{oui} (s);
\end{tikzpicture}
\legende{Organigramme de résolution : choisir entre article simple et expression partitive.}
\end{popfigure}

\begin{propriete}[Rappel : l'accord du verbe avec le partitif]
Le verbe s'accorde en nombre avec le \cle{partitif} (le récipient ou l'unité), jamais avec le nom indénombrable qui suit \eng{of} :
\begin{itemize}
\item \eng{A bottle of oil costs 1,500 FCFA.} « Une bouteille d'huile coûte 1 500 FCFA. » (une bouteille $\rightarrow$ singulier)
\item \eng{Two bottles of oil cost 3,000 FCFA.} « Deux bouteilles d'huile coûtent 3 000 FCFA. » (deux bouteilles $\rightarrow$ pluriel)
\end{itemize}
\end{propriete}

\courssub{Exercice 1 : Identification dans un texte de marché}

\begin{textebox}[Scène de marché à Yaoundé]
\eng{Every Saturday, Mama Ngono goes to Mokolo market in Yaoundé. She buys a bunch of plantains, a bag of rice and a bottle of groundnut oil. Then she stops at the bakery for a loaf of bread. Before going home, she asks the butcher for a piece of beef and buys a bar of soap at the provision store.}
\end{textebox}

\begin{exoresolu}[Souligner les expressions partitives]
\textbf{Énoncé.} Lisez le texte ci-dessus et soulignez les six expressions partitives.
\tcblower
\textbf{Solution.} Les six expressions partitives sont :
\begin{enumerate}
\item \eng{a \textbf{bunch} of plantains} — une grappe de plantains
\item \eng{a \textbf{bag} of rice} — un sac de riz
\item \eng{a \textbf{bottle} of groundnut oil} — une bouteille d'huile d'arachide
\item \eng{a \textbf{loaf} of bread} — un pain
\item \eng{a \textbf{piece} of beef} — un morceau de bœuf
\item \eng{a \textbf{bar} of soap} — une savonnette
\end{enumerate}
\textbf{Astuce de repérage :} cherchez la structure \textbf{a/an + nom-unité + of + nom indénombrable}. Le mot qui précède \eng{of} est presque toujours le partitif.
\end{exoresolu}

\courssub{Exercice 2 : Appariement partitif / nom}

\begin{exercice}[À chaque partitif son nom]
Reliez chaque partitif de la colonne A au nom qui convient dans la colonne B.
\begin{center}
\begin{tabularx}{\linewidth}{|X|X|}
\hline
\rowcolor{popblueL} \textbf{A — Partitif} & \textbf{B — Nom} \\
\hline
1. a bunch of & a. bread \\
\hline
2. a loaf of & b. soap \\
\hline
3. a bar of & c. bananas \\
\hline
4. a sheet of & d. paper \\
\hline
\end{tabularx}
\end{center}
\end{exercice}

\begin{corrige}[Exercice 2]
1-c : \eng{a bunch of bananas} « un régime de bananes » ; 2-a : \eng{a loaf of bread} « un pain » ; 3-b : \eng{a bar of soap} « une savonnette » ; 4-d : \eng{a sheet of paper} « une feuille de papier ».
\textbf{Clé mémorisation :} le partitif épouse la \emph{forme} de l'objet : les bananes viennent en grappe (\eng{bunch}), le pain en miche (\eng{loaf}), le savon en barre (\eng{bar}), le papier en feuille (\eng{sheet}).
\end{corrige}

\courssub{Exercice 3 : Transformation singulier $\rightarrow$ pluriel}

\begin{exercice}[Mettez au pluriel]
Transformez les expressions suivantes avec le nombre indiqué entre parenthèses. Attention : seul le partitif prend la marque du pluriel.
\begin{enumerate}
\item a cup of tea (3)
\item a loaf of bread (2)
\item a bottle of palm wine (4)
\item a piece of advice (2)
\item a glass of water (3)
\end{enumerate}
\end{exercice}

\begin{corrige}[Exercice 3]
\begin{enumerate}
\item \eng{three cups of tea} « trois tasses de thé »
\item \eng{two loaves of bread} « deux pains » — attention au pluriel irrégulier \eng{loaf} $\rightarrow$ \eng{loaves} (comme \eng{leaf/leaves}).
\item \eng{four bottles of palm wine} « quatre bouteilles de vin de palme »
\item \eng{two pieces of advice} « deux conseils » — \eng{advice} reste sans \eng{-s}.
\item \eng{three glasses of water} « trois verres d'eau » — \eng{glass} fait son pluriel en \eng{-es}.
\end{enumerate}
\end{corrige}

\begin{attention}[Le -s ne va jamais sur le nom indénombrable]
Au pluriel, c'est le \cle{partitif} qui change, pas le nom : \eng{two \textbf{cups} of tea}, jamais \eng{two cup of teas}. Le nom indénombrable (\eng{tea, bread, water, advice}) reste invariable.
\end{attention}

\courssub{Exercice 4 : Correction d'erreurs de francophones}

\begin{exercice}[Corrige ces phrases]
Chaque phrase contient une erreur typique des francophones. Corrige-la en appliquant la méthode en 4 étapes.
\begin{enumerate}
\item \eng{I need an information.}
\item \eng{She bought two breads at the bakery.}
\item \eng{He gave me good advices.}
\item \eng{We drank three waters.}
\end{enumerate}
\end{exercice}

\begin{corrige}[Exercice 4]
\begin{enumerate}
\item \eng{I need an information.} $\rightarrow$ \eng{I need a piece of information.} « J'ai besoin d'une information. » (\eng{information} est indénombrable : pas d'article \eng{an}.)
\item \eng{She bought two breads.} $\rightarrow$ \eng{She bought two loaves of bread.} « Elle a acheté deux pains. »
\item \eng{He gave me good advices.} $\rightarrow$ \eng{He gave me some good advice.} ou \eng{He gave me two pieces of advice.} « Il m'a donné de bons conseils. »
\item \eng{We drank three waters.} $\rightarrow$ \eng{We drank three glasses of water.} « Nous avons bu trois verres d'eau. »
\end{enumerate}
\end{corrige}

\begin{attention}[Le piège de \eng{news}]
Dans \eng{a piece of news} « une nouvelle », le mot \eng{news} garde \cle{toujours} son \eng{-s} final, même au singulier : ce n'est \emph{pas} un pluriel ! On dit \eng{The news \textbf{is} good} « Les nouvelles sont bonnes » avec un verbe au singulier. Même phénomène avec \eng{mathematics} et \eng{physics}.
\end{attention}

\courssub{Exercice 5 : Traduction français $\rightarrow$ anglais}

\begin{exercice}[Traduisez en anglais]
\begin{enumerate}
\item un verre d'eau
\item deux bouteilles d'huile
\item un conseil
\item une tranche de pain
\item un trousseau de clés
\end{enumerate}
\end{exercice}

\begin{corrige}[Exercice 5]
\begin{enumerate}
\item \eng{a glass of water}
\item \eng{two bottles of oil}
\item \eng{a piece of advice} — jamais \eng{an advice} !
\item \eng{a slice of bread} — \eng{slice} « tranche » se prononce « slaïss ».
\item \eng{a bunch of keys} — \eng{bunch} sert aussi pour les objets attachés ensemble (clés, fleurs).
\end{enumerate}
\end{corrige}

\begin{center}
\begin{tabularx}{\linewidth}{|X|X|X|}
\hline
\rowcolor{popblueL} \textbf{Français} & \textbf{English} & \textbf{Remarque} \\
\hline
un verre d'eau & a glass of water & \eng{glass}, pas \eng{cup} (réservé aux tasses) \\
\hline
deux bouteilles d'huile & two bottles of oil & \eng{-s} sur \eng{bottles} seulement \\
\hline
un conseil & a piece of advice & \eng{advice} indénombrable, sans \eng{-s} \\
\hline
une tranche de pain & a slice of bread & \eng{slice} = tranche fine \\
\hline
un trousseau de clés & a bunch of keys & \eng{bunch} = ensemble attaché \\
\hline
une information & a piece of news / information & \eng{news} garde son \eng{-s} \\
\hline
\end{tabularx}
\end{center}

\courssub{Exercice 6 : Production guidée — ma liste de courses}

\begin{experience}[My shopping list]
Rédigez votre propre liste de courses pour le marché de votre quartier (Mokolo, Marché Central, Buea market...) : \textbf{8 articles}, avec \textbf{8 partitifs différents}. Commencez par : \eng{This weekend, I am going to the market. I need to buy...}

\textbf{Modèle attendu (exemple de production) :}

\eng{This weekend, I am going to Mokolo market. I need to buy a bag of rice, a bunch of plantains, two bottles of palm oil, a loaf of bread, a bar of soap, a tin of tomatoes, a packet of sugar and a piece of smoked fish. I hope the prices will not be too high!}

« Ce week-end, je vais au marché Mokolo. Je dois acheter un sac de riz, un régime de plantains, deux bouteilles d'huile de palme, un pain, une savonnette, une boîte de tomates, un paquet de sucre et un morceau de poisson fumé. J'espère que les prix ne seront pas trop élevés ! »
\end{experience}

\courssub{Exercice 7 : Type Bac — texte à trous}

\begin{exercice}[Fill in the blanks with the correct partitive]
Complétez le texte avec les partitifs suivants : \eng{slice, piece, cup, bottle, bunch, loaf, sheet, glass} (chacun une seule fois).

\eng{Mr Etoundi runs a small provision store in Bafoussam. Every morning, he drinks a (1) ......... of coffee before opening his shop. His first customer, Mrs Ateba, buys a (2) ......... of bread and a (3) ......... of groundnut oil for the family lunch. A pupil from the neighbouring secondary school asks for a (4) ......... of paper to write a letter. Later, a teacher requests a (5) ......... of chalk for his class. At noon, Mr Etoundi eats a (6) ......... of bread with sardines and drinks a (7) ......... of fresh water. In the afternoon, a fruit seller brings him a (8) ......... of bananas to resell. Business is good today!}
\end{exercice}

\begin{corrige}[Exercice 7]
(1) \eng{cup} — a cup of coffee « une tasse de café » ; (2) \eng{loaf} — a loaf of bread « un pain » ; (3) \eng{bottle} — a bottle of groundnut oil « une bouteille d'huile d'arachide » ; (4) \eng{sheet} — a sheet of paper « une feuille de papier » ; (5) \eng{piece} — a piece of chalk « un morceau de craie » ; (6) \eng{slice} — a slice of bread « une tranche de pain » ; (7) \eng{glass} — a glass of fresh water « un verre d'eau fraîche » ; (8) \eng{bunch} — a bunch of bananas « un régime de bananes ».

\textbf{Compétence évaluée :} choisir le partitif en fonction de la nature et de la forme de l'objet, dans un texte cohérent sur la vie économique — exactement l'exercice de grammaire en contexte proposé à l'examen.
\end{corrige}

\begin{savaistu}[Les mesures traditionnelles au marché]
Au Cameroun anglophone comme francophone, les unités de mesure traditionnelles coexistent avec le système métrique. On achète encore souvent \eng{a "cup" of rice} (une tasse/tasseau, environ 250--300 g), \eng{a "bunch" of plantains} (un régime, 8--12 doigts), \eng{a "heap" of tomatoes/oignons} (un tas, volume variable). Ces mesures "de femme de marché" sont codifiées socialement : le \eng{cup} (souvent une boîte de conserve vide type lait concentré ou tomates) sert d'étalon. En anglais camerounais, on dira \eng{"Give me two cups of rice"} ou \eng{"One bunch of plantains"}. Le \eng{bunch} pour les clés (\eng{a bunch of keys}) est une image transposée du régime de bananes : objets attachés ensemble.
\end{savaistu}

\begin{aretenir}[L'essentiel pour réussir tous les exercices]
\begin{itemize}
\item Nom indénombrable $\rightarrow$ jamais de \eng{a/an} ni de \eng{-s} : utilisez un \cle{partitif} (\eng{a piece of, a bottle of, a bunch of...}).
\item Au pluriel, le \eng{-s} porte sur le partitif seul : \eng{two loaves of bread}.
\item Le verbe s'accorde avec le partitif : \eng{A bottle of oil costs...} / \eng{Two bottles of oil cost...}.
\item \eng{news} garde toujours son \eng{-s}, même au singulier.
\end{itemize}
\end{aretenir}

\newpage

\coursec{Activité d'intégration}

\coursec{Lesson 23 — Grammar: Partitive expressions}

\courssub{Objectifs de la leçon}

À l'issue de cette leçon, tu seras capable de :

\begin{itemize}
\item identifier et nommer les \cle{expressions partitives} courantes (\eng{a bag of}, \eng{a bottle of}, \eng{a bunch of}, \eng{a loaf of}, \eng{a bar of}, \eng{a piece of}, \eng{a box of}, \eng{a packet of}, \eng{a tin of}, \eng{a cup of}) ;
\item construire correctement la structure \textbf{quantité + partitif + \eng{of} + nom} en respectant la nature dénombrable ou indénombrable du nom ;
\item placer le marqueur du pluriel sur le partitif et non sur le nom indénombrable (\eng{two bars of soap}, jamais \emph{two soaps}) ;
\item accorder le verbe avec le partitif dans les questions \eng{How much is/are...?} ;
\item éviter les erreurs typiques des francophones : omission de \eng{of}, traduction mot à mot de « un pain » (\emph{a bread}), « un conseil » (\emph{an advice}), « une savonnette » (\emph{a soap}) ;
\item utiliser ces expressions dans un dialogue d'achat/vente et rédiger un \cle{reçu} (\eng{receipt}) en anglais correct.
\end{itemize}

\courssub{Observation et découverte}

Nous commençons par l'analyse d'un document authentique : une liste de prix affichée dans un commerce bien réel de Yaoundé.

\begin{textebox}[MOKOLO BEST PROVISION STORE — Price List]
WELCOME TO MOKOLO BEST PROVISION STORE!

Fresh goods every morning — good prices for students!

\begin{itemize}
\item a bag of rice (5 kg) — 3,500 FCFA
\item a bottle of groundnut oil (1 litre) — 1,800 FCFA
\item a bunch of bananas — 1,000 FCFA
\item a loaf of bread — 300 FCFA
\item a bar of soap — 500 FCFA
\item a box of matches — 100 FCFA
\item a packet of sugar (1 kg) — 900 FCFA
\item a tin of sardines — 650 FCFA
\item a piece of advice — free, always free!
\end{itemize}

“Tell me exactly what you need, and I will serve you with a smile!”
\end{textebox}

\begin{definition}[Glossaire utile]
\emph{stallholder} (n.) : vendeur d'un étal ; \emph{to bargain} (v.) : marchander ; \emph{groundnut oil} (n.) : huile d'arachide ; \emph{a receipt} (n.) : un reçu (prononcé « ri-site », le \emph{p} ne se prononce pas) ; \emph{change} (n.) : la monnaie (à rendre) ; \emph{provision store} (n.) : épicerie, magasin de denrées alimentaires.
\end{definition}

\begin{experience}[Observation guidée]
Travaille seul pendant 3 minutes, puis compare avec ton voisin.
\begin{enumerate}
\item Combien d'expressions partitives différentes comptes-tu dans la liste ? Souligne-les.
\item Quel mot relie systématiquement le partitif (\eng{bag}, \eng{bottle}...) au nom qui suit (\eng{rice}, \eng{oil}...) ?
\item Repère trois noms qui, en français, sont dénombrables (\emph{un pain}, \emph{une savonnette}, \emph{un conseil}) mais qui, en anglais, apparaissent \textbf{sans article} après \eng{of}. Pourquoi ?
\end{enumerate}
\end{experience}

\courssub{La règle : structure et formation}

\begin{propriete}[Structure de base de l'expression partitive]
\begin{center}
\textbf{Quantificateur / Déterminant + Partitif + \eng{of} + Nom}
\end{center}

\begin{tabular}{|l|l|l|l|}
\hline
\rowcolor{popblueL} \textbf{Quantificateur} & \textbf{Partitif} & \textbf{\eng{of}} & \textbf{Nom} \\
\hline
\eng{a / an} & \eng{bag} & \eng{of} & \eng{rice} (U) \\
\hline
\eng{two / three / many} & \eng{bottles} & \eng{of} & \eng{oil} (U) \\
\hline
\eng{a} & \eng{loaf} & \eng{of} & \eng{bread} (U) \\
\hline
\eng{several} & \eng{boxes} & \eng{of} & \eng{matches} (C, pl.) \\
\hline
\end{tabular}

\medskip
\textbf{Règles fondamentales :}
\begin{enumerate}
\item \textbf{\eng{of} est obligatoire.} Il ne se contracte pas, ne s'élide pas. \eng{a bottle of oil} (jamais \emph{a bottle oil}).
\item \textbf{Le pluriel se marque sur le partitif (le conteneur/la portion), jamais sur le nom indénombrable (le contenu).}
  \begin{itemize}
  \item \eng{two bottles of oil} (pas \emph{two bottle of oils} ni \emph{two bottles of oils}).
  \item \eng{three loaves of bread} (pluriel irrégulier de \eng{loaf} $\rightarrow$ \eng{loaves}).
  \item \eng{five bars of soap} (pas \emph{five bars of soaps}).
  \end{itemize}
\item \textbf{Accord du verbe :} Le verbe s'accorde avec le \textbf{partitif} (sujet grammatical), pas avec le nom indénombrable.
  \eng{How much \textbf{is} a bag of rice?} \quad mais \quad \eng{How much \textbf{are} two bags of rice?}
\item \textbf{Noms dénombrables au pluriel après \eng{of} :} Si le contenu est un objet comptable, il prend la marque du pluriel.
  \eng{a box of matches}, \eng{a bunch of bananas}, \eng{a packet of biscuits}.
\end{enumerate}
\end{propriete}

\courssub{Les principaux partitifs par catégories}

Il est plus facile de les mémoriser en les classant selon leur fonction sémantique.

\begin{center}
\begin{tabularx}{\linewidth}{|X|X|X|}
\hline
\rowcolor{popblueL} \textbf{Contenants (Containers)} & \textbf{Portions / Formes (Portions / Shapes)} & \textbf{Regroupements (Collections)} \\
\hline
\eng{a bag of} (sac de) & \eng{a loaf of} (miche/pain de) & \eng{a bunch of} (grappe/botte de) \\
\eng{a bottle of} (bouteille de) & \eng{a bar of} (barre/tablette de) & \eng{a pack of} (paquet de - cartes, cigarettes) \\
\eng{a box of} (boîte de) & \eng{a piece of} (morceau/pièce de) & \eng{a set of} (ensemble de) \\
\eng{a tin of} / \eng{a can of} (boîte métallique de) & \eng{a slice of} (tranche de) & \eng{a collection of} (collection de) \\
\eng{a packet of} / \eng{a pack of} (paquet/sachet de) & \eng{a lump of} (morceau compact de) & \eng{a heap of} (tas de) \\
\eng{a jar of} (pot/bocal de) & \eng{a sheet of} (feuille de) & \eng{a crowd of} (foule de) \\
\eng{a cup of} (tasse de) & \eng{a strip of} (bande/lamelle de) & \eng{a flock of} (troupeau d'oiseaux) \\
\eng{a mug of} (bol/mug de) & \eng{a chunk of} (gros morceau de) & \eng{a herd of} (troupeau de bêtes) \\
\eng{a glass of} (verre de) & & \\
\eng{a bowl of} (bol de) & & \\
\hline
\end{tabularx}
\end{center}

\begin{aretenir}[Partitifs « passe-partout » pour indénombrables]
\eng{a piece of} et \eng{a bit of} (plus familier) s'utilisent avec une très grande variété de noms indénombrables : \eng{information, advice, furniture, luggage, news, work, research, evidence, chocolate, cake, paper, wood, metal}.
\eng{“A piece of advice”, “Two bits of information”, “Three pieces of furniture”.}
\end{aretenir}

\courssub{Comparaison français / anglais : points de vigilance}

\begin{attention}[Les 5 pièges majeurs pour les francophones]
\begin{enumerate}
\item \textbf{Oubli de \eng{of}.} Fr : « une bouteille d'huile » $\rightarrow$ \eng{a bottle \textbf{of} oil}. L'\emph{of} anglais n'est pas l'article contracté français \emph{d'} ; c'est une préposition pleine, toujours écrite.
\item \textbf{Pluriel mal placé.} Fr : « deux savonnettes » $\rightarrow$ \eng{two \textbf{bars of soap}}. Le mot \eng{soap} reste au singulier (indénombrable). Le pluriel est sur \eng{bars}.
\item \textbf{Traduction mot à mot de noms indénombrables anglais.}
\begin{center}
\begin{tabular}{|l|l|l|}
\hline
\rowcolor{poporangeL} \textbf{Français (dénombrable)} & \textbf{Anglais (indénombrable)} & \textbf{Partitif correct} \\
\hline
un pain & \eng{bread} & \eng{a loaf of bread} \\
un conseil & \eng{advice} & \eng{a piece of advice} \\
une savonnette & \eng{soap} & \eng{a bar of soap} \\
un meuble & \eng{furniture} & \eng{a piece of furniture} \\
un bagage & \eng{luggage} & \eng{a piece of luggage} \\
une information & \eng{information} & \eng{a piece of information} \\
un éclair / un éclaircie & \eng{lightning} & \eng{a flash of lightning} \\
\hline
\end{tabular}
\end{center}
\item \textbf{Accord \eng{How much is/are}.} Fr : « Combien coûte... ? » (verbe au singulier). En : \eng{How much \textbf{is} a bottle?} mais \eng{How much \textbf{are} three bottles?}. Le sujet est \eng{a bottle} / \eng{three bottles}.
\item \textbf{Orthographe de pluriels irréguliers.}
\eng{loaf $\rightarrow$ loaves}, \eng{knife $\rightarrow$ knives} (a knife of... rare), \eng{leaf $\rightarrow$ leaves}, \eng{match $\rightarrow$ matches} (ajout \emph{-es}), \eng{box $\rightarrow$ boxes}.
\end{enumerate}
\end{attention}

\courssub{Emplois en contexte et cas particuliers}

\begin{exemplebox}[Formules de politesse au marché / au magasin]
\textbf{Demander (Client) :}
\begin{itemize}
\item \eng{“I'd like \textbf{a bag of rice}, please.”} (standard, poli)
\item \eng{“Can I have \textbf{two bottles of groundnut oil}?”} (direct)
\item \eng{“How much \textbf{is} \textbf{a loaf of bread}?”} (prix unitaire)
\item \eng{“How much \textbf{are} \textbf{three packets of sugar}?”} (prix multiple)
\item \eng{“Could you give me \textbf{a piece of advice} on the best soap?”} (demande de conseil)
\end{itemize}

\textbf{Répondre / Vendre (Vendeur) :}
\begin{itemize}
\item \eng{“That'll be 3,500 francs, please.”} (= \eng{That will be})
\item \eng{“Here is your change.”} (\eng{change} indénombrable)
\item \eng{“Anything else?”} / \eng{“Will that be all?”}
\item \eng{“Thank you, come again!”}
\end{itemize}
\end{exemplebox}

\begin{exemplebox}[Cas particuliers : \eng{a piece of} + nom d'action (V-ing / nom verbal)]
Certains noms d'action sont indénombrables. Pour les compter, on utilise \eng{a piece of} ou \eng{a bit of}.
\eng{“He did a piece of work.”} (un travail / une œuvre)
\eng{“She gave me a bit of help.”} (un peu d'aide)
\eng{“We need a piece of legislation on this.”} (un texte de loi)
\end{exemplebox}

\begin{savaistu}[Le \eng{receipt} vs \eng{recipe}]
Attention au \cle{faux ami} : \eng{a receipt} (prononcé /rɪˈsiːt/ « ri-site ») = un reçu (preuve de paiement). \eng{A recipe} (prononcé /ˈrɛsəpi/ « ré-sé-pi ») = une recette de cuisine. Le \emph{p} de \eng{receipt} est muet (héritage latin \emph{recepta}).
\end{savaistu}

\courssub{Exercices d'application et de transformation}

\courssub{Activité d'intégration : Au marché Mokolo}

\begin{experience}[Situation de communication : Achat et vente]
\textbf{Contexte :} You are at Mokolo Market in Yaoundé. Stallholders sell everything: rice, groundnut oil, bananas, soap, matches and sugar. Customers bargain, greet the sellers and always specify quantities carefully. Today, you will be both a customer and a stallholder.

\textbf{Consignes :} Travail en binômes. L'activité se déroule en cinq étapes.

\begin{enumerate}
\item \textbf{Préparation du vocabulaire (5 min).} Reprends la liste de prix (Observation). Recopie les neuf expressions partitives en les classant dans le tableau ci-dessous :

\begin{center}
\begin{tabularx}{\linewidth}{|X|X|X|}
\hline
\rowcolor{popgreenL} \textbf{Contenants} & \textbf{Portions / Formes} & \textbf{Regroupements} \\
\hline
\eng{a bottle of} & \eng{a loaf of} & \eng{a bunch of} \\
\eng{a box of} & \eng{a bar of} & \\
\eng{a tin of} & \eng{a piece of} & \\
\eng{a packet of} & & \\
\eng{a bag of} & & \\
\hline
\end{tabularx}
\end{center}

\item \textbf{Préparation des rôles (5 min).}
\begin{itemize}
\item \textbf{Élève A (Vendeur) :} Prépare des étiquettes-cartons pour l'affichage : \eng{bags of rice}, \eng{bottles of oil}, \eng{tins of sardines}, \eng{boxes of matches}, \eng{packets of sugar}, \eng{bars of soap}, \eng{loaves of bread}, \eng{bunches of bananas}. Note : le pluriel est sur le partitif !
\item \textbf{Élève B (Client) :} Reçois la liste de courses ci-dessous (en français). Traduis-la mentalement en demandes anglaises correctes (ne l'écris pas mot à mot).
\end{itemize}

\begin{center}
\begin{tabularx}{\linewidth}{|X|X|}
\hline
\rowcolor{popblueL} \textbf{Liste de courses (français)} & \textbf{Demande attendue (anglais)} \\
\hline
un sac de riz & \eng{a bag of rice} \\
\hline
deux bouteilles d'huile & \eng{two bottles of oil} \\
\hline
une grappe de bananes & \eng{a bunch of bananas} \\
\hline
un pain & \eng{a loaf of bread} \\
\hline
une savonnette & \eng{a bar of soap} \\
\hline
un conseil sur un produit & \eng{a piece of advice} \\
\hline
une boîte d'allumettes & \eng{a box of matches} \\
\hline
un paquet de sucre & \eng{a packet of sugar} \\
\hline
une boîte de sardines & \eng{a tin of sardines} \\
\hline
\end{tabularx}
\end{center}

\item \textbf{Jeu de rôle (10 min).} Jouez la scène au moins deux fois.
\begin{itemize}
\item Le client salue, formule \textbf{chaque demande} avec un partitif correct (\eng{I'd like a bag of rice, please. How much are two bottles of oil?}).
\item Il demande un conseil (\eng{Could you give me a piece of advice about the oil?}).
\item Le vendeur salue, annonce les prix en FCFA, rend la monnaie, remercie.
\item \textbf{Contrainte :} Le dialogue doit comporter au moins \textbf{neuf partitifs différents} (les 8 de la liste + un libre).
\end{itemize}

\item \textbf{Inversion des rôles (10 min).} Échangez les rôles. Le nouveau client compose une liste contenant au moins \textbf{quatre partitifs différents} de la première liste (ex. \eng{a cup of tea}, \eng{a piece of chalk}, \eng{a bottle of water}, \eng{a jar of honey}, \eng{a pack of cards}...). Rejouez.

\item \textbf{Production écrite : le reçu (10 min).} Chaque binôme rédige le \cle{reçu} (\eng{receipt}) de la \textbf{première} transaction, en anglais, avec les articles, leurs partitifs, quantités, prix unitaires et le total. Modèle :

\begin{center}
\begin{tabular}{|l|r|r|r|}
\hline
\rowcolor{popgreenL} \textbf{MOKOLO BEST PROVISION STORE — RECEIPT} & \textbf{Qty} & \textbf{Unit Price} & \textbf{Total} \\
\hline
Bag of rice (5 kg) & 1 & 3,500 FCFA & 3,500 FCFA \\
\hline
Bottles of groundnut oil & 2 & 1,800 FCFA & 3,600 FCFA \\
\hline
Bunch of bananas & 1 & 1,000 FCFA & 1,000 FCFA \\
\hline
Loaf of bread & 1 & 300 FCFA & 300 FCFA \\
\hline
Bar of soap & 1 & 500 FCFA & 500 FCFA \\
\hline
Box of matches & 1 & 100 FCFA & 100 FCFA \\
\hline
Packet of sugar (1 kg) & 1 & 900 FCFA & 900 FCFA \\
\hline
Tin of sardines & 1 & 650 FCFA & 650 FCFA \\
\hline
Piece of advice & 1 & Free & 0 FCFA \\
\hline
\rowcolor{popgoldL} \textbf{TOTAL} & & & \textbf{10,550 FCFA} \\
\hline
\end{tabular}
\end{center}
\end{enumerate}
\end{experience}

\courssub{Exercices de transformation écrits (Travail individuel)}

\begin{exercice}[Transformation : Singulier $\rightarrow$ Pluriel]
Mets les phrases suivantes au pluriel (quantité = \eng{three} ou \eng{several}). Attention au pluriel du partitif et à l'accord du verbe.
\begin{enumerate}
\item \eng{There is a bag of rice in the car.}
\item \eng{She bought a bottle of perfume.}
\item \eng{A loaf of bread costs 300 francs.}
\item \eng{He needs a piece of advice.}
\item \eng{There is a box of matches on the shelf.}
\item \eng{A bar of soap is useful.}
\item \eng{She drank a cup of tea.}
\item \eng{A bunch of bananas is heavy.}
\end{enumerate}
\end{exercice}

\begin{corrige}[Exercice Transformation Singulier $\rightarrow$ Pluriel]
\begin{enumerate}
\item \eng{There \textbf{are three bags of rice} in the car.}
\item \eng{She bought \textbf{several bottles of perfume}.}
\item \eng{\textbf{Three loaves of bread cost} 900 francs.} (\eng{loaves} pluriel irrégulier, verbe \eng{cost} au pluriel)
\item \eng{He needs \textbf{several pieces of advice}.}
\item \eng{There \textbf{are several boxes of matches} on the shelf.}
\item \eng{\textbf{Several bars of soap are} useful.}
\item \eng{She drank \textbf{three cups of tea}.}
\item \eng{\textbf{Several bunches of bananas are} heavy.}
\end{enumerate}
\end{corrige}

\begin{exercice}[Traduction : Français $\rightarrow$ Anglais]
Traduis en anglais en utilisant la structure partitive correcte.
\begin{enumerate}
\item Je voudrais un kilo de riz (a bag of rice).
\item Combien coûtent deux pains ? (How much...?)
\item Elle a acheté trois savonnettes.
\item Donne-moi un conseil, s'il te plaît.
\item Il n'y a plus d'allumettes (matches) dans la boîte.
\item Nous avons bu deux tasses de thé.
\item Ce meuble est lourd (a piece of furniture).
\item Une tranche de gâteau, s'il vous plaît (a slice of cake).
\end{enumerate}
\end{exercice}

\begin{corrige}[Exercice Traduction]
\begin{enumerate}
\item \eng{I'd like a bag of rice.} (ou \eng{a kilo of rice} si unité de mesure, mais leçon porte sur \eng{bag of}).
\item \eng{How much \textbf{are} two loaves of bread?}
\item \eng{She bought three bars of soap.}
\item \eng{Could you give me a piece of advice, please?}
\item \eng{There are no matches left in the box.} / \eng{The box of matches is empty.}
\item \eng{We drank two cups of tea.}
\item \eng{This piece of furniture is heavy.}
\item \eng{A slice of cake, please.}
\end{enumerate}
\end{corrige}

\begin{exercice}[Correction d'erreurs : « Chasse aux fautes »]
Chaque phrase contient une erreur typique de francophone. Corrige-la et explique la règle.
\begin{enumerate}
\item \eng{I would like a bread, please.}
\item \eng{She gave me an advice.}
\item \eng{He bought two bottles of oils.}
\item \eng{How much is three boxes of matches?}
\item \eng{Can I have a soap?}
\item \eng{There are many furnitures in this room.}
\item \eng{A bunch of flowers are beautiful.}
\item \eng{I need a packet of sugars.}
\end{enumerate}
\end{exercice}

\begin{corrige}[Exercice Correction d'erreurs]
\begin{enumerate}
\item \eng{I would like \textbf{a loaf of bread}, please.} (\eng{bread} indénombrable $\rightarrow$ partitif \eng{loaf of}).
\item \eng{She gave me \textbf{a piece of advice}.} (\eng{advice} indénombrable $\rightarrow$ partitif \eng{piece of}).
\item \eng{He bought two bottles of \textbf{oil}.} (Nom indénombrable $\rightarrow$ pas de \emph{s} sur \eng{oil}).
\item \eng{How much \textbf{are} three boxes of matches?} (Sujet = \eng{three boxes} $\rightarrow$ pluriel $\rightarrow$ \eng{are}).
\item \eng{Can I have \textbf{a bar of soap}?} (\eng{soap} indénombrable $\rightarrow$ partitif \eng{bar of}).
\item \eng{There is \textbf{much furniture} / \textbf{a lot of furniture} in this room.} (\eng{furniture} indénombrable, pas de pluriel ; \eng{many} $\rightarrow$ \eng{much/a lot of}).
\item \eng{A bunch of flowers \textbf{is} beautiful.} (Sujet = \eng{a bunch} singulier $\rightarrow$ verbe singulier \eng{is}).
\item \eng{I need a packet of \textbf{sugar}.} (\eng{sugar} indénombrable $\rightarrow$ pas de \emph{s}).
\end{enumerate}
\end{corrige}

\courssub{Bilan et auto-évaluation}

\begin{aretenir}[Check-list de la leçon — Objectifs atteints ?]
\begin{itemize}
\item[$\square$] Je sais lister au moins 10 partitifs classés par catégorie (contenants, portions, regroupements).
\item[$\square$] J'écris systématiquement \eng{of} entre le partitif et le nom.
\item[$\square$] Je mets le pluriel sur le partitif (\eng{two boxes}) et non sur le nom indénombrable (\eng{of soap}).
\item[$\square$] J'accorde le verbe avec le partitif : \eng{How much \textbf{is} a bag?} / \eng{How much \textbf{are} two bags?}.
\item[$\square$] Je ne dis jamais \emph{a bread}, \emph{an advice}, \emph{a soap}, \emph{a furniture}, \emph{a luggage} ; j'utilise \eng{a loaf of bread}, \eng{a piece of advice}, \eng{a bar of soap}, \eng{a piece of furniture}, \eng{a piece of luggage}.
\item[$\square$] Je sais jouer un dialogue d'achat/vente complet (salutations, demandes polies, prix, monnaie, reçu).
\item[$\square$] Je sais rédiger un \eng{receipt} clair avec quantités, partitifs, prix et total.
\end{itemize}
\end{aretenir}

\begin{methode}[Méthode : Réviser efficacement les partitifs]
\begin{enumerate}
\item \textbf{Fiches « Partitif + Nom type » :} Fabrique 15 petites cartes. Recto : \eng{a loaf of} / Verso : \eng{bread}. Teste-toi dans les deux sens.
\item \textbf{Tableau des indénombrables pièges :} Affiche dans ta chambre le tableau de la section 4 (Pain, Conseil, Savon, Meuble, Bagage, Information, Éclair). Récite le partitif associé.
\item \textbf{Dictée à l'oral :} Demande à un camarade de dicter 5 phrases au singulier ; tu les écris au pluriel (ex. \eng{A bag of rice} $\rightarrow$ \eng{Three bags of rice}).
\item \textbf{Contexte réel :} La prochaine fois que tu vas au marché (Mokolo, Nkolbisson, Mfoundi...), essaie de nommer mentalement les articles en anglais avec leur partitif : \eng{a bunch of plantains}, \eng{a bag of groundnuts}, \eng{a bottle of palm wine}...
\end{enumerate}
\end{methode}

\newpage

\coursec{Méthodes et exercices résolus}

\coursec{Lesson 23 — Grammar: Partitive Expressions}

\courssub{Observation et découverte}

\begin{textebox}[Dialogue au marché Mokolo — Yaoundé]
\textbf{Vendor:} Good morning, Madam! What can I get for you today? \\
\textbf{Customer:} Good morning. I need \textbf{a bag of rice}, \textbf{two bottles of palm oil} and \textbf{a bunch of plantains}. \\
\textbf{Vendor:} Certainly. Do you want \textbf{a loaf of bread} as well? The bakery just delivered fresh ones. \\
\textbf{Customer:} Yes, please. And \textbf{a bar of soap} — the green one, \textit{Cléopâtre} if you have it. \\
\textbf{Vendor:} Here you are. That makes 12,500 francs CFA. \\
\textbf{Customer:} Here is 15,000. Keep the change. \\
\textbf{Vendor:} Thank you very much, Madam. Have a nice day!
\end{textebox}

\begin{experience}[Jeu de rôle : Au marché]
\textbf{Consigne.} En binôme, rejouez le dialogue en changeant la liste de courses. Utilisez au moins \textbf{cinq expressions partitives} différentes. L'élève A est le vendeur, l'élève B le client. Inversez les rôles.
\begin{itemize}
\item \textbf{Aide-mémoire :} \eng{a packet of sugar / flour}, \eng{a glass of water / juice}, \eng{a slice of cake / bread}, \eng{a piece of chicken / meat}, \eng{a kilo of tomatoes / onions}.
\end{itemize}
\textbf{Débriefing :} La classe note les partitifs entendus au tableau. Vérifiez l'accord du verbe (\eng{is/are}) et le pluriel du partitif seulement.
\end{experience}

\courssub{La règle : structure et formation}

\begin{propriete}[Structure de l'expression partitive]
Une expression partitive permet de \textbf{quantifier un nom indénombrable} (que l'on ne peut pas compter directement : \eng{water, rice, advice, furniture}) ou un \textbf{ensemble} (\eng{keys, flowers, plantains}).

\medskip
\textbf{Forme canonique :}
\begin{center}
\textbf{Article (a/an/one/two...) + Partitif (nom comptable) + of + Nom indénombrable / ensemble}
\end{center}

\medskip
\textbf{Points clés :}
\begin{enumerate}
\item Le \textbf{partitif} est un nom comptable (\eng{a cup, two cups, a loaf, three loaves}). Il s'accorde en nombre.
\item La préposition \textbf{of} est \textbf{invariable} et obligatoire.
\item Le \textbf{nom qui suit of} (la substance, l'objet) \textbf{ne change jamais de forme} (pas de \eng{-s}, pas de \eng{a/an} devant lui).
\item L'\textbf{accord du verbe} se fait avec le \textbf{partitif} (le sujet grammatical), pas avec le nom après \eng{of}.
\end{enumerate}
\end{propriete}

\begin{aretenir}[Règle d'or — Le nom après \eng{of} est intouchable]
\begin{itemize}
\item \eng{Two cups of tea} (\textbf{pas} \emph{two cups of teas})
\item \eng{Three pieces of advice} (\textbf{pas} \emph{three advices} ni \emph{three pieces of advices})
\item \eng{A loaf of bread} (\textbf{pas} \emph{a bread})
\end{itemize}
\emph{En français on dit « deux thés, trois conseils, un pain » ; en anglais, le nom « matière » reste nu.}
\end{aretenir}

\begin{exemplebox}[Exemples de base]
\begin{itemize}
\item \eng{A glass of water.} « Un verre d'eau. » — \eng{water} indénombrable.
\item \eng{Two slices of bread.} « Deux tranches de pain. » — \eng{slices} pluriel, \eng{bread} inchangé.
\item \eng{A bunch of keys.} « Un trousseau de clés. » — ensemble naturel.
\item \eng{Three bottles of palm wine.} « Trois bouteilles de vin de palme. » — verbe \eng{are} si sujet pluriel.
\end{itemize}
\end{exemplebox}

\courssub{Les principaux partitifs par catégories}

\begin{center}
\begin{tabularx}{\linewidth}{|X|X|X|}
\hline
\rowcolor{popblueL}
\textbf{Catégorie} & \textbf{Partitif (singulier)} & \textbf{Exemples typiques (contexte camerounais)} \\
\hline
\textbf{Contenants liquides} & a glass of, a cup of, a bottle of, a jug of, a litre of & \eng{a glass of water}, \eng{a cup of tea/coffee}, \eng{a bottle of palm oil / palm wine / soda}, \eng{a litre of milk} \\
\hline
\textbf{Contenants solides / emballages} & a bag of, a packet of, a box of, a tin of, a can of, a sachet of & \eng{a bag of rice / cement}, \eng{a packet of sugar / flour / biscuits}, \eng{a tin of sardines / milk}, \eng{a box of matches / chalk} \\
\hline
\textbf{Formes / Portions} & a piece of, a slice of, a chunk of, a lump of, a loaf of, a bar of & \eng{a piece of meat / cake / advice / information / furniture}, \eng{a slice of bread / pineapple / cake}, \eng{a loaf of bread}, \eng{a bar of soap / chocolate / gold} \\
\hline
\textbf{Groupes naturels / Grappes} & a bunch of, a cluster of, a bundle of, a heap of, a pile of & \eng{a bunch of plantains / bananas / flowers / keys}, \eng{a bundle of firewood / newspapers}, \eng{a heap of sand / rubbish} \\
\hline
\textbf{Mesures de poids / volume} & a kilo of, a gram of, a pound of, a ton of & \eng{a kilo of tomatoes / onions / meat}, \eng{500 grams of flour} \\
\hline
\end{tabularx}
\end{center}

\begin{savaistu}[Mesures et culture au Cameroun]
Au Cameroun anglophone (régions du Nord-Ouest et du Sud-Ouest), on utilise souvent le système impérial \textbf{côté marché} (\eng{a pound of tomatoes}, \eng{a pint of palm wine}) et le système métrique \textbf{côté administration/école} (\eng{a kilo of rice}, \eng{a litre of petrol}). Le \eng{pint} (environ 0,57 L) est la mesure traditionnelle pour le vin de palme. Le \eng{cup} (tasse) sert d'unité de mesure en cuisine (\eng{two cups of flour}). Connaître les deux systèmes évite les malentendus à l'achat.
\end{savaistu}

\courssub{Comparaison français / anglais}

\begin{center}
\begin{tabularx}{\linewidth}{|X|X|X|}
\hline
\rowcolor{popblueL}
\textbf{Français (structure)} & \textbf{English (structure)} & \textbf{Remarque / Piège} \\
\hline
un pain & \eng{a loaf of bread} & \eng{bread} indénombrable ; \eng{loaf} $\rightarrow$ \eng{loaves} au pluriel. \\
un conseil & \eng{a piece of advice} & \eng{advice} indénombrable (jamais \emph{an advice}). \\
une information & \eng{a piece of information} & \eng{information} indénombrable (jamais \emph{an information}). \\
une nouvelle (nouvelles) & \eng{news} (singulier !) & \eng{The news is good.} Verbe au singulier. \\
un meuble & \eng{a piece of furniture} & \eng{furniture} indénombrable (jamais \emph{a furniture}). \\
du travail / des devoirs & \eng{work} / \eng{homework} & Indénombrables. \eng{a lot of homework}. \\
une feuille (de papier) & \eng{a sheet of paper} & \eng{paper} indénombrable = matière ; \eng{a paper} = un journal / un article scientifique. \\
de la monnaie (rendu) & \eng{change} & \eng{currency} = devise officielle (FCFA, Dollar). \\
\hline
\end{tabularx}
\end{center}

\begin{attention}[Faux amis et calques fréquents]
\begin{enumerate}
\item \textbf{\eng{Advice / Information / News / Furniture / Luggage / Equipment / Jewellery / Machinery}} : \textbf{jamais de \eng{a/an}, jamais de \eng{-s}}. Pour compter : \eng{a piece of} / \eng{an item of}.
\item \textbf{\eng{Bread / Cheese / Butter / Chocolate / Cake}} : indénombrables comme matières. \eng{A loaf/slice/piece/bar of...}
\item \textbf{\eng{Paper}} : indénombrable (matière) $\rightarrow$ \eng{a sheet of paper}. Comptable = journal (\eng{a newspaper}) ou examen (\eng{an exam paper}).
\item \textbf{\eng{Hair}} : indénombrable (chevelure) $\rightarrow$ \eng{She has long hair}. Comptable = un poil (\eng{a hair} / \eng{two hairs}).
\item \textbf{\eng{Change}} (monnaie rendue) vs \textbf{\eng{Currency}} (devise). \eng{Keep the change.} / \eng{The local currency is the CFA franc.}
\end{enumerate}
\end{attention}

\courssub{Emplois en contexte et cas particuliers}

\begin{methode}[Mettre au pluriel et accorder le verbe — 4 étapes]
\begin{enumerate}
\item \cle{Identifie} le partitif (le nom comptable avant \eng{of}).
\item \cle{Pluralise} \textbf{uniquement ce partitif} : \eng{a glass} $\rightarrow$ \eng{two glasses} ; \eng{a loaf} $\rightarrow$ \eng{two loaves} ; \eng{a bunch} $\rightarrow$ \eng{two bunches}.
\item \cle{Laisse} le nom après \eng{of} \textbf{inchangé} (singulier s'il est indénombrable, pluriel s'il est un nom d'ensemble déjà pluriel comme \eng{keys}, mais sans \eng{-s} ajouté).
\item \cle{Accorde} le verbe : singulier si partitif singulier (\eng{a/one}), pluriel si partitif pluriel (\eng{two, three, several, many}).
\end{enumerate}
\end{methode}

\begin{exemplebox}[Pluriels irréguliers à connaître]
\begin{center}
\begin{tabular}{|l|l|l|}
\hline
\textbf{Singulier} & \textbf{Pluriel} & \textbf{Exemple} \\
\hline
a loaf of bread & two \textbf{loaves} of bread & \eng{Two loaves of bread are on the table.} \\
a knife (partitif rare) & two \textbf{knives} & — \\
a half & two \textbf{halves} & \eng{Two halves of an orange.} \\
\hline
\end{tabular}
\end{center}
\emph{Les autres partitifs suivent la règle régulière (\eng{-s} ou \eng{-es} après sibilante : \eng{glasses, boxes, bunches, pieces, slices, bars, packets, bottles, bags, bundles, kilos}).}
\end{exemplebox}

\begin{methode}[Choisir entre \eng{How much} et \eng{How many}]
\begin{enumerate}
\item \textbf{\eng{How much}} + \textbf{nom indénombrable seul} : \eng{How much rice did you buy?} — \eng{How much palm oil is left?}
\item \textbf{\eng{How many}} + \textbf{partitif au pluriel} : \eng{How many bags of rice did you buy?} — \eng{How many bottles of oil are there?}
\item \textbf{\eng{How many}} + \textbf{nom comptable pluriel} : \eng{How many plantains did you buy?} (ici \eng{plantains} est comptable, pas besoin de partitif).
\end{enumerate}
\end{methode}

\begin{aretenir}[Cas particulier : \eng{a piece of} — le partitif « passe-partout »]
\eng{A piece of} s'utilise avec une très grande variété de noms indénombrables : \eng{advice, information, news, furniture, luggage, equipment, cake, meat, cheese, chocolate, bread, work, homework, evidence, research}. \\
\emph{Attention :} pour le pain, \eng{a slice of} ou \eng{a loaf of} sont plus précis. Pour le chocolat, \eng{a bar of}. Pour le savon, \eng{a bar of}.
\end{aretenir}

\courssub{Exercices d'application et de transformation}

\courssub{Méthode 1 — Identifier et construire une expression partitive}

\begin{methode}[Construire une expression partitive en 4 étapes]
Pour exprimer une quantité précise d'un nom indénombrable ou d'un ensemble, suis cette démarche :
\begin{enumerate}
\item \cle{Repère} le nom que tu veux quantifier. Demande-toi : est-il \textbf{comptable} (on peut dire \eng{two apples}) ou \textbf{indénombrable} (on ne peut pas dire \eng{two milks}) ?
\item \cle{Choisis} le partitif adapté à la \textbf{forme} ou au \textbf{contenant} : \eng{a piece of} (morceau), \eng{a cup of} (tasse), \eng{a bottle of} (bouteille), \eng{a bunch of} (grappe/bouquet), \eng{a slice of} (tranche), \eng{a bar of} (tablette/barre), \eng{a packet of} (paquet), \eng{a glass of} (verre), \eng{a loaf of} (miche), \eng{a bag of} (sac), \eng{a bundle of} (fagot/botte).
\item \cle{Respecte} la structure : \textbf{article (a/an) + partitif + of + nom}. Le partitif est toujours singulier avec \eng{a} ; pour le pluriel, on met le \textbf{partitif} au pluriel, pas le nom : \eng{two cups of tea} (et jamais \emph{two teas} dans ce sens).
\item \cle{Accorde} le verbe avec le \textbf{partitif} : \eng{A bunch of keys \textbf{is} on the table.} / \eng{Two bottles of palm wine \textbf{were} sold.}
\end{enumerate}
Réflexe à retenir : en anglais, un nom indénombrable (\eng{information, bread, water, money, news, advice, furniture}) ne prend \textbf{jamais} de \eng{a/an} ni de pluriel. On le « découpe » avec un partitif.
\end{methode}

\begin{exoresolu}[Exercice type — Compléter avec le bon partitif]
\textbf{Énoncé.} Complete each sentence with the correct partitive expression: \emph{a bunch of, a cup of, a piece of, a bottle of, a slice of, a bar of, a loaf of, a glass of, a bag of, a packet of}.
\begin{enumerate}
\item Manka drinks \dots\ of coffee every morning before going to school.
\item Can you buy \dots\ of bread at the bakery in Mvog-Mbi?
\item She gave me \dots\ of advice about my studies.
\item The farmer sold \dots\ of plantains at the Mokolo market.
\item I would like \dots\ of cake, please.
\item He bought \dots\ of soap and \dots\ of groundnut oil.
\item There is \dots\ of water on the table.
\item We need \dots\ of rice for the feast (50 kg).
\item She ate \dots\ of chocolate after lunch.
\item They collected \dots\ of firewood for the night.
\end{enumerate}
\tcblower
\textbf{Solution détaillée.}
\begin{enumerate}
\item \eng{a cup of coffee} — le café se boit dans une \textbf{tasse} ; \eng{coffee} est indénombrable, on ne peut pas dire \emph{a coffee} dans ce sens classique.
\item \eng{a loaf of bread} — le pain se vend en \textbf{miche} ; \eng{bread} est indénombrable (jamais \emph{a bread}).
\item \eng{a piece of advice} — \eng{advice} est indénombrable en anglais (contrairement au français « un conseil ») ; le partitif standard est \eng{a piece of}.
\item \eng{a bunch of plantains} — les bananes plantains poussent et se vendent en \textbf{régime} : \eng{a bunch of} est le partitif des grappes et bouquets.
\item \eng{a slice of cake} — on coupe un gâteau en \textbf{tranches}.
\item \eng{a bar of soap} (tablette/barre de savon) et \eng{a bottle of groundnut oil} (l'huile se vend en \textbf{bouteille}).
\item \eng{a glass of water} — l'eau à table se sert dans un \textbf{verre}.
\item \eng{a bag of rice} — le riz s'achète en \textbf{sac} (souvent 25 ou 50 kg).
\item \eng{a bar of chocolate} — tablette de chocolat.
\item \eng{a bundle of firewood} — le bois de chauffage se vend en \textbf{fagot/botte}.
\end{enumerate}
\textbf{Conclusion.} Le choix du partitif dépend du contenant ou de la forme naturelle de la chose. Vérifier d'abord si le nom est indénombrable : si oui, le partitif est \textbf{obligatoire} pour compter.
\end{exoresolu}

\courssub{Méthode 2 — Éviter les calques du français}

\begin{methode}[Traduire sans calquer : le réflexe « indénombrable »]
\begin{enumerate}
\item \cle{Repère} dans la phrase française un nom comme « conseil, information, nouvelle, meuble, pain, argent, travail, devoir, bagage ».
\item \cle{Vérifie} dans la liste des indénombrables anglais : \eng{advice, information, news, furniture, bread, money, work, homework, luggage, progress, knowledge, equipment, jewellery}.
\item \cle{Supprime} tout article \eng{a/an} devant ces noms et tout \eng{-s} de pluriel : « des informations » se dit \eng{information} ou \eng{some information}, jamais \emph{informations}.
\item \cle{Ajoute} un partitif si tu veux compter : « deux conseils » = \eng{two pieces of advice} ; « trois nouvelles » = \eng{three pieces of news}.
\item \cle{Accorde} le verbe au singulier avec ces noms : \eng{The news \textbf{is} good.} (et non \emph{are}).
\end{enumerate}
\end{methode}

\begin{exoresolu}[Exercice type — Traduire en anglais]
\textbf{Énoncé.} Translate into English.
\begin{enumerate}
\item Mon oncle m'a donné un bon conseil.
\item J'ai deux informations importantes pour toi.
\item Les nouvelles sont bonnes aujourd'hui.
\item Elle a acheté un pain et une bouteille d'huile de palme.
\item Nous avons beaucoup de devoirs ce soir.
\item Il a cassé un meuble en déplaçant la table.
\item Donne-moi une feuille de papier, s'il te plaît.
\item Le chauffeur m'a rendu ma monnaie.
\end{enumerate}
\tcblower
\textbf{Solution détaillée.}
\begin{enumerate}
\item \eng{My uncle gave me a good piece of advice.} — \eng{advice} indénombrable : impossible de dire \emph{an advice}. Le partitif \eng{a piece of} permet de compter.
\item \eng{I have two important pieces of information for you.} — \eng{information} indénombrable (pas de \emph{informations}) ; au pluriel, c'est le partitif qui prend le \eng{-s} : \eng{pieces}.
\item \eng{The news is good today.} — \eng{news} finit en \eng{-s} mais reste \textbf{singulier} : verbe \eng{is}.
\item \eng{She bought a loaf of bread and a bottle of palm oil.} — \eng{bread} indénombrable $\rightarrow$ \eng{a loaf of bread} ; l'huile se vend en bouteille $\rightarrow$ \eng{a bottle of palm oil}.
\item \eng{We have a lot of homework tonight.} — \eng{homework} indénombrable : jamais \emph{homeworks} ; on utilise \eng{a lot of} + singulier.
\item \eng{He broke a piece of furniture moving the table.} — \eng{furniture} indénombrable $\rightarrow$ \eng{a piece of furniture}.
\item \eng{Give me a sheet of paper, please.} — \eng{paper} (matière) indénombrable $\rightarrow$ \eng{a sheet of}. \eng{A paper} = un journal / un article.
\item \eng{The driver gave me my change.} — \eng{change} = monnaie (rendu) ; \eng{currency} = devise.
\end{enumerate}
\textbf{Conclusion.} La traduction correcte exige deux réflexes : reconnaître l'indénombrable anglais, puis choisir le partitif ou le quantifieur adapté.
\end{exoresolu}

\courssub{Méthode 3 — Mettre au pluriel et accorder le verbe}

\begin{methode}[Pluriel du partitif et accord du verbe]
\begin{enumerate}
\item \cle{Pluralise} uniquement le \textbf{partitif}, jamais le nom indénombrable : \eng{a cup of tea} $\rightarrow$ \eng{two cup\textbf{s} of tea} ; \eng{a piece of chalk} $\rightarrow$ \eng{three piece\textbf{s} of chalk}.
\item \cle{Accorde} le verbe avec le partitif : singulier si \eng{a/one}, pluriel si \eng{two, three, several} : \eng{A bottle of water \textbf{is} enough.} / \eng{Three bottles of water \textbf{are} enough.}
\item \cle{Attention} aux partitifs irréguliers au pluriel : \eng{a loaf of bread} $\rightarrow$ \eng{two loa\textbf{ves} of bread} ; \eng{a half} $\rightarrow$ \eng{two halves}.
\item \cle{Choisis} entre \eng{much/many} : avec le partitif comptable, on utilise \eng{many} : \eng{How many bottles of oil did you buy?} ; avec le nom indénombrable seul, \eng{much} : \eng{How much oil did you buy?}
\end{enumerate}
\end{methode}

\begin{exoresolu}[Exercice type — Mettre au pluriel]
\textbf{Énoncé.} Rewrite the sentences in the plural (change \eng{a/one} to \eng{two/three/several} and adapt the verb).
\begin{enumerate}
\item I bought a loaf of bread.
\item She drank a glass of water.
\item There is a bottle of palm wine in the kitchen.
\item He gave me a piece of information.
\item A bunch of keys was found near the school gate.
\item A bar of soap costs 300 francs.
\item She needs a packet of flour.
\end{enumerate}
\tcblower
\textbf{Solution détaillée.}
\begin{enumerate}
\item \eng{I bought two loaves of bread.} — pluriel irrégulier : \eng{loaf} $\rightarrow$ \eng{loaves} ; \eng{bread} reste inchangé.
\item \eng{She drank three glasses of water.} — \eng{glass} prend \eng{-es} (terminaison en \eng{-ss}) ; \eng{water} inchangé.
\item \eng{There are four bottles of palm wine in the kitchen.} — le verbe passe au pluriel \eng{are} car le partitif \eng{bottles} est pluriel.
\item \eng{He gave me several pieces of information.} — seul \eng{piece} se pluralise ; \eng{information} ne prend jamais de \eng{-s}.
\item \eng{Two bunches of keys were found near the school gate.} — \eng{bunch} $\rightarrow$ \eng{bunches} ; verbe au pluriel \eng{were} (accord avec le partitif).
\item \eng{Two bars of soap cost 300 francs each.} — \eng{bar} $\rightarrow$ \eng{bars} ; verbe \eng{cost} (pluriel). \emph{Note :} \eng{each} à la fin précise le prix unitaire.
\item \eng{She needs three packets of flour.} — \eng{packet} $\rightarrow$ \eng{packets} ; \eng{flour} inchangé.
\end{enumerate}
\textbf{Conclusion.} Au pluriel, tout change (article, partitif, verbe) \textbf{sauf} le nom indénombrable, qui reste invariable.
\end{exoresolu}

\courssub{Méthode 4 — Utiliser les partitifs dans un court texte}

\begin{methode}[Rédiger un paragraphe avec des partitifs]
Pour une composition (recette, liste de courses, récit de marché) :
\begin{enumerate}
\item \cle{Dresse} mentalement ta liste : quels noms sont indénombrables (\eng{rice, oil, flour, sugar, water, soap, bread}) ?
\item \cle{Associe} à chacun un partitif naturel : \eng{a bag of rice, a bottle of oil, a packet of flour, a cup of sugar, a glass of water, a bar of soap, a loaf of bread}.
\item \cle{Varie} les quantifieurs : \eng{a, two, several, a few} + partitif ; \eng{some, a lot of, much} + nom indénombrable.
\item \cle{Relis} en vérifiant trois points : pas de \eng{a/an} devant un indénombrable, pas de \eng{-s} sur un indénombrable, verbe accordé avec le partitif.
\end{enumerate}
\end{methode}

\begin{exoresolu}[Exercice type — Composition guidée]
\textbf{Énoncé.} Your mother sends you to the Mokolo market. Write a short paragraph (6–7 sentences) describing what you bought, using at least six different partitive expressions.
\tcblower
\textbf{Solution détaillée.}

\textbf{Raisonnement.} On choisit des produits du marché camerounais : riz (\eng{a bag of}), huile de palme (\eng{a bottle of}), savon (\eng{a bar of}), bananes plantains (\eng{a bunch of}), pain (\eng{a loaf of}), sucre (\eng{a packet of}), tomates (\eng{a pile of} ou \eng{a bowl of} — comptable mais souvent vendu en tas/bol), eau (\eng{a bottle of} ou \eng{a jerrycan of}). On accorde les verbes et on évite tout pluriel sur les indénombrables.

\textbf{Réponse rédigée.}
\eng{Last Saturday, my mother sent me to the Mokolo market with a long shopping list. First, I bought a large bag of rice and two bottles of red palm oil from a wholesaler near the main entrance. Then I found a beautiful bunch of ripe plantains at a very good price from a young lady selling under an umbrella. I also bought a bar of "Cléopâtre" soap and a packet of refined sugar for the house. Before leaving, I took a fresh loaf of bread from the bakery at the market exit. Finally, I carried a heavy jerrycan of water because the tap at home was dry. When I arrived home, my mother checked everything and smiled because the list was complete.}

\textbf{Vérification.} Six partitifs différents utilisés (\eng{a bag of, two bottles of, a bunch of, a bar of, a packet of, a loaf of, a jerrycan of}) ; aucun indénombrable ne porte de \eng{-s} (\eng{rice, oil, sugar, bread, water}) ; les verbes sont correctement accordés (\eng{I bought, I found, I took, my mother checked, was}).
\end{exoresolu}

\courssub{Méthode 5 — Transformer et reformuler}

\begin{methode}[Transformer une phrase avec un partitif]
Aux exercices de transformation type examen :
\begin{enumerate}
\item \cle{Identifie} la consigne : passer d'un nom seul à un partitif (\eng{water} $\rightarrow$ \eng{a glass of water}), ou l'inverse.
\item \cle{Conserve} le sens : le partitif précise la quantité sans changer la nature de la chose.
\item \cle{Adapte} l'article et le verbe : \eng{Water is on the table.} $\rightarrow$ \eng{A glass of water is on the table.} (le verbe reste singulier car \eng{glass} est singulier).
\item \cle{Pour une question}, utilise \eng{How many} + partitif pluriel ou \eng{How much} + indénombrable : \eng{How many cups of tea do you drink?} / \eng{How much tea do you drink?}
\end{enumerate}
\end{methode}

\begin{exoresolu}[Exercice type — Transformation]
\textbf{Énoncé.} Rewrite each sentence using a partitive expression, keeping the same meaning.
\begin{enumerate}
\item I drank some water. (one glass)
\item She bought some chocolate. (one bar)
\item We need some flour. (two packets)
\item He eats some bread every morning. (one slice)
\item They sold some firewood. (one bundle)
\item The teacher gave us some advice. (two pieces)
\item There is some furniture in the room. (three pieces)
\end{enumerate}
\tcblower
\textbf{Solution détaillée.}
\begin{enumerate}
\item \eng{I drank a glass of water.} — on remplace \eng{some} par \eng{a glass of} ; \eng{water} reste invariable.
\item \eng{She bought a bar of chocolate.} — le chocolat en tablette se dit \eng{a bar of chocolate}.
\item \eng{We need two packets of flour.} — quantité plurielle : \eng{packets} prend le \eng{-s}, pas \eng{flour}.
\item \eng{He eats a slice of bread every morning.} — une tranche de pain : \eng{a slice of bread} ; jamais \emph{a bread}.
\item \eng{They sold a bundle of firewood.} — le bois de chauffage se vend en fagot : \eng{a bundle of firewood} ; \eng{firewood} est indénombrable.
\item \eng{The teacher gave us two pieces of advice.} — \eng{advice} indénombrable $\rightarrow$ \eng{pieces of advice} (pluriel sur \eng{pieces}).
\item \eng{There are three pieces of furniture in the room.} — \eng{furniture} indénombrable $\rightarrow$ \eng{pieces of furniture} ; verbe \eng{are} (partitif pluriel).
\end{enumerate}
\textbf{Conclusion.} La transformation exige de garder le nom indénombrable intact et de reporter toute la marque du nombre sur le partitif.
\end{exoresolu}

\begin{exercice}[Exercices d'entraînement — À faire seul(e)]
\begin{enumerate}
\item \textbf{Complétez} avec le partitif approprié : \eng{a bunch of, a cup of, a piece of, a bottle of, a slice of, a bar of, a loaf of, a glass of, a bag of, a packet of, a sheet of, a jar of}.
\begin{enumerate}
\item Could you pass me \dots\ of jam for my toast?
\item The secretary typed \dots\ of paper for the report.
\item We need \dots\ of cement to finish the wall.
\item He drank \dots\ of fresh juice after the match.
\item She gave me \dots\ of useful information.
\end{enumerate}
\item \textbf{Mettez au pluriel} (verbe inclus) :
\begin{enumerate}
\item A loaf of bread is on the table.
\item There is a bottle of soda in the fridge.
\item A piece of chalk is broken.
\item A bunch of flowers was offered to the teacher.
\end{enumerate}
\item \textbf{Traduisez} en anglais :
\begin{enumerate}
\item J'ai acheté trois pains hier.
\item Il n'y a plus d'eau dans la bouteille.
\item Donne-moi un conseil, s'il te plaît.
\item Ces meubles sont très chers.
\item Combien de sacs de riz voulez-vous ?
\end{enumerate}
\item \textbf{Transformation} : Remplacez la quantité vague par le partitif indiqué.
\begin{enumerate}
\item She eats bread. (a slice of)
\item They need money. (a sum of)
\item We bought rice. (two bags of)
\item He gave me advice. (a piece of)
\end{enumerate}
\end{enumerate}
\end{exercice}

\begin{corrige}[Exercices d'entraînement]
\textbf{1. Complétion :}
\begin{enumerate}
\item \eng{a jar of jam} (pot de confiture).
\item \eng{a sheet of paper} (feuille de papier — \eng{paper} matière indénombrable).
\item \eng{a bag of cement} (sac de ciment).
\item \eng{a glass of fresh juice} (verre de jus).
\item \eng{a piece of information} (information indénombrable).
\end{enumerate}
\textbf{2. Pluriel :}
\begin{enumerate}
\item \eng{Two loaves of bread are on the table.} (\eng{loaf} $\rightarrow$ \eng{loaves}, verbe \eng{are}).
\item \eng{There are three bottles of soda in the fridge.} (\eng{bottles}, verbe \eng{are}).
\item \eng{Three pieces of chalk are broken.} (\eng{pieces}, verbe \eng{are}).
\item \eng{Two bunches of flowers were offered to the teacher.} (\eng{bunches}, verbe \eng{were}).
\end{enumerate}
\textbf{3. Traduction :}
\begin{enumerate}
\item \eng{I bought three loaves of bread yesterday.}
\item \eng{There is no more water in the bottle.} (ou \eng{The bottle is empty.})
\item \eng{Give me a piece of advice, please.}
\item \eng{This furniture is very expensive.} / \eng{These pieces of furniture are very expensive.}
\item \eng{How many bags of rice do you want?}
\end{enumerate}
\textbf{4. Transformation :}
\begin{enumerate}
\item \eng{She eats a slice of bread.}
\item \eng{They need a sum of money.}
\item \eng{We bought two bags of rice.}
\item \eng{He gave me a piece of advice.}
\end{enumerate}
\end{corrige}

\begin{attention}[Les 5 erreurs qui coûtent des points au Bac]
\begin{enumerate}
\item \textbf{Pluriel sur un indénombrable.} Écrire \emph{informations, advices, homeworks, breads, furnitures, luggages, equipments} est la faute la plus sanctionnée. Correct : \eng{information, advice, homework, bread, furniture, luggage, equipment}. Pour compter : \eng{two pieces of information}.
\item \textbf{Article devant un indénombrable.} Ne dis jamais \emph{a bread, an advice, a money, a furniture, a luggage}. Correct : \eng{a loaf of bread, a piece of advice, a sum of money, a piece of furniture, a piece of luggage}.
\item \textbf{Pluriel mis sur le nom au lieu du partitif.} \emph{Two cups of teas} est faux : c'est le partitif qui se pluralise : \eng{two cups of tea}. De même \eng{two loaves of bread} (attention au pluriel irrégulier \eng{loaves}), \eng{two pieces of advice} (pas \emph{advices}).
\item \textbf{Verbe mal accordé.} \eng{The news \textbf{is} good} (singulier malgré le \eng{-s}) ; \eng{A bunch of keys \textbf{is} missing} (accord avec \eng{bunch}, singulier) ; \eng{Two bottles of oil \textbf{are} on the table} (accord avec \eng{bottles}, pluriel).
\item \textbf{Faux amis et calques du français.} « Une information » n'est pas \emph{an information} mais \eng{a piece of information} ; « un meuble » n'est pas \emph{a furniture} mais \eng{a piece of furniture} ; « de la monnaie » (rendu) se dit \eng{change}, pas \emph{currency} dans la vie quotidienne. Méfie-toi aussi de \eng{a sheet of paper} (feuille) et non \emph{a paper} pour une simple feuille ; \eng{a paper} = un journal / un devoir / un article scientifique.
\end{enumerate}
\end{attention}

\newpage

\coursec{Exercices}

\courssub{Je vérifie mes connaissances}

\begin{exercice}[QCM : choisir la bonne option]
Pour chaque phrase, choisis la bonne réponse (a, b ou c).
\begin{enumerate}
\item She gave me \textbf{...} before the exam. \quad a) an advice \quad b) a piece of advice \quad c) advices
\item I need \textbf{...} about the cocoa market. \quad a) an information \quad b) a piece of information \quad c) informations
\item He bought \textbf{...} at the bakery. \quad a) a bread \quad b) a loaf of bread \quad c) breads
\item There is \textbf{...} on the table. \quad a) a water \quad b) a bottle of water \quad c) waters
\item My father gave me \textbf{...} for my birthday. \quad a) a money \quad b) a sum of money \quad c) moneys
\item She drank \textbf{...} this morning. \quad a) a tea \quad b) a cup of tea \quad c) teas
\end{enumerate}
\end{exercice}

\begin{exercice}[Vrai ou faux ? Justifie ta réponse]
Dis si chaque phrase est correcte (True) ou incorrecte (False), puis justifie en une phrase en français.
\begin{enumerate}
\item \eng{I bought two rices at the market.}
\item \eng{She cut a slice of bread for me.}
\item \eng{We saw many furnitures in the shop.}
\item \eng{He gave me a piece of good news.}
\item \eng{I need a pair of scissors to cut this paper.}
\item \eng{The teacher gave us three homeworks yesterday.}
\end{enumerate}
\end{exercice}

\begin{exercice}[Vocabulaire : quel partitif ?]
Complète chaque phrase avec le partitif qui convient, choisi dans la liste : \cle{a bunch of}, \cle{a tin of}, \cle{a bar of}, \cle{a sheet of}, \cle{a glass of}, \cle{a packet of}.
\begin{enumerate}
\item My grandmother keeps \textbf{...} sardines in the kitchen cupboard.
\item Please give me \textbf{...} paper; I want to write a letter.
\item She bought \textbf{...} soap to wash the children's clothes.
\item He offered his teacher \textbf{...} flowers at the end of the year.
\item After the football match, the players drank \textbf{...} cold water.
\item My little brother opened \textbf{...} biscuits before dinner.
\end{enumerate}
\end{exercice}

\begin{exercice}[Complète avec le bon partitif]
Complète chaque phrase avec l'expression partitive qui convient. Attention à l'article (\eng{a} ou \eng{an}).
\begin{enumerate}
\item Could you lend me \textbf{...} (une feuille de) paper, please?
\item My uncle grows \textbf{...} (un sac de) cocoa beans on his farm.
\item The journalist gave us \textbf{...} (un renseignement) about the new market.
\item She poured \textbf{...} (une tasse de) coffee for the visitor.
\item We heard \textbf{...} (un conseil) from the headmaster this morning.
\item The trader sells \textbf{...} (une grappe de) bananas every Saturday.
\end{enumerate}
\end{exercice}

\courssub{Je m'entraîne}

\begin{exercice}[Traduction : du français vers l'anglais]
Traduis les phrases suivantes en anglais correct.
\begin{enumerate}
\item Ma mère a acheté un sac de riz, deux grappes de bananes plantain, une bouteille d'huile de palme et une savonnette au marché.
\item Le fermier a récolté trois sacs de maïs cette année.
\item J'ai besoin d'un verre d'eau et d'un morceau de pain.
\item Elle m'a donné un bon conseil avant l'examen.
\item Nous avons reçu deux informations importantes ce matin.
\end{enumerate}
\end{exercice}

\begin{exercice}[Transformation : mets au pluriel]
Réécris chaque phrase en mettant l'expression partitive au pluriel. Fais tous les changements nécessaires (verbe, déterminants).

Exemple : \eng{There is a cup of tea on the table.} $\rightarrow$ \eng{There are two cups of tea on the table.}
\begin{enumerate}
\item \eng{There is a cup of tea on the tray.}
\item \eng{She bought a loaf of bread at the bakery.}
\item \eng{The teacher gave us a piece of information.}
\item \eng{He drank a bottle of water after the race.}
\item \eng{I need a sheet of paper for my drawing.}
\end{enumerate}
\end{exercice}

\begin{exercice}[Texte à trous : au marché de Mokolo]
Complète le texte suivant avec les partitifs de la liste (chaque partitif n'est utilisé qu'une seule fois) : \cle{a bag of}, \cle{a bunch of}, \cle{a piece of}, \cle{a bottle of}, \cle{a tin of}, \cle{a loaf of}, \cle{a bar of}, \cle{a pair of}.

\begin{textebox}[At Mokolo Market]
Last Saturday, Mama went to Mokolo market in Yaoundé very early. First, she bought (1) \textbf{...} rice from a trader near the entrance. Then she chose (2) \textbf{...} ripe plantains for the evening meal. At the next stall, she paid for (3) \textbf{...} palm oil and (4) \textbf{...} tomatoes. Before leaving, she also bought (5) \textbf{...} bread for breakfast, (6) \textbf{...} soap for the laundry, and (7) \textbf{...} new shoes for my little sister. When she came home, she gave me (8) \textbf{...} good advice: “Always compare prices before you buy anything!”
\end{textebox}
\end{exercice}

\begin{exercice}[Appariement : l'objet et son partitif]
Associe chaque élément de la colonne A au partitif qui convient dans la colonne B. Attention : un partitif peut convenir à plusieurs éléments.

\begin{center}
\begin{tabularx}{\linewidth}{|X|X|}
\hline
\rowcolor{popblueL} \textbf{Colonne A} & \textbf{Colonne B} \\
\hline
1. bread & a) a bunch of \\
\hline
2. keys & b) a loaf of \\
\hline
3. milk & c) a glass of \\
\hline
4. chocolate & d) a bar of \\
\hline
5. paper & e) a sheet of \\
\hline
6. bananas & f) a bottle of \\
\hline
7. news & g) a piece of \\
\hline
8. scissors & h) a pair of \\
\hline
\end{tabularx}
\end{center}
\end{exercice}

\courssub{Je me prépare au Bac}

\begin{exercice}[Type Bac — Compréhension et langue (10 points)]
Lis le texte suivant, puis réponds aux questions.

\begin{textebox}[Economic Life in a Cameroonian Village]
In the village of Obala, not far from Yaoundé, economic life begins before sunrise. Every morning, traders carry a bag of maize, a bunch of plantains or a basket of tomatoes to the small market in the centre of the village. Mama Béa, a well-known trader, sells a cup of rice for a few hundred francs and a bottle of groundnut oil for a little more. Her customers often ask for a piece of advice about cooking, and she always answers with a smile.

Not far from her stall, a carpenter repairs a pair of broken chairs, while a young boy sells a loaf of bread and a tin of sardines to passers-by. At the end of the day, each trader counts a sum of money earned and buys a bar of soap or a packet of salt for the family. In Obala, a piece of good news travels fast: last week, the whole village celebrated when the cooperative received a new lorry to transport the cocoa harvest to Douala.
\end{textebox}

\textbf{A. Comprehension (4 points).} Answer in complete sentences in English.
\begin{enumerate}
\item Where does the economic life of the village take place? (1 pt)
\item Name two products that Mama Béa sells. (1 pt)
\item What does each trader do at the end of the day? (1 pt)
\item Why did the village celebrate last week? (1 pt)
\end{enumerate}

\textbf{B. Language (6 points).}
\begin{enumerate}
\item Find in the text six different partitive expressions and copy them. (3 pts)
\item Rewrite in the plural: \eng{She sells a cup of rice.} / \eng{He buys a loaf of bread.} (2 pts)
\item Correct the mistake: \eng{Mama Béa gave me an advice about cooking.} (1 pt)
\end{enumerate}
\end{exercice}

\begin{exercice}[Type Bac — Traduction et correction d'erreurs (8 points)]
\textbf{A. Traduis en anglais (4 points).}
\begin{enumerate}
\item Mon père a acheté un sac de ciment et une feuille de tôle pour réparer la maison.
\item La vendeuse m'a donné un morceau de savon et une bouteille d'huile.
\end{enumerate}

\textbf{B. Corrige la faute dans chaque phrase (4 points).} Chaque phrase contient une erreur typique des francophones ; réécris la phrase correctement.
\begin{enumerate}
\item \eng{The teacher gave us an homework for the holidays.}
\item \eng{I bought three breads at the bakery this morning.}
\item \eng{She told me a news about the new market in Douala.}
\item \eng{There are many furnitures in the headmaster's office.}
\end{enumerate}
\end{exercice}

\begin{exercice}[Type Bac — Rédaction guidée (8 points)]
Write a short paragraph (60--80 words) entitled \textbf{“My mother's shopping day”}. Describe what your mother bought at the market. You must use \textbf{at least six different partitive expressions} (for example: \cle{a bag of}, \cle{a bunch of}, \cle{a bottle of}, \cle{a piece of}, \cle{a bar of}, \cle{a loaf of}, \cle{a tin of}, \cle{a pair of}). Pay attention to countable and uncountable nouns, and to the verb agreement.
\end{exercice}

\newpage

\coursec{Corrigés des exercices}

\begin{corrige}[Exercice 1 — QCM : choisir la bonne option]
\begin{enumerate}
\item \textbf{b) a piece of advice} — \eng{She gave me a piece of advice before the exam.} « Elle m'a donné un conseil avant l'examen. » \emph{Advice} est indénombrable : jamais d'article \eng{an} ni de pluriel \eng{advices}.
\item \textbf{b) a piece of information} — \eng{I need a piece of information about the cocoa market.} \emph{Information} est indénombrable : \eng{an information} et \eng{informations} sont impossibles.
\item \textbf{b) a loaf of bread} — \eng{He bought a loaf of bread at the bakery.} « Il a acheté une miche de pain à la boulangerie. » \emph{Bread} est indénombrable.
\item \textbf{b) a bottle of water} — \eng{There is a bottle of water on the table.} \emph{Water} est indénombrable ; \eng{a water} signifierait « une eau minérale » (sens différent).
\item \textbf{b) a sum of money} — \eng{My father gave me a sum of money for my birthday.} \emph{Money} est indénombrable : jamais \eng{a money} ni \eng{moneys}.
\item \textbf{b) a cup of tea} — \eng{She drank a cup of tea this morning.} « Elle a bu une tasse de thé ce matin. »
\end{enumerate}
\textbf{Règle générale :} les noms indénombrables (\emph{advice, information, bread, water, money, tea}) ne prennent ni \eng{a/an} ni pluriel ; on les compte avec un partitif : \cle{a piece of}, \cle{a loaf of}, \cle{a bottle of}, \cle{a cup of}, \cle{a sum of} + nom indénombrable.
\end{corrige}

\begin{corrige}[Exercice 2 — Vrai ou faux ?]
\begin{enumerate}
\item \textbf{False.} \emph{Rice} est indénombrable : on ne peut pas dire \eng{two rices}. Correction : \eng{I bought two bags of rice at the market.} « J'ai acheté deux sacs de riz au marché. »
\item \textbf{True.} \eng{She cut a slice of bread for me.} « Elle m'a coupé une tranche de pain. » \emph{Bread} est indénombrable, donc le partitif \cle{a slice of} est obligatoire.
\item \textbf{False.} \emph{Furniture} est indénombrable : pas de pluriel en \emph{-s}. Correction : \eng{We saw a lot of furniture in the shop.} ou \eng{many pieces of furniture}.
\item \textbf{True.} \eng{He gave me a piece of good news.} « Il m'a annoncé une bonne nouvelle. » \emph{News} est indénombrable malgré son \emph{-s} final.
\item \textbf{True.} \eng{I need a pair of scissors to cut this paper.} « J'ai besoin d'une paire de ciseaux. » \emph{Scissors} est un nom toujours pluriel ; on le compte avec \cle{a pair of}.
\item \textbf{False.} \emph{Homework} est indénombrable. Correction : \eng{The teacher gave us three pieces of homework yesterday.} ou simplement \eng{a lot of homework}.
\end{enumerate}
\end{corrige}

\begin{corrige}[Exercice 3 — Vocabulaire : quel partitif ?]
\begin{enumerate}
\item \eng{My grandmother keeps \textbf{a tin of} sardines in the kitchen cupboard.} « Ma grand-mère garde une boîte de sardines dans le placard. »
\item \eng{Please give me \textbf{a sheet of} paper; I want to write a letter.} « Donne-moi une feuille de papier, s'il te plaît. »
\item \eng{She bought \textbf{a bar of} soap to wash the children's clothes.} « Elle a acheté une savonnette. »
\item \eng{He offered his teacher \textbf{a bunch of} flowers at the end of the year.} « Il a offert un bouquet de fleurs à son professeur. »
\item \eng{After the football match, the players drank \textbf{a glass of} cold water.} « Après le match, les joueurs ont bu un verre d'eau froide. »
\item \eng{My little brother opened \textbf{a packet of} biscuits before dinner.} « Mon petit frère a ouvert un paquet de biscuits. »
\end{enumerate}
\textbf{À retenir :} \cle{a tin of} (boîte de conserve), \cle{a sheet of} (feuille de papier), \cle{a bar of} (savon, chocolat), \cle{a bunch of} (fleurs, clés, bananes), \cle{a glass of} (boisson dans un verre), \cle{a packet of} (paquet).
\end{corrige}

\begin{corrige}[Exercice 4 — Complète avec le bon partitif]
\begin{enumerate}
\item \eng{Could you lend me \textbf{a sheet of} paper, please?} « Pourrais-tu me prêter une feuille de papier ? »
\item \eng{My uncle grows \textbf{a bag of} cocoa beans on his farm.} « Mon oncle cultive un sac de fèves de cacao dans sa ferme. » (On accepte aussi \eng{a sack of}.)
\item \eng{The journalist gave us \textbf{a piece of information} about the new market.} « Le journaliste nous a donné un renseignement. »
\item \eng{She poured \textbf{a cup of} coffee for the visitor.} « Elle a versé une tasse de café pour le visiteur. »
\item \eng{We heard \textbf{a piece of advice} from the headmaster this morning.} « Nous avons entendu un conseil du proviseur ce matin. »
\item \eng{The trader sells \textbf{a bunch of} bananas every Saturday.} « Le vendeur vend une régime/grappe de bananes chaque samedi. »
\end{enumerate}
\textbf{Point de vigilance :} aucun de ces partitifs ne commence par une voyelle, donc l'article est toujours \eng{a} (jamais \eng{an}) : \eng{a sheet}, \eng{a bag}, \eng{a piece}, \eng{a cup}, \eng{a bunch}.
\end{corrige}

\begin{corrige}[Exercice 5 — Traduction : du français vers l'anglais]
\begin{enumerate}
\item « Ma mère a acheté un sac de riz, deux grappes de bananes plantain, une bouteille d'huile de palme et une savonnette au marché. » \\
\eng{My mother bought a bag of rice, two bunches of plantains, a bottle of palm oil and a bar of soap at the market.} \\
Attention : au pluriel, c'est le partitif qui prend le \emph{-s} : \eng{two bunch\textbf{es} of plantains} ; \emph{rice}, \emph{oil} et \emph{soap} restent sans \emph{-s}.
\item « Le fermier a récolté trois sacs de maïs cette année. » \\
\eng{The farmer harvested three bags of maize this year.} (\emph{maize} est indénombrable.)
\item « J'ai besoin d'un verre d'eau et d'un morceau de pain. » \\
\eng{I need a glass of water and a piece of bread.} (On accepte \eng{a slice of bread} pour « une tranche de pain ».)
\item « Elle m'a donné un bon conseil avant l'examen. » \\
\eng{She gave me a piece of good advice before the exam.} — Jamais \eng{a good advice}.
\item « Nous avons reçu deux informations importantes ce matin. » \\
\eng{We received two pieces of important information this morning.} — Jamais \eng{two informations}.
\end{enumerate}
\textbf{Méthode :} repérer le nom indénombrable français (riz, pain, conseil, information), choisir le partitif anglais adapté, mettre le pluriel sur le partitif et non sur le nom indénombrable.
\end{corrige}

\begin{corrige}[Exercice 6 — Transformation : mets au pluriel]
\begin{enumerate}
\item \eng{There is a cup of tea on the tray.} $\rightarrow$ \eng{There \textbf{are} two \textbf{cups} of tea on the tray.} « Il y a deux tasses de thé sur le plateau. »
\item \eng{She bought a loaf of bread at the bakery.} $\rightarrow$ \eng{She bought two \textbf{loaves} of bread at the bakery.} Attention au pluriel irrégulier : \emph{loaf} $\rightarrow$ \emph{loaves}.
\item \eng{The teacher gave us a piece of information.} $\rightarrow$ \eng{The teacher gave us two \textbf{pieces} of information.}
\item \eng{He drank a bottle of water after the race.} $\rightarrow$ \eng{He drank two \textbf{bottles} of water after the race.}
\item \eng{I need a sheet of paper for my drawing.} $\rightarrow$ \eng{I need two \textbf{sheets} of paper for my drawing.}
\end{enumerate}
\textbf{Règle :} au pluriel, seul le partitif change (\eng{cups, loaves, pieces, bottles, sheets}) ; le nom indénombrable reste invariable (\eng{tea, bread, information, water, paper}). Avec \eng{there}, le verbe s'accorde avec le partitif : \eng{There \textbf{is} a cup} / \eng{There \textbf{are} two cups}.
\end{corrige}

\begin{corrige}[Exercice 7 — Texte à trous : au marché de Mokolo]
\begin{enumerate}
\item \textbf{a bag of} rice — « un sac de riz »
\item \textbf{a bunch of} ripe plantains — « une grappe de plantains mûrs »
\item \textbf{a bottle of} palm oil — « une bouteille d'huile de palme »
\item \textbf{a tin of} tomatoes — « une boîte de tomates »
\item \textbf{a loaf of} bread — « une miche de pain »
\item \textbf{a bar of} soap — « une savonnette »
\item \textbf{a pair of} new shoes — « une paire de chaussures neuves »
\item \textbf{a piece of} good advice — « un bon conseil »
\end{enumerate}
Texte complet vérifié : \eng{Last Saturday, Mama went to Mokolo market in Yaoundé very early. First, she bought a bag of rice from a trader near the entrance. Then she chose a bunch of ripe plantains for the evening meal. At the next stall, she paid for a bottle of palm oil and a tin of tomatoes. Before leaving, she also bought a loaf of bread for breakfast, a bar of soap for the laundry, and a pair of new shoes for my little sister. When she came home, she gave me a piece of good advice: “Always compare prices before you buy anything!”}

\textbf{Astuce de résolution :} commencer par les associations les plus évidentes (\eng{a pair of shoes}, \eng{a bar of soap}, \eng{a piece of advice}), puis éliminer les partitifs restants un à un.
\end{corrige}

\begin{corrige}[Exercice 8 — Appariement : l'objet et son partitif]
\begin{center}
\begin{tabularx}{\linewidth}{|X|X|X|}
\hline
\rowcolor{popblueL} \textbf{Colonne A} & \textbf{Réponse} & \textbf{Traduction} \\
\hline
1. bread & b) a loaf of bread & une miche de pain \\
\hline
2. keys & a) a bunch of keys & un trousseau de clés \\
\hline
3. milk & c) a glass of milk (ou f) a bottle of milk) & un verre / une bouteille de lait \\
\hline
4. chocolate & d) a bar of chocolate & une tablette de chocolat \\
\hline
5. paper & e) a sheet of paper & une feuille de papier \\
\hline
6. bananas & a) a bunch of bananas & une régime de bananes \\
\hline
7. news & g) a piece of news & une nouvelle \\
\hline
8. scissors & h) a pair of scissors & une paire de ciseaux \\
\hline
\end{tabularx}
\end{center}
\textbf{Remarques :} \emph{milk} accepte deux réponses (\cle{a glass of} ou \cle{a bottle of}) selon le contenant ; l'énoncé précisait qu'un partitif pouvait convenir à plusieurs éléments. \emph{News} est indénombrable malgré son \emph{-s} : \eng{a piece of news}, jamais \eng{a new}.
\end{corrige}

\begin{corrige}[Exercice 9 — Type Bac : Compréhension et langue (10 points)]
\textbf{A. Comprehension (4 points — 1 pt par réponse correcte, en phrase complète).}
\begin{enumerate}
\item \eng{The economic life of the village takes place at the small market in the centre of the village.} « La vie économique du village se déroule au petit marché du centre. »
\item \eng{Mama Béa sells rice (by the cup) and groundnut oil (by the bottle).} « Elle vend du riz et de l'huile d'arachide. »
\item \eng{At the end of the day, each trader counts the sum of money earned and buys something for the family, like a bar of soap or a packet of salt.}
\item \eng{The village celebrated because the cooperative received a new lorry to transport the cocoa harvest to Douala.}
\end{enumerate}
\textbf{Barème :} 1 pt par réponse ; retirer 0,5 pt si la réponse n'est pas une phrase complète ou si le sens est déformé.

\textbf{B. Language (6 points).}
\begin{enumerate}
\item Six partitifs du texte (3 pts — 0,5 pt par partitif correctement recopié) : \eng{a bag of maize}, \eng{a bunch of plantains}, \eng{a basket of tomatoes}, \eng{a cup of rice}, \eng{a bottle of groundnut oil}, \eng{a piece of advice}, \eng{a pair of broken chairs}, \eng{a loaf of bread}, \eng{a tin of sardines}, \eng{a sum of money}, \eng{a bar of soap}, \eng{a packet of salt}, \eng{a piece of good news}. (Six suffisent.)
\item Pluriel (2 pts — 1 pt par phrase) : \\
\eng{She sells a cup of rice.} $\rightarrow$ \eng{She sells \textbf{two cups} of rice.} \\
\eng{He buys a loaf of bread.} $\rightarrow$ \eng{He buys \textbf{two loaves} of bread.} (pluriel irrégulier \emph{loaves} exigé)
\item Correction (1 pt) : \eng{Mama Béa gave me \textbf{a piece of advice} about cooking.} — \emph{advice} est indénombrable : \eng{an advice} est impossible.
\end{enumerate}
\textbf{Points de vigilance :} ne pas recopier un partitif incomplet (\eng{a bag} sans \eng{of maize} ne vaut pas 0,5 pt entier) ; vérifier l'orthographe de \emph{loaves} ; ne pas oublier \eng{of} dans la correction.
\end{corrige}

\begin{corrige}[Exercice 10 — Type Bac : Traduction et correction d'erreurs (8 points)]
\textbf{A. Traduction (4 points — 2 pts par phrase).}
\begin{enumerate}
\item « Mon père a acheté un sac de ciment et une feuille de tôle pour réparer la maison. » \\
\eng{My father bought a bag of cement and a sheet of corrugated iron to repair the house.} \\
(\emph{cement} et \emph{iron} sont indénombrables ; \cle{a sheet of} convient aux surfaces plates : papier, tôle, métal.)
\item « La vendeuse m'a donné un morceau de savon et une bouteille d'huile. » \\
\eng{The saleswoman gave me a piece of soap and a bottle of oil.} \\
(On accepte \eng{a bar of soap} pour « une savonnette » ; \eng{a piece of soap} traduit bien « un morceau de savon ».)
\end{enumerate}
\textbf{Barème :} 1 pt pour le partitif correct, 0,5 pt pour la grammaire générale, 0,5 pt pour le vocabulaire.

\textbf{B. Correction d'erreurs (4 points — 1 pt par phrase).}
\begin{enumerate}
\item \eng{The teacher gave us an homework for the holidays.} $\rightarrow$ \eng{The teacher gave us \textbf{some homework} for the holidays.} (ou \eng{a piece of homework}). \emph{Homework} est indénombrable : jamais \eng{an homework}.
\item \eng{I bought three breads at the bakery this morning.} $\rightarrow$ \eng{I bought \textbf{three loaves of bread} at the bakery this morning.} \emph{Bread} ne prend pas de pluriel.
\item \eng{She told me a news about the new market in Douala.} $\rightarrow$ \eng{She told me \textbf{a piece of news} about the new market in Douala.} \emph{News} est indénombrable malgré le \emph{-s}.
\item \eng{There are many furnitures in the headmaster's office.} $\rightarrow$ \eng{There \textbf{is} a lot of \textbf{furniture} in the headmaster's office.} (ou \eng{many pieces of furniture}). \emph{Furniture} est indénombrable : verbe au singulier avec \eng{a lot of furniture}.
\end{enumerate}
\textbf{Points de vigilance :} ces quatre erreurs sont les plus fréquentes des candidats francophones ; les mémoriser comme des blocs : \emph{homework, bread, news, furniture} = toujours indénombrables.
\end{corrige}

\begin{corrige}[Exercice 11 — Type Bac : Rédaction guidée (8 points)]
\textbf{Proposition de rédaction modèle (environ 75 mots) :}

\begin{textebox}[My mother's shopping day]
Last Saturday, my mother went to the market early in the morning. She bought a bag of rice and a bunch of plantains for lunch. Then she chose a bottle of palm oil and a tin of tomatoes for the sauce. Before coming home, she also bought a loaf of bread, a bar of soap and a pair of shoes for my little brother. When she arrived, she gave me a piece of advice: “Never buy the first thing you see!”
\end{textebox}

Traduction : « Samedi dernier, ma mère est allée au marché tôt le matin. Elle a acheté un sac de riz et une grappe de plantains pour le déjeuner. Puis elle a choisi une bouteille d'huile de palme et une boîte de tomates pour la sauce. Avant de rentrer, elle a aussi acheté une miche de pain, une savonnette et une paire de chaussures pour mon petit frère. En arrivant, elle m'a donné un conseil : “N'achète jamais la première chose que tu vois !” »

\textbf{Barème indicatif (8 points) :}
\begin{itemize}
\item Respect de la consigne (60--80 mots, sujet traité, titre) : 2 pts
\item Au moins six partitifs différents et corrects : 3 pts (0,5 pt chacun)
\item Correction grammaticale (accord du verbe, indénombrables sans \emph{-s}, passé) : 2 pts
\item Vocabulaire et cohérence du récit : 1 pt
\end{itemize}

\textbf{Points de vigilance :}
\begin{itemize}
\item Ne pas écrire \eng{a bread}, \eng{a soap}, \eng{an advice} : toujours passer par le partitif.
\item Le verbe s'accorde avec le partitif : \eng{a bottle of oil \textbf{is}}, mais \eng{two bottles of oil \textbf{are}}.
\item Varier les partitifs : répéter \eng{a piece of} six fois ne rapporte pas tous les points.
\item Employer le prétérit pour raconter la journée : \eng{went, bought, chose, gave}.
\end{itemize}
\end{corrige}

\newpage

\coursec{Fiche bilan}

\begin{popfigure}
\begin{tikzpicture}[mindmap, grow cyclic, every node/.style={concept, text=white, font=\small\bfseries},
root concept/.append style={concept color=popblue, font=\large\bfseries},
level 1/.append style={level distance=3.4cm, sibling angle=51.5},
level 2/.append style={level distance=2.1cm, sibling angle=40, font=\scriptsize}]
\node[root concept]{Partitive expressions}
  child[concept color=poporange]{ node {Structure} 
    child { node {a / an + container + of} }
    child { node {+ noun (uncountable)} }
  }
  child[concept color=popgreen]{ node {Liquids}
    child { node {a cup of tea} }
    child { node {a bottle of water} }
    child { node {a glass of milk} }
  }
  child[concept color=poppurple]{ node {Food}
    child { node {a piece of cake} }
    child { node {a slice of bread} }
    child { node {a bar of soap} }
  }
  child[concept color=poppink]{ node {Groups}
    child { node {a bunch of keys} }
    child { node {a bunch of flowers} }
    child { node {a crowd of people} }
  }
  child[concept color=popgold]{ node {Abstract}
    child { node {a piece of advice} }
    child { node {a piece of news} }
    child { node {an item of information} }
  }
  child[concept color=popblue!70]{ node {Agreement}
    child { node {verb agrees with container} }
    child { node {plural: two cups of} }
  }
  child[concept color=popmuted]{ node {Pitfalls}
    child { node {no plural on uncountables} }
    child { node {never ``an advice''} }
  };
\end{tikzpicture}
\legende{Carte mentale : les expressions partitives (a cup of, a bunch of, a piece of...)}
\end{popfigure}

\begin{bilan}
\textbf{1. Lexique clé (English -- Français)}
\begin{itemize}
\item \eng{a cup of tea / coffee} : une tasse de thé / de café
\item \eng{a glass of water / wine} : un verre d'eau / de vin
\item \eng{a bottle of oil / beer} : une bouteille d'huile / de bière
\item \eng{a piece of cake / advice / news / furniture} : un morceau de gâteau / un conseil / une nouvelle / un meuble
\item \eng{a slice of bread / meat} : une tranche de pain / de viande
\item \eng{a bunch of keys / flowers / bananas} : un trousseau de clés / un bouquet de fleurs / un régime de bananes
\item \eng{a bar of soap / chocolate} : une barre (un pain) de savon / une tablette de chocolat
\item \eng{a loaf of bread} : un pain ; \eng{a packet of rice} : un paquet de riz ; \eng{a tin of tomatoes} : une boîte de tomates
\item \eng{a sheet of paper} : une feuille de papier ; \eng{a lump of sugar} : un morceau de sucre
\end{itemize}

\textbf{2. Règles de grammaire à connaître par cœur}
\begin{center}
\begin{tabularx}{\linewidth}{|X|X|X|}
\hline
\rowcolor{popblueL} Règle & Exemple & Remarque \\
\hline
Structure : \textbf{a/an + contenant + of + nom indénombrable} & \eng{a bottle of palm wine} & le contenant rend le nom comptable \\
\hline
Pluriel : on pluralise le \textbf{contenant}, jamais l'indénombrable & \eng{two cups of tea} (jamais \emph{two teas} au sens de « deux tasses de thé ») & \eng{three pieces of advice} \\
\hline
Accord du verbe : avec le \textbf{contenant} & \eng{This bunch of keys \textbf{is} mine.} / \eng{Two bottles of oil \textbf{are} missing.} & singulier / pluriel selon le contenant \\
\hline
Noms indénombrables pièges : advice, news, information, furniture, bread, money, luggage & \eng{a piece of information} & jamais de \textbf{-s}, jamais de \textbf{a/an} seul \\
\hline
\end{tabularx}
\end{center}

\textbf{3. Méthode express}
\begin{enumerate}
\item Repérer si le nom est \cle{indénombrable} (advice, bread, water, news...).
\item Choisir le \cle{contenant} adapté (cup, bottle, piece, bunch...).
\item Construire : \textbf{a/an + contenant + of + nom}.
\item Au pluriel : mettre le \textbf{-s} sur le contenant uniquement, et accorder le verbe avec lui.
\end{enumerate}
\end{bilan}

\begin{aretenir}[Checklist avant le Bac]
\begin{itemize}
\item[$\square$] Je connais la structure : \eng{a/an + container + of + uncountable noun}.
\item[$\square$] Je sais que \eng{advice}, \eng{news}, \eng{information}, \eng{furniture}, \eng{bread}, \eng{luggage} sont indénombrables en anglais.
\item[$\square$] Je dis \eng{a piece of advice} et jamais \emph{an advice}.
\item[$\square$] Je pluralise le contenant : \eng{two bottles of water}, pas le nom indénombrable.
\item[$\square$] J'accorde le verbe avec le contenant : \eng{The bunch of keys is on the table.}
\item[$\square$] Je connais les contenant courants : cup, glass, bottle, piece, slice, bar, loaf, sheet, lump, packet, tin, bunch.
\item[$\square$] Je sais que \eng{a bunch of} s'emploie pour keys, flowers, bananas, grapes.
\item[$\square$] Je ne traduis pas mot à mot « un conseil » : en anglais c'est \eng{a piece of advice}.
\item[$\square$] Je peux transformer : \eng{some bread} $\rightarrow$ \eng{two slices of bread}.
\item[$\square$] Je peux utiliser au moins cinq partitifs dans un court texte sur le marché ou la vie économique.
\end{itemize}
\end{aretenir}

\findecours

\end{document}
